\documentclass[12pt,oneside]{book}

%============================================================
% PACKAGES
%============================================================

%--------------------
% Page + font
%--------------------
\usepackage[margin=1in]{geometry}
\usepackage[T1]{fontenc}
\usepackage[utf8]{inputenc}
\usepackage{lmodern}
\usepackage{microtype}
\usepackage{parskip}


%--------------------
% Math
%--------------------
\usepackage{amsmath, amssymb, amsfonts, amsthm, mathtools}
\usepackage{bm}
\usepackage{bbm}
\usepackage{mathrsfs}
\usepackage{dsfont}

%--------------------
% Tables + figures
%--------------------
\usepackage{graphicx}
\usepackage{float}
\usepackage{booktabs}
\usepackage{array}
\usepackage{xltabular}
\usepackage{multirow}
\usepackage{caption}
\usepackage{subcaption}

%--------------------
% Lists
%--------------------
\usepackage{enumitem}
\setlist[itemize]{topsep=2pt,itemsep=2pt}
\setlist[enumerate]{topsep=2pt,itemsep=2pt}

%--------------------
% Colors + hyperlinks
%--------------------
\usepackage[dvipsnames]{xcolor}
\usepackage[
    colorlinks=true,
    linkcolor=MidnightBlue,
    citecolor=MidnightBlue,
    urlcolor=MidnightBlue
]{hyperref}
\usepackage[nameinlink,capitalise]{cleveref}

%--------------------
% Algorithms
%--------------------
\usepackage{algorithm}
\usepackage{algpseudocode}

%--------------------
% Code blocks
%--------------------
\usepackage{listings}

\lstdefinestyle{codestyle}{
    basicstyle=\ttfamily\small,
    breaklines=true,
    frame=single,
    numbers=left,
    numberstyle=\tiny,
    keywordstyle=\color{MidnightBlue},
    commentstyle=\color{ForestGreen},
    stringstyle=\color{BrickRed},
    showstringspaces=false,
    tabsize=4
}

\lstset{style=codestyle}

%--------------------
% Section styling
%--------------------
\usepackage{titlesec}
\usepackage{tocloft}

\setcounter{tocdepth}{2}
\setcounter{secnumdepth}{3}

%============================================================
% THEOREMS
%============================================================

\numberwithin{equation}{section}

\theoremstyle{definition}
\newtheorem{definition}{Definition}[section]
\newtheorem{example}{Example}[section]
\newtheorem{exercise}{Exercise}[section]
\newtheorem{remark}{Remark}[section]

\theoremstyle{plain}
\newtheorem{theorem}{Theorem}[section]
\newtheorem{lemma}{Lemma}[section]
\newtheorem{proposition}{Proposition}[section]
\newtheorem{corollary}{Corollary}[section]

\theoremstyle{remark}
\newtheorem*{note}{Note}

%============================================================
% CUSTOM COMMANDS
%============================================================

% Sets
\newcommand{\R}{\mathbb{R}}
\newcommand{\N}{\mathbb{N}}
\newcommand{\Z}{\mathbb{Z}}
\newcommand{\Q}{\mathbb{Q}}
\newcommand{\C}{\mathbb{C}}

% Probability
\newcommand{\E}{\mathbb{E}}
\newcommand{\Var}{\mathrm{Var}}
\newcommand{\Cov}{\mathrm{Cov}}
\newcommand{\Corr}{\mathrm{Corr}}
\newcommand{\Prob}{\mathbb{P}}

\newcommand{\ind}{\mathbbm{1}}
\newcommand{\iid}{\overset{\text{iid}}{\sim}}
\newcommand{\convd}{\overset{d}{\to}}
\newcommand{\convp}{\overset{p}{\to}}
\newcommand{\as}{\overset{a.s.}{\to}}

% Operators
\newcommand{\given}{\,\middle|\,}
\newcommand{\argmin}{\operatorname*{arg\,min}}
\newcommand{\argmax}{\operatorname*{arg\,max}}

% Linear algebra / analysis
\newcommand{\norm}[1]{\left\lVert #1 \right\rVert}
\newcommand{\abs}[1]{\left\lvert #1 \right\rvert}
\newcommand{\inner}[2]{\left\langle #1, #2 \right\rangle}

\newcommand{\dd}{\,\mathrm{d}}
\newcommand{\grad}{\nabla}
\newcommand{\Hess}{\nabla^2}

\newcommand{\tr}{\operatorname{tr}}
\newcommand{\rank}{\operatorname{rank}}
\newcommand{\diag}{\operatorname{diag}}

% Distributions
\newcommand{\Normal}{\mathcal{N}}
\newcommand{\Bern}{\mathrm{Bernoulli}}
\newcommand{\Bin}{\mathrm{Binomial}}
\newcommand{\Pois}{\mathrm{Poisson}}
\newcommand{\Gam}{\mathrm{Gamma}}
\newcommand{\BetaDist}{\mathrm{Beta}}
\newcommand{\Exp}{\mathrm{Exponential}}
\newcommand{\Dir}{\mathrm{Dirichlet}}
\newcommand{\InvGam}{\mathrm{Inv\text{-}Gamma}}

% Title block
\title{EN.553.639 Time Series Analysis}
\author{Jonathan Y. Ma}
\date{June 2026}

\begin{document}

\maketitle

\noindent \emph{This course was taught by Dr. Sergey Kushnarev in the Fall 2025 Semester. The notes are transcribed by Jonathan Ma and may contain errors and omitted content, so any issues are to be directed towards me.}

\tableofcontents

\chapter{Time Series Introduction}

\section{Time Series Notation}

We use lower case letters to denote an observed time series, such as
$x_0,x_1,x_2,\ldots$ or $y_0,y_1,y_2,\ldots$. An observed time series is a realized sequence of numbers indexed by time.

We use capital letters to denote the probabilistic model for a time series. For example,
\[
Y_0,Y_1,Y_2,\ldots
\]
denotes a sequence of random variables. We may also write the full sequence as
\[
\{Y_t\}_{t \in T},
\]
where $T$ is the time index set. Common examples include $T=\Z$, $T=\Z_{\geq 0}$, or $T=\{0,1,\ldots,n\}$.

\begin{definition}
A time series model for observed data $\{y_t\}$ is a specification of the joint distribution, or possibly only the means, variances, and covariances, of a sequence of random variables $\{Y_t:t\in T\}$ for which $\{y_t\}$ is a realization.
\end{definition}

The distinction between observed data and a probabilistic model is important. The observed time series is fixed once it has been collected. The probabilistic time series model describes a random mechanism that could have generated the observed path.

In this chapter, the stochastic process is written as a sequence of random variables $\{Y_t(\omega)\}$ indexed by time $t$. For a fixed outcome $\omega$, the function $t \mapsto Y_t(\omega)$ is one possible sample path of the process. The randomness comes from $\omega$, while the time index $t$ organizes the dependence structure across observations.

\section{Functions Used to Characterize a Time Series}

To characterize a time series, we often summarize its first and second order behavior using the mean function, variance function, autocovariance function, and autocorrelation function.

\subsection{Mean, Variance, Autocovariance, and Autocorrelation}

\begin{definition}
For a time series $\{Y_t\}$, the mean function is
\[
\mu_t := \E[Y_t].
\]
\end{definition}

\begin{definition}
For any two time points $s,t \in T$, the autocovariance function is
\[
\gamma_{t,s}
:=
\Cov(Y_t,Y_s)
=
\E[(Y_t-\mu_t)(Y_s-\mu_s)].
\]
Equivalently,
\[
\gamma_{t,s}
=
\E[Y_tY_s]-\mu_t\mu_s.
\]
\end{definition}

\begin{proof}
Starting from the definition of covariance,
\begin{align*}
\Cov(Y_t,Y_s)
&=
\E[(Y_t-\mu_t)(Y_s-\mu_s)] \\
&=
\E[Y_tY_s-\mu_sY_t-\mu_tY_s+\mu_t\mu_s] \\
&=
\E[Y_tY_s]-\mu_s\E[Y_t]-\mu_t\E[Y_s]+\mu_t\mu_s \\
&=
\E[Y_tY_s]-\mu_s\mu_t-\mu_t\mu_s+\mu_t\mu_s \\
&=
\E[Y_tY_s]-\mu_t\mu_s.
\end{align*}
\end{proof}

\begin{definition}
The variance function is the special case of the autocovariance function obtained when $s=t$:
\[
\gamma_{t,t}
=
\Cov(Y_t,Y_t)
=
\Var(Y_t).
\]
\end{definition}

\begin{definition}
For any two time points $s,t\in T$, the autocorrelation function is
\[
\rho_{t,s}
:=
\Corr(Y_t,Y_s)
=
\frac{\Cov(Y_t,Y_s)}
{\sqrt{\Var(Y_t)}\sqrt{\Var(Y_s)}}
=
\frac{\gamma_{t,s}}{\sqrt{\gamma_{t,t}\gamma_{s,s}}},
\]
provided $\gamma_{t,t}>0$ and $\gamma_{s,s}>0$.
\end{definition}

\subsection*{Basic Properties}

When $t=s$, we have
\[
\gamma_{t,t}=\Var(Y_t),
\qquad
\rho_{t,t}=1.
\]

The autocovariance and autocorrelation functions are symmetric:
\[
\gamma_{t,s}=\gamma_{s,t},
\qquad
\rho_{t,s}=\rho_{s,t}.
\]

\begin{proof}
Using the symmetry of multiplication,
\[
\gamma_{t,s}
=
\E[(Y_t-\mu_t)(Y_s-\mu_s)]
=
\E[(Y_s-\mu_s)(Y_t-\mu_t)]
=
\gamma_{s,t}.
\]
The symmetry of $\rho_{t,s}$ follows because the denominator $\sqrt{\gamma_{t,t}\gamma_{s,s}}$ is also symmetric in $t$ and $s$.
\end{proof}

By the Cauchy--Schwarz inequality,
\[
|\Cov(X,Y)| \leq \sqrt{\Var(X)\Var(Y)}.
\]
Therefore,
\[
|\gamma_{t,s}|
\leq
\sqrt{\gamma_{t,t}\gamma_{s,s}},
\qquad
|\rho_{t,s}|\leq 1.
\]

\begin{proof}
Let $X_c=X-\E[X]$ and $Y_c=Y-\E[Y]$. Then
\[
\Cov(X,Y)=\E[X_cY_c].
\]
By Cauchy--Schwarz,
\[
|\E[X_cY_c]|
\leq
\sqrt{\E[X_c^2]}\sqrt{\E[Y_c^2]}
=
\sqrt{\Var(X)\Var(Y)}.
\]
Applying this to $X=Y_t$ and $Y=Y_s$ gives
\[
|\gamma_{t,s}|
\leq
\sqrt{\gamma_{t,t}\gamma_{s,s}}.
\]
Dividing by $\sqrt{\gamma_{t,t}\gamma_{s,s}}$ gives $|\rho_{t,s}|\leq 1$.
\end{proof}

Some useful interpretations are:

\begin{itemize}
\item If $\rho_{t,s}\approx 1$, then $Y_t$ and $Y_s$ are strongly positively linearly related.
\item If $\rho_{t,s}\approx -1$, then $Y_t$ and $Y_s$ are strongly negatively linearly related.
\item If $\rho_{t,s}=0$, then $Y_t$ and $Y_s$ are uncorrelated.
\item If $Y_t$ and $Y_s$ are independent and both have finite second moments, then they are uncorrelated.
\end{itemize}

\begin{proof}
If $Y_t$ and $Y_s$ are independent, then $\E[Y_tY_s]=\E[Y_t]\E[Y_s]$. Hence
\[
\Cov(Y_t,Y_s)
=
\E[Y_tY_s]-\E[Y_t]\E[Y_s]
=
0.
\]
Thus $\gamma_{t,s}=0$ and, if the variances are nonzero, $\rho_{t,s}=0$.
\end{proof}

However, the converse is not true in general. Uncorrelated random variables need not be independent unless additional assumptions, such as joint normality, are imposed.

\subsection*{Bilinearity of Covariance}

For constants $a_1,\ldots,a_n,b_1,\ldots,b_m \in \R$ and time points $t_1,\ldots,t_n,s_1,\ldots,s_m$, covariance is bilinear:
\[
\Cov\left(\sum_{i=1}^n a_iY_{t_i},\sum_{j=1}^m b_jY_{s_j}\right)
=
\sum_{i=1}^n\sum_{j=1}^m a_ib_j\Cov(Y_{t_i},Y_{s_j}).
\]

\begin{proof}
Using the centered form of covariance,
\begin{align*}
&\Cov\left(\sum_{i=1}^n a_iY_{t_i},\sum_{j=1}^m b_jY_{s_j}\right) \\
&=
\E\left[
\left(
\sum_{i=1}^n a_iY_{t_i}
-
\E\left[\sum_{i=1}^n a_iY_{t_i}\right]
\right)
\left(
\sum_{j=1}^m b_jY_{s_j}
-
\E\left[\sum_{j=1}^m b_jY_{s_j}\right]
\right)
\right] \\
&=
\E\left[
\left(
\sum_{i=1}^n a_i(Y_{t_i}-\mu_{t_i})
\right)
\left(
\sum_{j=1}^m b_j(Y_{s_j}-\mu_{s_j})
\right)
\right] \\
&=
\E\left[
\sum_{i=1}^n\sum_{j=1}^m
a_ib_j(Y_{t_i}-\mu_{t_i})(Y_{s_j}-\mu_{s_j})
\right] \\
&=
\sum_{i=1}^n\sum_{j=1}^m
a_ib_j\E[(Y_{t_i}-\mu_{t_i})(Y_{s_j}-\mu_{s_j})] \\
&=
\sum_{i=1}^n\sum_{j=1}^m
a_ib_j\Cov(Y_{t_i},Y_{s_j}).
\end{align*}
\end{proof}

For example,
\[
\Cov(a_1Y_{t_1}+a_2Y_{t_2},Y_s)
=
a_1\Cov(Y_{t_1},Y_s)+a_2\Cov(Y_{t_2},Y_s).
\]

Also,
\begin{align*}
&\Cov(a_1Y_{t_1}+a_2Y_{t_2},b_1Y_{s_1}+b_2Y_{s_2}) \\
&=
a_1b_1\Cov(Y_{t_1},Y_{s_1})
+
a_1b_2\Cov(Y_{t_1},Y_{s_2})
+
a_2b_1\Cov(Y_{t_2},Y_{s_1})
+
a_2b_2\Cov(Y_{t_2},Y_{s_2}).
\end{align*}

\section{Examples of Time Series}

\subsection*{Linear Trend Model}

Consider the time series
\[
Y_t = a+bt+\varepsilon_t,
\qquad
\varepsilon_t \iid \Normal(0,\sigma^2).
\]
The terms $\varepsilon_t$ are often called innovations, errors, or noise terms.

The mean function is
\[
\mu_t
=
\E[Y_t]
=
\E[a+bt+\varepsilon_t]
=
a+bt.
\]

Thus the mean changes with time unless $b=0$.

For the autocovariance function,
\begin{align*}
\gamma_{t,s}
&=
\Cov(Y_t,Y_s) \\
&=
\Cov(a+bt+\varepsilon_t,a+bs+\varepsilon_s) \\
&=
\Cov(\varepsilon_t,\varepsilon_s).
\end{align*}
Since the innovations are independent with variance $\sigma^2$,
\[
\gamma_{t,s}
=
\begin{cases}
\sigma^2, & t=s,\\
0, & t\neq s.
\end{cases}
\]

The autocorrelation function is therefore
\[
\rho_{t,s}
=
\begin{cases}
1, & t=s,\\
0, & t\neq s.
\end{cases}
\]

Although the errors are uncorrelated over time, the process itself has a time-varying mean. Therefore, when $b\neq 0$, this process is not stationary in the usual weak sense.

\subsection*{Random Walk}

Let $\varepsilon_1,\varepsilon_2,\ldots \iid \Normal(0,\sigma^2)$. Define a random walk by
\[
Y_0=0,
\qquad
Y_t=\sum_{i=1}^t \varepsilon_i,
\qquad t\geq 1.
\]
Equivalently,
\[
Y_t=Y_{t-1}+\varepsilon_t.
\]

The mean function is
\[
\mu_t
=
\E[Y_t]
=
\E\left[\sum_{i=1}^t \varepsilon_i\right]
=
\sum_{i=1}^t \E[\varepsilon_i]
=
0.
\]

The variance function is
\[
\gamma_{t,t}
=
\Var(Y_t)
=
\Var\left(\sum_{i=1}^t \varepsilon_i\right).
\]
Using independence,
\[
\Var\left(\sum_{i=1}^t \varepsilon_i\right)
=
\sum_{i=1}^t \Var(\varepsilon_i)
+
2\sum_{1\leq i<j\leq t}\Cov(\varepsilon_i,\varepsilon_j).
\]
Since $\Cov(\varepsilon_i,\varepsilon_j)=0$ for $i\neq j$,
\[
\gamma_{t,t}
=
\sum_{i=1}^t \sigma^2
=
t\sigma^2.
\]
Thus the variance grows linearly with time.

Now suppose $1\leq s\leq t$. Since
\[
Y_t
=
Y_s+\sum_{i=s+1}^t \varepsilon_i,
\]
we have
\begin{align*}
\gamma_{t,s}
&=
\Cov(Y_t,Y_s) \\
&=
\Cov\left(Y_s+\sum_{i=s+1}^t\varepsilon_i,Y_s\right) \\
&=
\Cov(Y_s,Y_s)
+
\Cov\left(\sum_{i=s+1}^t\varepsilon_i,Y_s\right).
\end{align*}
The second covariance is zero because the future innovations $\varepsilon_{s+1},\ldots,\varepsilon_t$ are independent of $Y_s$, which depends only on $\varepsilon_1,\ldots,\varepsilon_s$. Therefore,
\[
\gamma_{t,s}
=
\Var(Y_s)
=
s\sigma^2.
\]
By symmetry, for arbitrary $s,t\geq 1$,
\[
\gamma_{t,s}
=
\min\{t,s\}\sigma^2.
\]

The autocorrelation function is
\[
\rho_{t,s}
=
\frac{\gamma_{t,s}}{\sqrt{\gamma_{t,t}\gamma_{s,s}}}
=
\frac{\min\{t,s\}\sigma^2}{\sqrt{t\sigma^2}\sqrt{s\sigma^2}}
=
\frac{\min\{t,s\}}{\sqrt{ts}}.
\]

For example,
\[
\rho_{1,2}
=
\frac{1}{\sqrt{2}},
\qquad
\rho_{2,3}
=
\frac{2}{\sqrt{6}},
\qquad
\rho_{3,4}
=
\frac{3}{\sqrt{12}}.
\]
Also,
\[
\rho_{t,t+1}
=
\frac{t}{\sqrt{t(t+1)}}
=
\sqrt{\frac{t}{t+1}}
\to 1
\qquad
\text{as } t\to\infty.
\]
This illustrates that adjacent values of a random walk become highly correlated at large time indices.

\begin{remark}
A random walk has constant mean equal to zero, but its variance depends on $t$. Therefore, the random walk is not weakly stationary.
\end{remark}

\begin{lstlisting}[language=R, caption={Simulating sample paths from a random walk in R}]
set.seed(123)

n <- 100
sigma <- 1

eps1 <- rnorm(n, mean = 0, sd = sigma)
eps2 <- rnorm(n, mean = 0, sd = sigma)

y1 <- cumsum(eps1)
y2 <- cumsum(eps2)

plot(
  1:n, y1, type = "l",
  xlab = "t", ylab = "Y_t",
  main = "Two Sample Paths of a Random Walk"
)
lines(1:n, y2, lty = 2)
legend(
  "topleft",
  legend = c("Path 1", "Path 2"),
  lty = c(1, 2),
  bty = "n"
)
\end{lstlisting}

\subsection*{Moving Average Process}

Let $\varepsilon_t \iid \Normal(0,\sigma^2)$ and define
\[
Y_t
=
\frac{1}{2}(\varepsilon_t+\varepsilon_{t-1}).
\]
This is a moving average process of order $1$.

The mean function is
\[
\mu_t
=
\E[Y_t]
=
\frac{1}{2}\E[\varepsilon_t+\varepsilon_{t-1}]
=
0.
\]
The mean does not depend on $t$.

The variance function is
\begin{align*}
\gamma_{t,t}
&=
\Var(Y_t) \\
&=
\Var\left(\frac{1}{2}(\varepsilon_t+\varepsilon_{t-1})\right) \\
&=
\frac{1}{4}\Var(\varepsilon_t+\varepsilon_{t-1}) \\
&=
\frac{1}{4}\left[
\Var(\varepsilon_t)+\Var(\varepsilon_{t-1})
+
2\Cov(\varepsilon_t,\varepsilon_{t-1})
\right].
\end{align*}
Since $\varepsilon_t$ and $\varepsilon_{t-1}$ are independent,
\[
\gamma_{t,t}
=
\frac{1}{4}(\sigma^2+\sigma^2)
=
\frac{1}{2}\sigma^2.
\]

Next, compute $\gamma_{t,t-1}$. We have
\[
Y_t=\frac{1}{2}(\varepsilon_t+\varepsilon_{t-1}),
\qquad
Y_{t-1}=\frac{1}{2}(\varepsilon_{t-1}+\varepsilon_{t-2}).
\]
Therefore,
\begin{align*}
\gamma_{t,t-1}
&=
\Cov(Y_t,Y_{t-1}) \\
&=
\Cov\left(
\frac{1}{2}(\varepsilon_t+\varepsilon_{t-1}),
\frac{1}{2}(\varepsilon_{t-1}+\varepsilon_{t-2})
\right) \\
&=
\frac{1}{4}\Cov(\varepsilon_t+\varepsilon_{t-1},
\varepsilon_{t-1}+\varepsilon_{t-2}) \\
&=
\frac{1}{4}\{
\Cov(\varepsilon_t,\varepsilon_{t-1})
+
\Cov(\varepsilon_t,\varepsilon_{t-2})
+
\Cov(\varepsilon_{t-1},\varepsilon_{t-1})
+
\Cov(\varepsilon_{t-1},\varepsilon_{t-2})
\}.
\end{align*}
All terms are zero except $\Cov(\varepsilon_{t-1},\varepsilon_{t-1})=\sigma^2$. Hence
\[
\gamma_{t,t-1}
=
\frac{1}{4}\sigma^2.
\]

\begin{exercise}
Show that $\gamma_{t,t-k}=0$ for $k\geq 2$.
\end{exercise}

\begin{proof}
For $k\geq 2$,
\[
Y_{t-k}
=
\frac{1}{2}(\varepsilon_{t-k}+\varepsilon_{t-k-1}).
\]
Thus
\begin{align*}
\gamma_{t,t-k}
&=
\Cov(Y_t,Y_{t-k}) \\
&=
\Cov\left(
\frac{1}{2}(\varepsilon_t+\varepsilon_{t-1}),
\frac{1}{2}(\varepsilon_{t-k}+\varepsilon_{t-k-1})
\right) \\
&=
\frac{1}{4}
\Cov(\varepsilon_t+\varepsilon_{t-1},
\varepsilon_{t-k}+\varepsilon_{t-k-1}).
\end{align*}
Expanding by bilinearity gives
\begin{align*}
\gamma_{t,t-k}
=
\frac{1}{4}\{
&\Cov(\varepsilon_t,\varepsilon_{t-k})
+
\Cov(\varepsilon_t,\varepsilon_{t-k-1}) \\
&+
\Cov(\varepsilon_{t-1},\varepsilon_{t-k})
+
\Cov(\varepsilon_{t-1},\varepsilon_{t-k-1})
\}.
\end{align*}
When $k\geq 2$, the indices $t,t-1,t-k,t-k-1$ are all distinct in the pairings above. Since the innovations are independent, all four covariance terms are zero. Therefore,
\[
\gamma_{t,t-k}=0.
\]
\end{proof}

Combining the cases, the autocovariance function is
\[
\gamma_{t,s}
=
\begin{cases}
\frac{1}{2}\sigma^2, & t=s,\\
\frac{1}{4}\sigma^2, & |t-s|=1,\\
0, & |t-s|\geq 2.
\end{cases}
\]

The autocorrelation function is
\[
\rho_{t,s}
=
\frac{\gamma_{t,s}}{\sqrt{\gamma_{t,t}\gamma_{s,s}}}.
\]
Since $\gamma_{t,t}=\gamma_{s,s}=\frac{1}{2}\sigma^2$, we get
\[
\rho_{t,s}
=
\begin{cases}
1, & t=s,\\
\frac{1}{2}, & |t-s|=1,\\
0, & |t-s|\geq 2.
\end{cases}
\]

In this example, both the autocovariance and autocorrelation functions depend only on the lag $|t-s|$, not on the particular values of $t$ and $s$.

\section*{Stationarity}

Stationarity formalizes the idea that the probabilistic behavior of a time series does not change over time. There are two common forms of stationarity:

\begin{enumerate}
\item Strict stationarity.
\item Weak stationarity, also called second order stationarity.
\end{enumerate}

Strict stationarity requires invariance of the full joint distribution under time shifts. Weak stationarity only requires invariance of the mean and autocovariance structure.

\subsection*{Strict Stationarity}

\begin{definition}
A stochastic process is strictly stationary if all finite-dimensional joint distributions are invariant under shifts in time. That is, for any positive integer $n$, any time indices $t_1,\ldots,t_n$, and any integer shift $h$,
\[
(Y_{t_1},\ldots,Y_{t_n})
\overset{d}{=}
(Y_{t_1+h},\ldots,Y_{t_n+h}),
\]
whenever the shifted indices are in the time index set.
\end{definition}

For a time series $\{Y_t\}_{t\in\Z}$, strict stationarity means
\[
F_{Y_{t_1},\ldots,Y_{t_n}}(y_1,\ldots,y_n)
=
F_{Y_{t_1+h},\ldots,Y_{t_n+h}}(y_1,\ldots,y_n)
\]
for every $n$, every $t_1,\ldots,t_n\in\Z$, every $h\in\Z$, and every $(y_1,\ldots,y_n)\in\R^n$.

Several necessary consequences follow immediately from strict stationarity. For example,
\[
Y_t \overset{d}{=} Y_s
\qquad
\text{for all } t,s.
\]
Also,
\[
(Y_1,Y_2)\overset{d}{=}(Y_{10},Y_{11}),
\]
and
\[
(Y_1,Y_3,Y_7)\overset{d}{=}(Y_{10},Y_{12},Y_{16}).
\]
These are only necessary implications; by themselves, they are not sufficient to verify strict stationarity.

\begin{exercise}
Show that for a strictly stationary time series with finite second moments, the autocovariance $\gamma_{t,s}$ depends only on the lag $t-s$.
\end{exercise}

\begin{proof}
Assume $\{Y_t\}$ is strictly stationary and has finite second moments. First, strict stationarity with $n=1$ implies
\[
Y_t \overset{d}{=} Y_0
\]
for every $t$. Hence
\[
\E[Y_t]=\E[Y_0]=:\mu
\]
for every $t$.

Now take any two time points $t$ and $s$. Let $h=-s$. By strict stationarity with $n=2$,
\[
(Y_t,Y_s)
\overset{d}{=}
(Y_{t-s},Y_0).
\]
Since equal joint distributions imply equal expectations of measurable functions whenever the expectations exist,
\[
\E[Y_tY_s]
=
\E[Y_{t-s}Y_0].
\]
Therefore,
\begin{align*}
\gamma_{t,s}
&=
\Cov(Y_t,Y_s) \\
&=
\E[Y_tY_s]-\E[Y_t]\E[Y_s] \\
&=
\E[Y_{t-s}Y_0]-\mu^2.
\end{align*}
Define
\[
\gamma(t-s)
:=
\Cov(Y_{t-s},Y_0)
=
\E[Y_{t-s}Y_0]-\mu^2.
\]
Then
\[
\gamma_{t,s}=\gamma(t-s).
\]
Thus the autocovariance depends only on the lag $t-s$.
\end{proof}

By symmetry of covariance,
\[
\gamma(t-s)=\gamma(s-t).
\]
Therefore, for a strictly stationary time series with finite second moments, we often write the autocovariance as a function of a nonnegative lag:
\[
\gamma_h := \Cov(Y_{t+h},Y_t),
\qquad h\geq 0.
\]

\subsection*{Weak Stationarity}

\begin{definition}
A time series $\{Y_t\}$ is weakly stationary if:
\begin{enumerate}
\item $\E[Y_t]=\mu$ for all $t$.
\item $\Var(Y_t)<\infty$ for all $t$.
\item $\Cov(Y_t,Y_s)$ depends only on the lag $t-s$.
\end{enumerate}
\end{definition}

For a weakly stationary time series, we write
\[
\gamma_h
:=
\Cov(Y_{t+h},Y_t),
\]
and
\[
\rho_h
:=
\frac{\gamma_h}{\gamma_0},
\]
where
\[
\gamma_0=\Var(Y_t).
\]

\begin{remark}
Strict stationarity does not automatically imply weak stationarity unless second moments exist. If second moments are finite, strict stationarity implies weak stationarity. Conversely, weak stationarity does not generally imply strict stationarity. However, for Gaussian time series, weak stationarity implies strict stationarity because the finite-dimensional distributions are fully determined by their means and covariances.
\end{remark}

\begin{example}
The linear trend model $Y_t=a+bt+\varepsilon_t$ is not weakly stationary when $b\neq 0$ because $\E[Y_t]=a+bt$ depends on $t$.
\end{example}

\begin{example}
The random walk $Y_t=\sum_{i=1}^t\varepsilon_i$ is not weakly stationary because
\[
\Var(Y_t)=t\sigma^2,
\]
which depends on $t$.
\end{example}

\begin{example}
The moving average process
\[
Y_t=\frac{1}{2}(\varepsilon_t+\varepsilon_{t-1})
\]
is weakly stationary because
\[
\E[Y_t]=0,
\]
and
\[
\gamma_{t,s}
=
\begin{cases}
\frac{1}{2}\sigma^2, & t=s,\\
\frac{1}{4}\sigma^2, & |t-s|=1,\\
0, & |t-s|\geq 2,
\end{cases}
\]
which depends only on $|t-s|$.
\end{example}

\section{Stationarity}

\subsection{Strict Stationarity}

\begin{definition}
A time series $\{Y_t\}_{t\in\Z}$ is strictly stationary if all finite-dimensional joint distributions are invariant under time shifts. That is, for every positive integer $n$, every collection of time points $t_1,\ldots,t_n\in\Z$, and every shift $k\in\Z$,
\[
(Y_{t_1},\ldots,Y_{t_n})
\overset{d}{=}
(Y_{t_1-k},\ldots,Y_{t_n-k}).
\]
Equivalently,
\[
F_{Y_{t_1},\ldots,Y_{t_n}}(y_1,\ldots,y_n)
=
F_{Y_{t_1-k},\ldots,Y_{t_n-k}}(y_1,\ldots,y_n)
\]
for all $(y_1,\ldots,y_n)\in\R^n$.
\end{definition}

Strict stationarity is a distributional invariance condition. It says that shifting all time indices by the same amount does not change the joint law of the process.

\subsection*{Properties of Strict Stationarity}

Suppose $\{Y_t\}$ is strictly stationary.

Taking $n=1$ in the definition gives
\[
Y_t \overset{d}{=} Y_s
\]
for all $t,s\in\Z$. Therefore, all one-dimensional marginal distributions are identical.

If the distribution of $Y_t$ has a finite second moment, then
\[
\E[Y_t^2] < \infty
\]
for one value of $t$ implies
\[
\E[Y_s^2]<\infty
\]
for every $s\in\Z$, because all $Y_s$ have the same marginal distribution.

Since $Y_t \overset{d}{=} Y_s$, we also have
\[
\mu_t
=
\E[Y_t]
=
\E[Y_s]
=
\mu,
\]
so the mean function does not depend on time. Similarly,
\[
\sigma_t^2
=
\Var(Y_t)
=
\Var(Y_s)
=
\sigma^2,
\]
so the variance function does not depend on time.

Now take $n=2$. For any $t,s\in\Z$, strict stationarity gives
\[
(Y_t,Y_s)
\overset{d}{=}
(Y_{t-s},Y_0).
\]
Therefore, if second moments exist,
\begin{align*}
\Cov(Y_t,Y_s)
&=
\E[Y_tY_s]-\E[Y_t]\E[Y_s] \\
&=
\E[Y_{t-s}Y_0]-\mu^2 \\
&=
\Cov(Y_{t-s},Y_0).
\end{align*}
Thus the autocovariance function only depends on the lag $t-s$.

\begin{definition}
For a strictly stationary time series with finite second moments, define
\[
\gamma_k
:=
\Cov(Y_t,Y_{t-k}).
\]
By strict stationarity, $\gamma_k$ does not depend on $t$.
\end{definition}

By symmetry of covariance,
\[
\gamma_k
=
\Cov(Y_t,Y_{t-k})
=
\Cov(Y_{t-k},Y_t)
=
\gamma_{-k}.
\]
Thus the autocovariance function is an even function of the lag.

In summary, for a strictly stationary time series with finite second moments,
\[
\mu_t=\mu,
\qquad
\gamma_0=\Var(Y_t)=\sigma^2,
\qquad
\gamma_k=\Cov(Y_t,Y_{t-k}),
\]
and
\[
\rho_k
=
\frac{\gamma_k}{\gamma_0}.
\]

\subsection{Weak Stationarity}

\begin{definition}
A time series $\{Y_t\}_{t\in\Z}$ is weakly stationary, also called second-order stationary or covariance stationary, if:
\begin{enumerate}
\item $\E[Y_t]=\mu$ for some finite constant $\mu$ and all $t\in\Z$.
\item $\Var(Y_t)=\sigma^2<\infty$ for all $t\in\Z$.
\item $\Cov(Y_t,Y_{t-k})=\gamma_k$ for some function $\gamma_k$ that depends only on the lag $k$ and not on the time index $t$.
\end{enumerate}
\end{definition}

Weak stationarity only controls the first two moments of the process. Strict stationarity controls the full joint distribution. Therefore, strict stationarity with finite second moments implies weak stationarity, but weak stationarity does not generally imply strict stationarity.

\subsection{Examples of Stationary and Nonstationary Time Series}

\subsubsection{Example 1: IID Noise}

Let $\varepsilon_t \iid \Normal(0,\sigma^2)$. We show that $\{\varepsilon_t\}$ is strictly stationary.

For any $t_1,\ldots,t_n\in\Z$ and any shift $k\in\Z$,
\[
(\varepsilon_{t_1},\ldots,\varepsilon_{t_n})
\]
and
\[
(\varepsilon_{t_1-k},\ldots,\varepsilon_{t_n-k})
\]
have the same joint distribution because both are vectors of independent $\Normal(0,\sigma^2)$ random variables. More explicitly,
\begin{align*}
F_{\varepsilon_{t_1},\ldots,\varepsilon_{t_n}}(a_1,\ldots,a_n)
&=
\Prob(\varepsilon_{t_1}\leq a_1,\ldots,\varepsilon_{t_n}\leq a_n) \\
&=
\prod_{j=1}^n \Prob(\varepsilon_{t_j}\leq a_j) \\
&=
\prod_{j=1}^n \Prob(\varepsilon_{t_j-k}\leq a_j) \\
&=
\Prob(\varepsilon_{t_1-k}\leq a_1,\ldots,\varepsilon_{t_n-k}\leq a_n) \\
&=
F_{\varepsilon_{t_1-k},\ldots,\varepsilon_{t_n-k}}(a_1,\ldots,a_n).
\end{align*}
Hence $\{\varepsilon_t\}$ is strictly stationary.

It is also weakly stationary. The mean is
\[
\E[\varepsilon_t]=0,
\]
the variance is
\[
\Var(\varepsilon_t)=\sigma^2,
\]
and the autocovariance function is
\[
\Cov(\varepsilon_t,\varepsilon_{t-k})
=
\begin{cases}
\sigma^2, & k=0,\\
0, & k\neq 0.
\end{cases}
\]
Thus the autocovariance function depends only on the lag $k$.

\subsubsection{Example 2: Weakly Stationary but Not Strictly Stationary}

Let $U_t \iid \Normal(0,1)$, and define
\[
X_t
=
\begin{cases}
U_t, & t \text{ is even},\\
\dfrac{1}{\sqrt{2}}(U_t^2-1), & t \text{ is odd}.
\end{cases}
\]

First, compute the mean. For even $t$,
\[
\E[X_t]=\E[U_t]=0.
\]
For odd $t$,
\[
\E[X_t]
=
\frac{1}{\sqrt{2}}\E[U_t^2-1]
=
\frac{1}{\sqrt{2}}(\E[U_t^2]-1)
=
\frac{1}{\sqrt{2}}(1-1)
=
0.
\]
Thus $\mu_t=0$ for all $t$.

Next, compute the variance. For even $t$,
\[
\Var(X_t)=\Var(U_t)=1.
\]
For odd $t$,
\[
\Var(X_t)
=
\Var\left(\frac{1}{\sqrt{2}}(U_t^2-1)\right)
=
\frac{1}{2}\Var(U_t^2-1)
=
\frac{1}{2}\Var(U_t^2).
\]
Since $U_t\sim \Normal(0,1)$,
\[
\E[U_t^2]=1,
\qquad
\E[U_t^4]=3.
\]
Therefore,
\[
\Var(U_t^2)
=
\E[U_t^4]-(\E[U_t^2])^2
=
3-1
=
2.
\]
Hence
\[
\Var(X_t)=\frac{1}{2}\cdot 2=1.
\]
So $\Var(X_t)=1$ for all $t$.

For $k\neq 0$, $X_t$ and $X_{t-k}$ are functions of independent random variables $U_t$ and $U_{t-k}$. Hence,
\[
\Cov(X_t,X_{t-k})=0.
\]
For $k=0$,
\[
\Cov(X_t,X_t)=\Var(X_t)=1.
\]
Thus
\[
\Cov(X_t,X_{t-k})
=
\begin{cases}
1, & k=0,\\
0, & k\neq 0.
\end{cases}
\]
Therefore, $\{X_t\}$ is weakly stationary.

\begin{exercise}
Show that $\{X_t\}$ is not strictly stationary.
\end{exercise}

\begin{proof}
To show that the process is not strictly stationary, it is enough to show that the one-dimensional marginal distributions are not identical for all $t$.

For even $t$, $X_t=U_t$, so $X_t\sim\Normal(0,1)$. In particular,
\[
\Prob(X_t\leq 0)=\frac{1}{2}.
\]

For odd $t$,
\[
X_t=\frac{1}{\sqrt{2}}(U_t^2-1).
\]
Then
\[
X_t\leq 0
\iff
\frac{1}{\sqrt{2}}(U_t^2-1)\leq 0
\iff
U_t^2\leq 1
\iff
-1\leq U_t\leq 1.
\]
Therefore,
\[
\Prob(X_t\leq 0)
=
\Prob(-1\leq U_t\leq 1).
\]
Since $U_t\sim\Normal(0,1)$,
\[
\Prob(-1\leq U_t\leq 1)\neq \frac{1}{2}.
\]
Thus the marginal distribution of $X_t$ depends on whether $t$ is even or odd. Hence the process is not strictly stationary.
\end{proof}

\begin{lstlisting}[language=R, caption={Checking the different marginal distributions in R}]
set.seed(123)

n <- 100000
u <- rnorm(n)

x_even <- u
x_odd <- (u^2 - 1) / sqrt(2)

mean(x_even <= 0)
mean(x_odd <= 0)

var(x_even)
var(x_odd)
\end{lstlisting}

\subsubsection*{Clarification on IID Noise and White Noise}

IID noise and white noise are related but distinct concepts.

\begin{definition}
A time series $\{\varepsilon_t\}$ is IID noise if the random variables $\varepsilon_t$ are independent and identically distributed.
\end{definition}

\begin{definition}
A time series $\{\varepsilon_t\}$ is white noise with mean $0$ and variance $\sigma^2$, written
\[
\varepsilon_t \sim WN(0,\sigma^2),
\]
if
\[
\E[\varepsilon_t]=0,
\qquad
\Var(\varepsilon_t)=\sigma^2,
\qquad
\Cov(\varepsilon_t,\varepsilon_s)=0
\quad
\text{for } t\neq s.
\]
\end{definition}

IID noise implies white noise when the variance is finite, because independence implies zero covariance. However, white noise does not necessarily imply IID noise, because uncorrelatedness is weaker than independence.

If the process is Gaussian, then uncorrelatedness implies independence. Thus, for Gaussian processes,
\[
\text{IID Gaussian noise}
\iff
\text{Gaussian white noise}.
\]

\begin{remark}
In many applied time series courses, IID noise and white noise are often used interchangeably when Gaussianity is assumed. Without Gaussianity, they are not the same concept.
\end{remark}

\subsubsection*{More Examples}

Consider
\[
Y_t=a+bt+\varepsilon_t,
\qquad
\varepsilon_t\sim WN(0,\sigma^2).
\]
Then
\[
\E[Y_t]=a+bt.
\]
Therefore, if $b\neq 0$, the process is not stationary because its mean depends on time.

A random walk is defined by
\[
Y_t=\sum_{i=1}^t \varepsilon_i,
\qquad
\varepsilon_t\sim WN(0,\sigma^2).
\]
Then
\[
\E[Y_t]=0,
\qquad
\Var(Y_t)=t\sigma^2.
\]
Thus the process is not stationary because the variance depends on time.

A moving average process of order $1$ is
\[
Y_t=\frac{1}{2}\varepsilon_t+\frac{1}{2}\varepsilon_{t-1},
\qquad
\varepsilon_t\sim WN(0,\sigma^2).
\]
As shown earlier,
\[
\E[Y_t]=0,
\]
and
\[
\gamma_k
=
\begin{cases}
\frac{1}{2}\sigma^2, & k=0,\\
\frac{1}{4}\sigma^2, & |k|=1,\\
0, & |k|\geq 2.
\end{cases}
\]
Thus the process is weakly stationary.

\subsubsection{Example 3: A Random Amplitude Sinusoid}

Let $A$ and $B$ be random variables satisfying
\[
\E[A]=\E[B]=0,
\qquad
\Var(A)=\Var(B)=\sigma^2.
\]
Let $\omega\in\R$ be fixed. For each $t\in\Z$, define
\[
Y_t=A\cos(\omega t)+B\sin(\omega t).
\]
The same random variables $A$ and $B$ are used for all time points.

The mean function is
\begin{align*}
\mu_t
&=
\E[Y_t] \\
&=
\E[A\cos(\omega t)+B\sin(\omega t)] \\
&=
\cos(\omega t)\E[A]+\sin(\omega t)\E[B] \\
&=
0.
\end{align*}

Assume additionally that $\Cov(A,B)=0$. Then
\begin{align*}
\Var(Y_t)
&=
\Var(A\cos(\omega t)+B\sin(\omega t)) \\
&=
\cos^2(\omega t)\Var(A)
+
\sin^2(\omega t)\Var(B)
+
2\cos(\omega t)\sin(\omega t)\Cov(A,B) \\
&=
\sigma^2\cos^2(\omega t)+\sigma^2\sin^2(\omega t) \\
&=
\sigma^2.
\end{align*}

Now compute the autocovariance:
\begin{align*}
\Cov(Y_t,Y_{t-k})
&=
\Cov(A\cos(\omega t)+B\sin(\omega t),
A\cos(\omega(t-k))+B\sin(\omega(t-k))) \\
&=
\cos(\omega t)\cos(\omega(t-k))\Var(A) \\
&\quad
+
\sin(\omega t)\sin(\omega(t-k))\Var(B) \\
&\quad
+
\cos(\omega t)\sin(\omega(t-k))\Cov(A,B) \\
&\quad
+
\sin(\omega t)\cos(\omega(t-k))\Cov(B,A).
\end{align*}
Using $\Var(A)=\Var(B)=\sigma^2$ and $\Cov(A,B)=0$,
\begin{align*}
\Cov(Y_t,Y_{t-k})
&=
\sigma^2\{\cos(\omega t)\cos(\omega(t-k))
+
\sin(\omega t)\sin(\omega(t-k))\}.
\end{align*}
Using the trigonometric identity
\[
\cos(x-y)=\cos x\cos y+\sin x\sin y,
\]
we obtain
\[
\Cov(Y_t,Y_{t-k})
=
\sigma^2\cos(\omega k).
\]
This depends only on the lag $k$, so $\{Y_t\}$ is weakly stationary.

\begin{exercise}
Show that $Y_0$ and $Y_1$ are not identically distributed in general.
\end{exercise}

\begin{proof}
We have
\[
Y_0=A\cos(0)+B\sin(0)=A.
\]
Also,
\[
Y_1=A\cos(\omega)+B\sin(\omega).
\]
Unless the joint distribution of $(A,B)$ is rotationally invariant, the distribution of $A$ need not equal the distribution of $A\cos(\omega)+B\sin(\omega)$.

For example, let $A$ and $B$ be independent Rademacher random variables taking values $-1$ and $1$ with probability $1/2$ each. Then $\E[A]=\E[B]=0$ and $\Var(A)=\Var(B)=1$.

Take $\omega=\pi/4$. Then
\[
Y_0=A\in\{-1,1\}.
\]
But
\[
Y_1
=
\frac{1}{\sqrt{2}}A+\frac{1}{\sqrt{2}}B
=
\frac{A+B}{\sqrt{2}}.
\]
Thus $Y_1$ takes values
\[
-\sqrt{2},\quad 0,\quad \sqrt{2}.
\]
Therefore $Y_0$ and $Y_1$ are not identically distributed. Hence the process is not strictly stationary in general.
\end{proof}

\section*{$q$-Dependent Central Limit Theorem}

Suppose $\{Y_t\}$ is a stationary time series with
\[
\mu_t=\mu,
\qquad
\Var(Y_t)=\sigma^2,
\qquad
\Cov(Y_t,Y_{t-k})=\gamma_k.
\]
For the sample mean
\[
\overline{Y}
=
\frac{1}{n}\sum_{t=1}^n Y_t,
\]
we want to study the behavior of $\overline{Y}$ as an estimator of $\mu$.

If $Y_1,\ldots,Y_n$ are IID, then the classical central limit theorem gives
\[
\sqrt{n}(\overline{Y}-\mu)
\convd
\Normal(0,\sigma^2).
\]
For time series data, this result may fail or require modification because observations can be dependent.

\subsection*{Motivating Example}

Let
\[
Z\sim\Normal(0,1),
\qquad
U_t \iid \Normal(0,1),
\]
where $Z$ is independent of all $\{U_t\}$. Define
\[
Y_t=Z+U_t.
\]
Then
\[
\E[Y_t]=0,
\]
and
\[
\Var(Y_t)=\Var(Z)+\Var(U_t)=2.
\]
For $k>0$,
\[
\gamma_k
=
\Cov(Y_t,Y_{t-k})
=
\Cov(Z+U_t,Z+U_{t-k})
=
\Var(Z)=1.
\]
Thus the process is stationary with persistent dependence at every lag.

The sample mean is
\[
\overline{Y}
=
\frac{1}{n}\sum_{t=1}^n (Z+U_t)
=
Z+\overline{U}.
\]
Therefore,
\[
\overline{Y}\convd Z
\]
because $\overline{U}\convp 0$. Also,
\[
\Var(\overline{Y})
=
\Var(Z)+\Var(\overline{U})
=
1+\frac{1}{n}.
\]
Hence $\overline{Y}$ does not converge in probability to $0$, and the usual law of large numbers fails.

If one incorrectly applied the IID CLT, one might expect $\Var(\overline{Y})$ to shrink like $2/n$. But the actual variance converges to $1$, not to $0$. The common random component $Z$ creates dependence that does not vanish as the sample size grows.

This motivates the need for dependence restrictions that are weak enough to allow time series dependence but strong enough for asymptotic averaging.

\subsection*{$q$-Dependence}

\begin{definition}
A time series $\{Y_t\}$ is $q$-dependent if $Y_t$ and $Y_{t-k}$ are independent for every $t$ and every $|k|>q$.
\end{definition}

\begin{definition}
A time series $\{Y_t\}$ is $q$-correlated if
\[
\gamma_k=0
\]
for every $|k|>q$.
\end{definition}

$q$-dependence means that observations more than $q$ time steps apart are independent. $q$-correlation is weaker: it only says that observations more than $q$ time steps apart are uncorrelated.

In the motivating example $Y_t=Z+U_t$, we have
\[
\gamma_k=1
\]
for all $k>0$. Therefore, this process is not $q$-correlated and not $q$-dependent for any finite $q$.

\subsection*{Variance of the Sample Mean Under Stationarity}

Let $\{Y_t\}$ be weakly stationary with autocovariance function $\gamma_k$. Then
\[
\Var(\overline{Y})
=
\Var\left(\frac{1}{n}\sum_{t=1}^n Y_t\right)
=
\frac{1}{n^2}\sum_{t=1}^n\sum_{s=1}^n \Cov(Y_t,Y_s).
\]
By stationarity,
\[
\Cov(Y_t,Y_s)=\gamma_{t-s}.
\]
For each lag $k$, there are $n-|k|$ pairs $(t,s)$ such that $t-s=k$. Therefore,
\[
\Var(\overline{Y})
=
\frac{1}{n^2}
\sum_{k=-(n-1)}^{n-1}
(n-|k|)\gamma_k.
\]
Using symmetry $\gamma_k=\gamma_{-k}$,
\[
\Var(\overline{Y})
=
\frac{1}{n}
\left[
\gamma_0
+
2\sum_{k=1}^{n-1}
\left(1-\frac{k}{n}\right)\gamma_k
\right].
\]

If the process is $q$-correlated, then $\gamma_k=0$ for $|k|>q$, so for $n>q$,
\[
\Var(\overline{Y})
=
\frac{1}{n}
\left[
\gamma_0
+
2\sum_{k=1}^{q}
\left(1-\frac{k}{n}\right)\gamma_k
\right].
\]
Thus,
\[
n\Var(\overline{Y})
\to
\gamma_0+2\sum_{k=1}^q \gamma_k.
\]

\begin{definition}
For a stationary $q$-correlated process, the long-run variance is
\[
\tau^2
=
\gamma_0+2\sum_{k=1}^q \gamma_k.
\]
\end{definition}

\subsection*{Central Limit Theorem for $q$-Dependent Processes}

\begin{theorem}
Let $\{Y_t\}$ be a stationary $q$-dependent time series with $\E[Y_t]=\mu$ and $\Var(Y_t)<\infty$. Under appropriate regularity conditions,
\[
\sqrt{n}(\overline{Y}-\mu)
\convd
\Normal(0,\tau^2),
\]
where
\[
\tau^2
=
\gamma_0+2\sum_{k=1}^q \gamma_k.
\]
\end{theorem}

\begin{proof}
The key idea is to group observations into blocks separated far enough apart that the blocks are independent.

For simplicity, suppose $q$ is fixed and $n$ is large. Divide the series into blocks whose lengths exceed $q$. Blocks separated by more than $q$ observations are independent by $q$-dependence. One can then apply the classical central limit theorem to sums of approximately independent blocks.

The variance of the normalized sum is determined by the covariance structure:
\[
\Var\left(\frac{1}{\sqrt{n}}\sum_{t=1}^n (Y_t-\mu)\right)
=
n\Var(\overline{Y}).
\]
From the previous calculation,
\[
n\Var(\overline{Y})
=
\gamma_0
+
2\sum_{k=1}^{q}
\left(1-\frac{k}{n}\right)\gamma_k
\to
\gamma_0+2\sum_{k=1}^q\gamma_k
=
\tau^2.
\]
The block CLT argument gives asymptotic normality, and the limiting variance is $\tau^2$.
\end{proof}

\section{$q$-Dependent Central Limit Theorem}

Last time, we defined $q$-dependent and $q$-correlated time series. The main point is that if dependence only extends for a finite number of lags, then a central limit theorem can still hold, but the limiting variance must account for autocovariance.

\subsection{Statement of the Theorem}

\begin{theorem}[$q$-Dependent Central Limit Theorem]
Let $\{Y_t\}$ be a stationary $q$-dependent time series with
\[
\mu=\E[Y_t],
\qquad
\sigma^2=\Var(Y_t),
\]
and suppose
\[
\sigma^2+2\Cov(Y_1,Y_2)+\cdots+2\Cov(Y_1,Y_{q+1})>0.
\]
Then
\[
\frac{\sum_{i=1}^n Y_i-n\mu}
{\sqrt{\Var\left(\sum_{i=1}^nY_i\right)}}
\convd
\Normal(0,1),
\qquad
n\to\infty.
\]
Equivalently,
\[
\frac{\overline{Y}-\mu}{\sqrt{\Var(\overline{Y})}}
\convd
\Normal(0,1),
\qquad
n\to\infty.
\]
\end{theorem}

\begin{remark}
The second statement can be read heuristically as
\[
\overline{Y}\approx \Normal(\mu,\Var(\overline{Y})).
\]
Unlike the IID case, $\Var(\overline{Y})$ is not generally equal to $\sigma^2/n$. It must include the effect of serial dependence.
\end{remark}

\subsection{Sketch of the Proof}

Suppose $n>q$. Since $\{Y_t\}$ is stationary, the autocovariance function satisfies
\[
\Cov(Y_i,Y_j)=\gamma_{i-j}.
\]
Therefore,
\begin{align*}
\Var\left(\sum_{i=1}^nY_i\right)
&=
\sum_{i=1}^n\sum_{j=1}^n \Cov(Y_i,Y_j) \\
&=
\sum_{i=1}^n\sum_{j=1}^n \gamma_{i-j}.
\end{align*}

For lag $0$, there are $n$ terms, each equal to $\gamma_0$. For lag $j>0$, there are $n-j$ pairs with distance $j$ in one direction and another $n-j$ pairs with distance $-j$. Since $\gamma_j=\gamma_{-j}$, we get
\[
\Var\left(\sum_{i=1}^nY_i\right)
=
n\gamma_0+2\sum_{j=1}^{n-1}(n-j)\gamma_j.
\]
Because the process is $q$-dependent, $\gamma_j=0$ for $j>q$. Hence
\[
\Var\left(\sum_{i=1}^nY_i\right)
=
n\gamma_0+2\sum_{j=1}^{q}(n-j)\gamma_j.
\]
Dividing by $n^2$ gives
\begin{align*}
\Var(\overline{Y})
&=
\frac{1}{n^2}\Var\left(\sum_{i=1}^nY_i\right) \\
&=
\frac{1}{n}\gamma_0
+
\frac{2}{n^2}\sum_{j=1}^{q}(n-j)\gamma_j \\
&=
\frac{1}{n}
\left[
\gamma_0+2\sum_{j=1}^{q}\left(1-\frac{j}{n}\right)\gamma_j
\right].
\end{align*}
Thus,
\[
n\Var(\overline{Y})
\to
\gamma_0+2\sum_{j=1}^{q}\gamma_j.
\]

Define the long-run variance
\[
\tau^2
:=
\gamma_0+2\sum_{j=1}^{q}\gamma_j.
\]
The positivity assumption in the theorem is exactly the assumption that $\tau^2>0$. Under the $q$-dependence condition, blocks of observations that are separated by more than $q$ time steps are independent. This allows one to apply a block version of the classical central limit theorem, giving
\[
\frac{\overline{Y}-\mu}{\sqrt{\Var(\overline{Y})}}
\convd
\Normal(0,1).
\]

\begin{remark}
A central limit theorem can also hold for some time series that are not $q$-dependent, as long as the dependence decays sufficiently fast. In that case, the long-run variance typically has the form
\[
\tau^2
=
\gamma_0+2\sum_{j=1}^{\infty}\gamma_j,
\]
provided the infinite sum is well-defined.
\end{remark}

\chapter{Model Specification}

\section{Moving Average (MA) Processes}

\subsection*{The MA($q$) Model}

MA($q$) stands for moving average of order $q$.

\begin{definition}
A time series $\{Y_t\}$ is an MA($q$) process if
\[
Y_t
=
\varepsilon_t-\theta_1\varepsilon_{t-1}
-\theta_2\varepsilon_{t-2}
-\cdots
-\theta_q\varepsilon_{t-q},
\]
where
\[
\varepsilon_t \iid \Normal(0,\sigma^2).
\]
\end{definition}

The coefficient on $\varepsilon_t$ is normalized to be $1$, while $\theta_1,\ldots,\theta_q$ are unrestricted constants. Some books and software use the alternative sign convention
\[
Y_t
=
\varepsilon_t+\theta_1\varepsilon_{t-1}
+\cdots
+\theta_q\varepsilon_{t-q}.
\]
The sign convention does not change the substance of the model, but it changes the signs of the reported coefficients.

For notational convenience, write
\[
\theta_0=-1.
\]
Then
\[
Y_t
=
-\sum_{j=0}^{q}\theta_j\varepsilon_{t-j}.
\]

\subsection*{Autocovariance Function of MA($q$)}

First, note that
\[
\E[Y_t]=0.
\]
For the autocovariance, consider
\[
\gamma_k
=
\Cov(Y_t,Y_{t-k}).
\]
Using the MA($q$) representation,
\[
Y_t
=
-\sum_{j=0}^{q}\theta_j\varepsilon_{t-j},
\qquad
Y_{t-k}
=
-\sum_{\ell=0}^{q}\theta_\ell\varepsilon_{t-k-\ell}.
\]
Therefore,
\begin{align*}
\gamma_k
&=
\Cov\left(
-\sum_{j=0}^{q}\theta_j\varepsilon_{t-j},
-\sum_{\ell=0}^{q}\theta_\ell\varepsilon_{t-k-\ell}
\right) \\
&=
\sum_{j=0}^{q}\sum_{\ell=0}^{q}
\theta_j\theta_\ell
\Cov(\varepsilon_{t-j},\varepsilon_{t-k-\ell}).
\end{align*}
Since the innovations are IID,
\[
\Cov(\varepsilon_{t-j},\varepsilon_{t-k-\ell})
=
\begin{cases}
\sigma^2, & t-j=t-k-\ell,\\
0, & t-j\neq t-k-\ell.
\end{cases}
\]
The equality $t-j=t-k-\ell$ is equivalent to
\[
\ell=j-k.
\]
Hence only terms with $\ell=j-k$ contribute. For $k\geq 0$, this gives
\[
\gamma_k
=
\sigma^2
\sum_{j=k}^{q}\theta_j\theta_{j-k},
\qquad
0\leq k\leq q.
\]
If $k>q$, there is no overlap between the innovations appearing in $Y_t$ and $Y_{t-k}$, so
\[
\gamma_k=0.
\]
Thus
\[
\gamma_k
=
\begin{cases}
0, & |k|>q,\\
\sigma^2\sum_{j=0}^{q-|k|}\theta_j\theta_{j+|k|}, & |k|\leq q.
\end{cases}
\]
With the convention $\theta_0=-1$, this matches the chosen sign convention above.

\begin{remark}
This calculation shows that every MA($q$) process is $q$-correlated. In fact, because $Y_t$ and $Y_{t-k}$ are functions of disjoint innovation sets whenever $|k|>q$, an MA($q$) process is also $q$-dependent.
\end{remark}

\begin{theorem}
If $\{Y_t\}$ is a stationary $q$-correlated time series with mean $0$, then there exist an uncorrelated stationary sequence $\{\varepsilon_t\}$ and constants $\tilde{\theta}_1,\ldots,\tilde{\theta}_q$ such that, if
\[
X_t
=
\varepsilon_t-\tilde{\theta}_1\varepsilon_{t-1}
-\cdots
-\tilde{\theta}_q\varepsilon_{t-q},
\]
then $\{X_t\}$ has the same autocovariance function as $\{Y_t\}$.
\end{theorem}

\begin{remark}
The theorem is a second-order representation result. It says that, at the level of means and autocovariances, every stationary $q$-correlated process can be matched by an MA($q$)-type construction. It does not say that the full joint distributions must be identical.
\end{remark}

\section{General Linear Processes (GLP)}

\subsection{Definition}

\begin{definition}
Let $\{\varepsilon_t\}$ be IID with
\[
\E[\varepsilon_t]=0,
\qquad
\Var(\varepsilon_t)=\sigma^2.
\]
A process $\{Y_t\}$ is called a general linear process with mean $0$ if
\[
Y_t
=
\sum_{j=-\infty}^{\infty}\psi_j\varepsilon_{t-j},
\]
where the coefficients satisfy
\[
\sum_{j=-\infty}^{\infty}|\psi_j|<\infty.
\]
\end{definition}

The condition
\[
\sum_{j=-\infty}^{\infty}|\psi_j|<\infty
\]
is called absolute summability. It ensures that the infinite linear combination is well behaved.

Explicitly,
\[
Y_t
=
\cdots
+
\psi_{-2}\varepsilon_{t+2}
+
\psi_{-1}\varepsilon_{t+1}
+
\psi_0\varepsilon_t
+
\psi_1\varepsilon_{t-1}
+
\psi_2\varepsilon_{t-2}
+
\cdots.
\]

\begin{definition}
A general linear process is causal, or future-independent, if
\[
\psi_j=0
\qquad
\text{for all } j<0.
\]
In that case,
\[
Y_t
=
\sum_{j=0}^{\infty}\psi_j\varepsilon_{t-j}.
\]
If some $\psi_j\neq 0$ for $j<0$, then the process is noncausal or future-dependent.
\end{definition}

\subsection{Backshift Operator}

\begin{definition}
For a time series $\{Y_t\}$, the backshift operator $B$ is defined by
\[
BY_t=Y_{t-1}.
\]
More generally,
\[
B^jY_t=Y_{t-j}.
\]
\end{definition}

For example,
\[
B^2Y_t
=
B(BY_t)
=
BY_{t-1}
=
Y_{t-2}.
\]
The inverse operator $B^{-1}$ is the forward shift:
\[
B^{-1}Y_t=Y_{t+1},
\qquad
B^{-2}Y_t=Y_{t+2}.
\]

\begin{definition}
A linear filter is an operator of the form
\[
\psi(B)
=
\sum_{j=-\infty}^{\infty}\psi_jB^j.
\]
\end{definition}

Using this notation, a general linear process can be written compactly as
\[
Y_t=\psi(B)\varepsilon_t,
\]
because
\[
\psi(B)\varepsilon_t
=
\left(\sum_{j=-\infty}^{\infty}\psi_jB^j\right)\varepsilon_t
=
\sum_{j=-\infty}^{\infty}\psi_jB^j\varepsilon_t
=
\sum_{j=-\infty}^{\infty}\psi_j\varepsilon_{t-j}.
\]

\subsection*{Stationarity of a General Linear Process}

We now derive the mean, variance, autocovariance function, and autocorrelation function of a general linear process. Let
\[
Y_t
=
\sum_{j=-\infty}^{\infty}\psi_j\varepsilon_{t-j},
\qquad
\sum_{j=-\infty}^{\infty}|\psi_j|<\infty.
\]

The mean function is
\begin{align*}
\mu_t
&=
\E[Y_t] \\
&=
\E\left[
\sum_{j=-\infty}^{\infty}\psi_j\varepsilon_{t-j}
\right] \\
&=
\sum_{j=-\infty}^{\infty}\psi_j\E[\varepsilon_{t-j}] \\
&=
0.
\end{align*}

The variance function is
\begin{align*}
\Var(Y_t)
&=
\Var\left(
\sum_{j=-\infty}^{\infty}\psi_j\varepsilon_{t-j}
\right) \\
&=
\sum_{j=-\infty}^{\infty}\psi_j^2\Var(\varepsilon_{t-j})
+
2\sum_{i<j}\psi_i\psi_j
\Cov(\varepsilon_{t-i},\varepsilon_{t-j}).
\end{align*}
Since the innovations are independent, the covariance terms vanish when $i\neq j$. Therefore,
\[
\Var(Y_t)
=
\sigma^2\sum_{j=-\infty}^{\infty}\psi_j^2.
\]

\begin{exercise}
Show that
\[
\sum_{j=-\infty}^{\infty}|\psi_j|<\infty
\]
implies
\[
\sum_{j=-\infty}^{\infty}\psi_j^2<\infty.
\]
\end{exercise}

\begin{proof}
Since $\sum_{j=-\infty}^{\infty}|\psi_j|<\infty$, we must have $|\psi_j|\to 0$ as $|j|\to\infty$. Hence there exists an integer $N$ such that for all $|j|>N$,
\[
|\psi_j|\leq 1.
\]
For those terms,
\[
\psi_j^2\leq |\psi_j|.
\]
Therefore,
\[
\sum_{|j|>N}\psi_j^2
\leq
\sum_{|j|>N}|\psi_j|
<\infty.
\]
The remaining terms
\[
\sum_{|j|\leq N}\psi_j^2
\]
form a finite sum. Hence
\[
\sum_{j=-\infty}^{\infty}\psi_j^2<\infty.
\]
\end{proof}

Now compute the autocovariance:
\begin{align*}
\Cov(Y_t,Y_{t+k})
&=
\E[Y_tY_{t+k}]
-
\E[Y_t]\E[Y_{t+k}] \\
&=
\E[Y_tY_{t+k}] \\
&=
\E\left[
\left(\sum_{j=-\infty}^{\infty}\psi_j\varepsilon_{t-j}\right)
\left(\sum_{i=-\infty}^{\infty}\psi_i\varepsilon_{t+k-i}\right)
\right] \\
&=
\sum_{j=-\infty}^{\infty}
\sum_{i=-\infty}^{\infty}
\psi_j\psi_i
\E[\varepsilon_{t-j}\varepsilon_{t+k-i}].
\end{align*}
Since $\varepsilon_t$ has mean $0$ and is uncorrelated across time,
\[
\E[\varepsilon_{t-j}\varepsilon_{t+k-i}]
=
\begin{cases}
\sigma^2, & t-j=t+k-i,\\
0, & t-j\neq t+k-i.
\end{cases}
\]
The equality $t-j=t+k-i$ is equivalent to
\[
i=j+k.
\]
Therefore,
\[
\gamma_k
=
\Cov(Y_t,Y_{t+k})
=
\sigma^2
\sum_{j=-\infty}^{\infty}
\psi_j\psi_{j+k}.
\]
This depends only on the lag $k$, not on $t$.

The autocorrelation function is
\[
\rho_k
=
\frac{\gamma_k}{\gamma_0}
=
\frac{\sum_{j=-\infty}^{\infty}\psi_j\psi_{j+k}}
{\sum_{j=-\infty}^{\infty}\psi_j^2}.
\]

Therefore, every general linear process with absolutely summable coefficients is weakly stationary.

\begin{remark}
For a general linear process with mean $\mu$, define
\[
Y_t
=
\mu+\sum_{j=-\infty}^{\infty}\psi_j\varepsilon_{t-j}.
\]
The autocovariance and autocorrelation functions are unchanged by adding the constant mean $\mu$.
\end{remark}

\section{Autoregressive (AR) Processes or order 1}

\subsection{The AR(1) Model}

AR(1) stands for autoregressive time series of order $1$.

\begin{definition}
Let $\phi\in\R$ satisfy $|\phi|<1$, and let
\[
\varepsilon_t \iid \Normal(0,\sigma^2).
\]
An AR(1) process is a stationary solution to
\[
Y_t=\phi Y_{t-1}+\varepsilon_t.
\]
\end{definition}

Equivalently,
\[
Y_t-\phi Y_{t-1}=\varepsilon_t.
\]
Using the backshift operator,
\[
(1-\phi B)Y_t=\varepsilon_t.
\]

\begin{remark}
Some texts write the AR(1) model as
\[
Y_t=\beta_0+\beta_1Y_{t-1}+\varepsilon_t.
\]
The version above has mean $0$. A nonzero intercept is handled later by centering the process around its mean.
\end{remark}

\subsection*{AR(1) as a General Linear Process}

Starting from
\[
Y_t=\phi Y_{t-1}+\varepsilon_t,
\]
substitute recursively:
\begin{align*}
Y_t
&=
\phi(\phi Y_{t-2}+\varepsilon_{t-1})+\varepsilon_t \\
&=
\phi^2Y_{t-2}+\phi\varepsilon_{t-1}+\varepsilon_t \\
&=
\phi^2(\phi Y_{t-3}+\varepsilon_{t-2})+\phi\varepsilon_{t-1}+\varepsilon_t \\
&=
\phi^3Y_{t-3}+\phi^2\varepsilon_{t-2}+\phi\varepsilon_{t-1}+\varepsilon_t.
\end{align*}
Continuing this recursion for $n$ steps gives
\[
Y_t
=
\phi^{n+1}Y_{t-n-1}
+
\sum_{j=0}^{n}\phi^j\varepsilon_{t-j}.
\]
Thus,
\[
Y_t-\sum_{j=0}^{n}\phi^j\varepsilon_{t-j}
=
\phi^{n+1}Y_{t-n-1}.
\]

Now examine the remainder term. Since the process is stationary,
\[
\Var(Y_{t-n-1})=\Var(Y_t)<\infty.
\]
Therefore,
\[
\Var(\phi^{n+1}Y_{t-n-1})
=
\phi^{2n+2}\Var(Y_{t-n-1})
=
\phi^{2n+2}\Var(Y_t).
\]
Because $|\phi|<1$,
\[
\phi^{2n+2}\Var(Y_t)\to 0.
\]
Hence
\[
\phi^{n+1}Y_{t-n-1}\to 0
\]
in mean square, and therefore in probability.

Taking the limit gives
\[
Y_t
=
\sum_{j=0}^{\infty}\phi^j\varepsilon_{t-j}.
\]
Thus the stationary AR(1) process is a causal general linear process with coefficients
\[
\psi_j=
\begin{cases}
\phi^j, & j\geq 0,\\
0, & j<0.
\end{cases}
\]

Since
\[
\sum_{j=0}^{\infty}|\phi|^j
=
\frac{1}{1-|\phi|}
<\infty,
\]
the coefficient sequence is absolutely summable.

\subsection{Moments of the AR(1) Process}

Using the general linear representation,
\[
Y_t=\sum_{j=0}^{\infty}\phi^j\varepsilon_{t-j}.
\]

The mean is
\[
\E[Y_t]
=
\sum_{j=0}^{\infty}\phi^j\E[\varepsilon_{t-j}]
=
0.
\]

The variance is
\begin{align*}
\Var(Y_t)
&=
\sigma^2\sum_{j=0}^{\infty}\phi^{2j} \\
&=
\sigma^2\frac{1}{1-\phi^2}.
\end{align*}
Thus
\[
\gamma_0
=
\frac{\sigma^2}{1-\phi^2}.
\]

For $k\geq 0$, the autocovariance is
\begin{align*}
\gamma_k
&=
\Cov(Y_t,Y_{t+k}) \\
&=
\sigma^2\sum_{j=0}^{\infty}\psi_j\psi_{j+k} \\
&=
\sigma^2\sum_{j=0}^{\infty}\phi^j\phi^{j+k} \\
&=
\sigma^2\phi^k\sum_{j=0}^{\infty}\phi^{2j} \\
&=
\frac{\sigma^2\phi^k}{1-\phi^2}.
\end{align*}
Therefore,
\[
\gamma_k=\phi^k\gamma_0,
\qquad
k\geq 0.
\]
By symmetry,
\[
\gamma_k=\phi^{|k|}\gamma_0.
\]

The autocorrelation function is
\[
\rho_k
=
\frac{\gamma_k}{\gamma_0}
=
\phi^{|k|}.
\]

\begin{remark}
The stationary AR(1) process is not $q$-dependent for any finite $q$ when $\phi\neq 0$, because
\[
\gamma_k=\phi^{|k|}\gamma_0
\]
is nonzero for every $k$. However, its dependence decays geometrically when $|\phi|<1$.
\end{remark}

\begin{lstlisting}[language=R, caption={Simulating an AR(1) process in R}]
set.seed(123)

n <- 300
phi <- 0.7
sigma <- 1

eps <- rnorm(n, mean = 0, sd = sigma)
y <- numeric(n)

for (t in 2:n) {
  y[t] <- phi * y[t - 1] + eps[t]
}

plot(
  y, type = "l",
  xlab = "t", ylab = "Y_t",
  main = "Simulated AR(1) Process"
)

acf(y, main = "Sample ACF of AR(1)")
\end{lstlisting}

\subsection*{More on AR(1)}

Last time, we defined an autoregressive process of order $1$ by
\[
Y_t=\phi Y_{t-1}+\varepsilon_t,
\qquad
\varepsilon_t \iid \Normal(0,\sigma_\varepsilon^2),
\]
and derived the general linear process representation
\[
Y_t=\sum_{j=0}^{\infty}\phi^j\varepsilon_{t-j},
\]
valid when $|\phi|<1$.

\subsubsection*{Applying the GLP Results to AR(1)}

Assume $|\phi|<1$. Since the AR(1) process has GLP coefficients
\[
\psi_j=\phi^j,
\qquad j\geq 0,
\]
we can compute its moments directly.

The mean function is
\[
\mu_t=\E[Y_t]=0.
\]

The variance function is
\begin{align*}
\sigma_Y^2
=
\Var(Y_t)
&=
\sigma_\varepsilon^2\sum_{j=0}^{\infty}\psi_j^2 \\
&=
\sigma_\varepsilon^2\sum_{j=0}^{\infty}\phi^{2j} \\
&=
\frac{\sigma_\varepsilon^2}{1-\phi^2}.
\end{align*}

The autocovariance function is
\begin{align*}
\gamma_k
=
\Cov(Y_t,Y_{t-k})
&=
\sigma_\varepsilon^2\sum_{j=0}^{\infty}\psi_j\psi_{j+k} \\
&=
\sigma_\varepsilon^2\sum_{j=0}^{\infty}\phi^j\phi^{j+k} \\
&=
\sigma_\varepsilon^2\phi^k\sum_{j=0}^{\infty}\phi^{2j} \\
&=
\frac{\sigma_\varepsilon^2\phi^k}{1-\phi^2}.
\end{align*}
Therefore,
\[
\gamma_k=\phi^k\gamma_0,
\qquad k\geq 0.
\]
By symmetry,
\[
\gamma_k=\phi^{|k|}\gamma_0.
\]

The autocorrelation function is
\[
\rho_k=\frac{\gamma_k}{\gamma_0}=\phi^{|k|}.
\]

\begin{remark}
Since $\rho_k\to 0$ geometrically as $k\to\infty$, the dependence in a stationary AR(1) process decays geometrically. However, if $\phi\neq 0$, then $\rho_k\neq 0$ for every finite $k$, so the AR(1) process is not $q$-dependent for any finite $q$.
\end{remark}

\subsubsection*{Using the Backshift Operator}

Using the backshift operator $B$, the AR(1) equation can be written as
\[
Y_t=\phi Y_{t-1}+\varepsilon_t
\iff
Y_t-\phi BY_t=\varepsilon_t
\iff
(1-\phi B)Y_t=\varepsilon_t.
\]
Formally,
\[
Y_t=(1-\phi B)^{-1}\varepsilon_t.
\]
When $|\phi|<1$, we use the geometric expansion
\[
(1-\phi B)^{-1}
=
1+\phi B+\phi^2B^2+\phi^3B^3+\cdots.
\]
Therefore,
\begin{align*}
Y_t
&=
(1+\phi B+\phi^2B^2+\phi^3B^3+\cdots)\varepsilon_t \\
&=
\varepsilon_t+\phi\varepsilon_{t-1}+\phi^2\varepsilon_{t-2}
+\phi^3\varepsilon_{t-3}+\cdots \\
&=
\sum_{j=0}^{\infty}\phi^j\varepsilon_{t-j}.
\end{align*}
This agrees with the GLP representation derived by recursive substitution.

\subsubsection{The Case $|\phi|>1$}

The condition $|\phi|<1$ gives the causal stationary representation
\[
Y_t=\sum_{j=0}^{\infty}\phi^j\varepsilon_{t-j}.
\]
Now suppose $|\phi|>1$. In this case, the causal expansion does not converge because
\[
\sum_{j=0}^{\infty}|\phi|^j=\infty.
\]
However, we can rewrite
\[
(1-\phi B)Y_t=\varepsilon_t
\]
as
\[
Y_t=(1-\phi B)^{-1}\varepsilon_t.
\]
For $|\phi|>1$, factor the operator differently:
\[
1-\phi B
=
-\phi B\left(1-\frac{1}{\phi}B^{-1}\right).
\]
Hence
\[
(1-\phi B)^{-1}
=
-\frac{1}{\phi}B^{-1}
\left(1-\frac{1}{\phi}B^{-1}\right)^{-1}.
\]
Since $|1/\phi|<1$,
\[
\left(1-\frac{1}{\phi}B^{-1}\right)^{-1}
=
\sum_{j=0}^{\infty}\phi^{-j}B^{-j}.
\]
Therefore,
\begin{align*}
Y_t
&=
-\frac{1}{\phi}B^{-1}
\sum_{j=0}^{\infty}\phi^{-j}B^{-j}\varepsilon_t \\
&=
-\sum_{j=0}^{\infty}\phi^{-(j+1)}B^{-(j+1)}\varepsilon_t \\
&=
-\sum_{j=1}^{\infty}\phi^{-j}\varepsilon_{t+j}.
\end{align*}
This is a noncausal GLP representation because it depends on future innovations.

\begin{remark}
For $|\phi|>1$, the AR(1) equation has a stationary solution, but that solution is noncausal:
\[
Y_t=-\sum_{j=1}^{\infty}\phi^{-j}\varepsilon_{t+j}.
\]
It can still be written as a GLP with absolutely summable coefficients because $\sum_{j=1}^{\infty}|\phi|^{-j}<\infty$.
\end{remark}

In summary, for AR(1):
\begin{itemize}
\item If $|\phi|<1$, the process has a causal stationary GLP representation.
\item If $|\phi|>1$, the process has a noncausal stationary GLP representation.
\item If $|\phi|=1$, the process is not stationary.
\end{itemize}

\section{Wold's Decomposition Theorem}

\begin{theorem}[Wold's Decomposition]
Any covariance stationary time series can be represented as the sum of a general linear process and a deterministic component.
\end{theorem}

The theorem says that, under stationarity, the purely nondeterministic part of a time series can be represented as an infinite moving average of white noise innovations. This provides the theoretical motivation for studying ARMA and general linear process representations.

\section{Yule--Walker Method for AR(1)}

The Yule--Walker method finds the autocovariance and autocorrelation functions of an autoregressive process directly from the defining equation. For this method, we assume the process is mean zero, stationary, and causal.

\subsection*{Yule--Walker Equations for AR(1)}

Consider the causal AR(1) process
\[
Y_t-\phi Y_{t-1}=\varepsilon_t.
\]
Multiply both sides by $Y_{t-k}$, where $k\geq 0$:
\[
Y_tY_{t-k}-\phi Y_{t-1}Y_{t-k}
=
\varepsilon_tY_{t-k}.
\]
Taking expectations gives
\[
\E[Y_tY_{t-k}]
-
\phi\E[Y_{t-1}Y_{t-k}]
=
\E[\varepsilon_tY_{t-k}].
\]
Since the process is mean zero and stationary,
\[
\E[Y_tY_{t-k}]=\gamma_k,
\]
and
\[
\E[Y_{t-1}Y_{t-k}]
=
\Cov(Y_{t-1},Y_{t-k})
=
\gamma_{k-1}.
\]
Thus,
\[
\gamma_k-\phi\gamma_{k-1}
=
\E[\varepsilon_tY_{t-k}].
\]

For $k=0$,
\[
\E[\varepsilon_tY_t]
=
\E[\varepsilon_t(\phi Y_{t-1}+\varepsilon_t)]
=
\phi\E[\varepsilon_tY_{t-1}]+\E[\varepsilon_t^2].
\]
By causality, $Y_{t-1}$ depends only on past innovations, so it is independent of $\varepsilon_t$. Therefore,
\[
\E[\varepsilon_tY_{t-1}]=0,
\]
and
\[
\E[\varepsilon_tY_t]=\sigma_\varepsilon^2.
\]
Hence the $k=0$ Yule--Walker equation is
\[
\gamma_0-\phi\gamma_1=\sigma_\varepsilon^2.
\]

For $k\geq 1$, $Y_{t-k}$ depends only on innovations up to time $t-k$, so it is independent of $\varepsilon_t$. Therefore,
\[
\E[\varepsilon_tY_{t-k}]=0,
\]
and
\[
\gamma_k-\phi\gamma_{k-1}=0.
\]
The Yule--Walker equations are therefore
\[
\begin{cases}
\gamma_0-\phi\gamma_1=\sigma_\varepsilon^2, \\
\gamma_1-\phi\gamma_0=0, \\
\gamma_2-\phi\gamma_1=0, \\
\gamma_3-\phi\gamma_2=0, \\
\cdots
\end{cases}
\]
From the second equation,
\[
\gamma_1=\phi\gamma_0.
\]
Substituting into the first equation,
\[
\gamma_0-\phi^2\gamma_0=\sigma_\varepsilon^2.
\]
Therefore,
\[
\gamma_0=\frac{\sigma_\varepsilon^2}{1-\phi^2}.
\]
The recursive equations then give
\[
\gamma_k=\phi^k\gamma_0
=
\frac{\sigma_\varepsilon^2\phi^k}{1-\phi^2},
\qquad k\geq 0.
\]
Hence
\[
\rho_k=\phi^k,
\qquad k\geq 0.
\]

\begin{remark}
The Yule--Walker derivation uses causality. Without causality, the term $\E[\varepsilon_tY_{t-k}]$ may not vanish for $k\geq 1$ because $Y_{t-k}$ could depend on future innovations.
\end{remark}

\section*{AR(1) with Mean $\mu$}

Suppose an AR(1) process has mean $\mu$. Then $Y_t-\mu$ is a mean-zero AR(1) process:
\[
Y_t-\mu=\phi(Y_{t-1}-\mu)+\varepsilon_t.
\]
Expanding,
\[
Y_t-\mu=\phi Y_{t-1}-\phi\mu+\varepsilon_t,
\]
so
\[
Y_t=(\mu-\phi\mu)+\phi Y_{t-1}+\varepsilon_t.
\]
Thus
\[
Y_t=c+\phi Y_{t-1}+\varepsilon_t,
\]
where
\[
c=\mu(1-\phi).
\]
Equivalently,
\[
\mu=\frac{c}{1-\phi}.
\]

\section{Autoregressive Processes (AR) of Order $p$}

\subsection{Definition}

\begin{definition}
A time series $\{Y_t\}$ is an AR($p$) process with mean zero if it is a stationary solution to
\[
Y_t
=
\phi_1Y_{t-1}
+
\phi_2Y_{t-2}
+\cdots+
\phi_pY_{t-p}
+\varepsilon_t,
\]
where
\[
\varepsilon_t \iid \Normal(0,\sigma_\varepsilon^2).
\]
\end{definition}

Equivalently,
\[
Y_t-\phi_1Y_{t-1}-\cdots-\phi_pY_{t-p}
=
\varepsilon_t.
\]

\begin{remark}
This convention uses plus signs in the time-domain equation. Some books and software write the coefficients with the opposite sign.
\end{remark}

\subsection{AR Polynomial}

\begin{definition}
The AR polynomial for an AR($p$) process is
\[
\Phi(z)
=
1-\phi_1z-\phi_2z^2-\cdots-\phi_pz^p.
\]
\end{definition}

Using the backshift operator,
\[
\Phi(B)
=
1-\phi_1B-\phi_2B^2-\cdots-\phi_pB^p.
\]
Therefore, the AR($p$) equation can be written as
\[
\Phi(B)Y_t=\varepsilon_t.
\]

\subsection{Causality Condition for AR($p$)}

\begin{definition}
An AR($p$) process is causal if it has a representation
\[
Y_t
=
\sum_{j=0}^{\infty}\psi_j\varepsilon_{t-j},
\]
where
\[
\sum_{j=0}^{\infty}|\psi_j|<\infty.
\]
\end{definition}

\begin{theorem}[Causality Condition]
An AR($p$) process is causal if and only if all complex roots of the AR polynomial
\[
\Phi(z)=1-\phi_1z-\cdots-\phi_pz^p
\]
lie outside the unit circle. That is, if $z_1,\ldots,z_p$ are the roots of $\Phi(z)$, then
\[
|z_i|>1
\qquad
\text{for all } i=1,\ldots,p.
\]
\end{theorem}

For $p=1$, the AR polynomial is
\[
\Phi(z)=1-\phi z.
\]
The root is
\[
z=\frac{1}{\phi}.
\]
The causality condition requires
\[
\left|\frac{1}{\phi}\right|>1,
\]
which is equivalent to
\[
|\phi|<1.
\]
Thus the AR($p$) causality condition reduces to the usual AR(1) condition when $p=1$.

\subsection*{Deriving the GLP Form for Causal AR($p$)}

Suppose the AR polynomial has roots $z_1,\ldots,z_p$, all satisfying
\[
|z_i|>1.
\]
Assume for simplicity that the roots are distinct. Then
\[
\Phi(z)
=
1-\phi_1z-\cdots-\phi_pz^p
=
\prod_{i=1}^p\left(1-\frac{z}{z_i}\right).
\]
Therefore,
\[
\Phi(B)
=
\prod_{i=1}^p\left(1-\frac{B}{z_i}\right).
\]
Since $|1/z_i|<1$, each inverse factor has the convergent expansion
\[
\left(1-\frac{B}{z_i}\right)^{-1}
=
\sum_{j=0}^{\infty}z_i^{-j}B^j.
\]
Thus
\[
Y_t
=
\Phi(B)^{-1}\varepsilon_t
=
\left[
\prod_{i=1}^p
\left(1-\frac{B}{z_i}\right)^{-1}
\right]\varepsilon_t.
\]
Expanding the product gives
\[
Y_t
=
\left[
\prod_{i=1}^p
\left(\sum_{j=0}^{\infty}z_i^{-j}B^j\right)
\right]\varepsilon_t.
\]
This has the form
\[
Y_t=\sum_{n=0}^{\infty}\psi_n\varepsilon_{t-n}.
\]

To see the coefficient formula, first consider $p=2$. Then
\begin{align*}
\left(\sum_{j=0}^{\infty}a_jB^j\right)
\left(\sum_{j=0}^{\infty}b_jB^j\right)
&=
a_0b_0+(a_0b_1+a_1b_0)B \\
&\quad+
(a_0b_2+a_1b_1+a_2b_0)B^2+\cdots \\
&=
\sum_{n=0}^{\infty}
\left(
\sum_{j=0}^{n}a_jb_{n-j}
\right)B^n.
\end{align*}
For general $p$, the product is a $p$-fold convolution:
\[
\prod_{i=1}^p
\left(\sum_{j=0}^{\infty}z_i^{-j}B^j\right)
=
\sum_{n=0}^{\infty}\psi_nB^n,
\]
where
\[
\psi_n
=
\sum_{j_1+\cdots+j_p=n}
z_1^{-j_1}z_2^{-j_2}\cdots z_p^{-j_p}.
\]
Thus,
\[
Y_t
=
\sum_{n=0}^{\infty}\psi_n\varepsilon_{t-n}.
\]

\subsection*{Sketch That the GLP Coefficients Are Absolutely Summable}

We need to show that
\[
\sum_{n=0}^{\infty}|\psi_n|<\infty.
\]
Using the coefficient formula,
\[
|\psi_n|
\leq
\sum_{j_1+\cdots+j_p=n}
|z_1|^{-j_1}\cdots |z_p|^{-j_p}.
\]
Let $z_\ast$ be the root with smallest modulus. Since all roots are outside the unit circle,
\[
|z_\ast|>1.
\]
Therefore,
\[
|z_i|^{-j_i}\leq |z_\ast|^{-j_i}
\]
for every $i$. Hence
\[
|\psi_n|
\leq
\sum_{j_1+\cdots+j_p=n}
|z_\ast|^{-n}.
\]
The number of nonnegative integer solutions to
\[
j_1+\cdots+j_p=n
\]
is
\[
\binom{n+p-1}{p-1},
\]
which is a polynomial in $n$ of degree $p-1$. Therefore, there exists a constant $C(p)$ such that
\[
|\psi_n|
\leq
C(p)n^{p-1}|z_\ast|^{-n}.
\]
Since $|z_\ast|>1$,
\[
\sum_{n=0}^{\infty}n^{p-1}|z_\ast|^{-n}<\infty.
\]
Thus
\[
\sum_{n=0}^{\infty}|\psi_n|<\infty.
\]
So the causal AR($p$) process is a valid general linear process.

\subsection*{Other Root Cases}

For AR($p$), the following cases are important:

\begin{itemize}
\item If all roots of $\Phi(z)$ lie outside the unit circle, then $\{Y_t\}$ is causal and stationary.
\item If at least one root lies inside the unit circle, and no roots lie on the unit circle, then $\{Y_t\}$ has a stationary noncausal representation.
\item If at least one root lies on the unit circle, then $\{Y_t\}$ is not stationary.
\end{itemize}

\subsection{Example: AR(2)}

Consider the AR(2) process
\[
Y_t=Y_{t-1}+Y_{t-2}+\varepsilon_t.
\]
The AR polynomial is
\[
\Phi(z)=1-z-z^2.
\]
Solving $\Phi(z)=0$ gives
\[
1-z-z^2=0
\iff
z^2+z-1=0,
\]
so
\[
z_{1,2}
=
\frac{-1\pm \sqrt{5}}{2}.
\]
One root has absolute value greater than $1$, while the other has absolute value less than $1$. Therefore, this AR(2) process is stationary but noncausal.

For a general AR(2) process
\[
Y_t=\phi_1Y_{t-1}+\phi_2Y_{t-2}+\varepsilon_t,
\]
the AR polynomial is
\[
\Phi(z)=1-\phi_1z-\phi_2z^2.
\]
The roots are
\[
z_{1,2}
=
\frac{-\phi_1\pm \sqrt{\phi_1^2+4\phi_2}}{2\phi_2},
\]
when $\phi_2\neq 0$. The causality condition can be equivalently written as
\[
|\phi_2|<1,
\qquad
\phi_1+\phi_2<1,
\qquad
\phi_2-\phi_1<1.
\]

\section{Finding the GLP Coefficients for AR($p$)}

\subsection{Method 1: Using Convolution}

Suppose the roots of the AR polynomial are known and are outside the unit circle. Then the GLP coefficients are
\[
\psi_n
=
\sum_{j_1+\cdots+j_p=n}
z_1^{-j_1}\cdots z_p^{-j_p}.
\]

\begin{example}
Consider the AR(2) process
\[
Y_t=\frac{1}{6}Y_{t-1}+\frac{1}{6}Y_{t-2}+\varepsilon_t.
\]
The AR polynomial is
\[
\Phi(z)=1-\frac{1}{6}z-\frac{1}{6}z^2.
\]
Solving $\Phi(z)=0$ gives roots
\[
z_1=-3,
\qquad
z_2=2.
\]
Both roots satisfy $|z_i|>1$, so the process is causal.

The GLP coefficients are
\[
\psi_n
=
\sum_{j=0}^{n}z_1^{-j}z_2^{-(n-j)}.
\]
Thus,
\[
\psi_0=1,
\]
\[
\psi_1
=
z_1^{-1}+z_2^{-1}
=
-\frac{1}{3}+\frac{1}{2}
=
\frac{1}{6},
\]
and
\[
\psi_2
=
z_1^{-2}+z_1^{-1}z_2^{-1}+z_2^{-2}
=
\frac{1}{9}-\frac{1}{6}+\frac{1}{4}
=
\frac{7}{36}.
\]
\end{example}

\subsection{Method 2: Using the AR($p$) Equation}

For the same example, suppose
\[
Y_t
=
\frac{1}{6}Y_{t-1}
+
\frac{1}{6}Y_{t-2}
+
\varepsilon_t.
\]
Assume the GLP representation
\[
Y_t
=
\psi_0\varepsilon_t
+
\psi_1\varepsilon_{t-1}
+
\psi_2\varepsilon_{t-2}
+\cdots.
\]
Substitute this into the AR equation:
\begin{align*}
\psi_0\varepsilon_t+\psi_1\varepsilon_{t-1}+\psi_2\varepsilon_{t-2}+\cdots
&=
\frac{1}{6}
(\psi_0\varepsilon_{t-1}+\psi_1\varepsilon_{t-2}+\psi_2\varepsilon_{t-3}+\cdots) \\
&\quad+
\frac{1}{6}
(\psi_0\varepsilon_{t-2}+\psi_1\varepsilon_{t-3}+\psi_2\varepsilon_{t-4}+\cdots)
+
\varepsilon_t.
\end{align*}
Matching coefficients of each innovation gives
\[
\psi_0=1,
\]
\[
\psi_1=\frac{1}{6}\psi_0,
\]
and for $n\geq 2$,
\[
\psi_n
=
\frac{1}{6}\psi_{n-1}
+
\frac{1}{6}\psi_{n-2}.
\]
Therefore,
\[
\psi_0=1,
\qquad
\psi_1=\frac{1}{6},
\]
and
\[
\psi_2
=
\frac{1}{6}\psi_1+\frac{1}{6}\psi_0
=
\frac{1}{36}+\frac{1}{6}
=
\frac{7}{36}.
\]
This agrees with the convolution method.

\section{Yule--Walker Method for AR($p$)}

The Yule--Walker method generalizes from AR(1) to AR($p$). We assume the process is mean zero, stationary, and causal.

The AR($p$) equation is
\[
Y_t-\phi_1Y_{t-1}-\cdots-\phi_pY_{t-p}
=
\varepsilon_t.
\]
Multiply both sides by $Y_{t-k}$, where $k\geq 0$:
\[
Y_tY_{t-k}
-
\phi_1Y_{t-1}Y_{t-k}
-\cdots
-
\phi_pY_{t-p}Y_{t-k}
=
\varepsilon_tY_{t-k}.
\]
Taking expectations gives
\[
\E[Y_tY_{t-k}]
-
\phi_1\E[Y_{t-1}Y_{t-k}]
-\cdots
-
\phi_p\E[Y_{t-p}Y_{t-k}]
=
\E[\varepsilon_tY_{t-k}].
\]
By stationarity,
\[
\E[Y_tY_{t-k}]=\gamma_k,
\]
and
\[
\E[Y_{t-j}Y_{t-k}]
=
\gamma_{k-j}.
\]
Thus,
\[
\gamma_k-\phi_1\gamma_{k-1}-\cdots-\phi_p\gamma_{k-p}
=
\E[\varepsilon_tY_{t-k}].
\]

By causality,
\[
\E[\varepsilon_tY_{t-k}]
=
\begin{cases}
\sigma_\varepsilon^2, & k=0,\\
0, & k\geq 1.
\end{cases}
\]
Therefore,
\[
\gamma_k-\phi_1\gamma_{k-1}-\cdots-\phi_p\gamma_{k-p}
=
\begin{cases}
\sigma_\varepsilon^2, & k=0,\\
0, & k\geq 1.
\end{cases}
\]

Using $\gamma_{-k}=\gamma_k$, the first $p+1$ equations are
\[
\begin{cases}
\gamma_0-\phi_1\gamma_1-\phi_2\gamma_2-\cdots-\phi_p\gamma_p
=
\sigma_\varepsilon^2, \\
\gamma_1-\phi_1\gamma_0-\phi_2\gamma_1-\cdots-\phi_p\gamma_{p-1}
=
0, \\
\gamma_2-\phi_1\gamma_1-\phi_2\gamma_0-\cdots-\phi_p\gamma_{p-2}
=
0, \\
\vdots \\
\gamma_p-\phi_1\gamma_{p-1}-\phi_2\gamma_{p-2}-\cdots-\phi_p\gamma_0
=
0.
\end{cases}
\]
This is a system of $p+1$ equations in the $p+1$ unknowns
\[
\gamma_0,\gamma_1,\ldots,\gamma_p.
\]

For $k\geq p$, the recursion simplifies to
\[
\gamma_k
=
\phi_1\gamma_{k-1}
+
\phi_2\gamma_{k-2}
+\cdots+
\phi_p\gamma_{k-p}.
\]
Thus, once $\gamma_0,\ldots,\gamma_p$ are known, the remaining autocovariances can be computed recursively.

\begin{exercise}
Verify that
\[
\E[\varepsilon_tY_{t-k}]
=
\begin{cases}
\sigma_\varepsilon^2, & k=0,\\
0, & k\geq 1.
\end{cases}
\]
\end{exercise}

\begin{proof}
For $k=0$,
\[
Y_t
=
\phi_1Y_{t-1}
+\cdots+
\phi_pY_{t-p}
+\varepsilon_t.
\]
Therefore,
\begin{align*}
\E[\varepsilon_tY_t]
&=
\E\left[
\varepsilon_t
\left(
\phi_1Y_{t-1}
+\cdots+
\phi_pY_{t-p}
+\varepsilon_t
\right)
\right] \\
&=
\sum_{j=1}^p\phi_j\E[\varepsilon_tY_{t-j}]
+
\E[\varepsilon_t^2].
\end{align*}
Because the process is causal, each $Y_{t-j}$ depends only on innovations dated $t-j,t-j-1,\ldots$, and is independent of $\varepsilon_t$. Hence
\[
\E[\varepsilon_tY_{t-j}]=0,
\qquad j=1,\ldots,p.
\]
Thus
\[
\E[\varepsilon_tY_t]=\sigma_\varepsilon^2.
\]

For $k\geq 1$, $Y_{t-k}$ depends only on innovations dated $t-k,t-k-1,\ldots$, all of which occur before time $t$. By causality, $Y_{t-k}$ is independent of $\varepsilon_t$, so
\[
\E[\varepsilon_tY_{t-k}]
=
\E[\varepsilon_t]\E[Y_{t-k}]
=
0.
\]
\end{proof}

\subsection{A Property of the Recursion}

For large $k$, computing $\gamma_k$ recursively from
\[
\gamma_k
=
\phi_1\gamma_{k-1}
+
\phi_2\gamma_{k-2}
+\cdots+
\phi_p\gamma_{k-p}
\]
may be difficult. The recursion has a useful closed-form structure.

\begin{proposition}
If the roots of the AR polynomial $\Phi(z)$ are distinct, then there exist constants $A_1,\ldots,A_p$ such that
\[
\gamma_k
=
A_1z_1^{-k}+A_2z_2^{-k}+\cdots+A_pz_p^{-k},
\qquad k\geq 0.
\]
\end{proposition}

\begin{proof}
For $k\geq p$, the autocovariances satisfy the linear difference equation
\[
\gamma_k
-
\phi_1\gamma_{k-1}
-\cdots
-
\phi_p\gamma_{k-p}
=
0.
\]
Look for solutions of the form
\[
\gamma_k=r^k.
\]
Substitution gives
\[
r^k-\phi_1r^{k-1}-\cdots-\phi_pr^{k-p}=0.
\]
Dividing by $r^{k-p}$ gives
\[
r^p-\phi_1r^{p-1}-\cdots-\phi_p=0.
\]
This is the characteristic equation of the recursion. Its roots are
\[
r_i=\frac{1}{z_i},
\]
where $z_i$ are the roots of the AR polynomial $\Phi(z)$. Therefore, when the roots are distinct, the general solution is
\[
\gamma_k
=
A_1r_1^k+\cdots+A_pr_p^k
=
A_1z_1^{-k}+\cdots+A_pz_p^{-k}.
\]
The constants $A_1,\ldots,A_p$ are determined by the initial values $\gamma_0,\ldots,\gamma_{p-1}$.
\end{proof}

\begin{remark}
For a causal AR($p$) process,
\[
|z_i|>1
\]
for every root. Therefore,
\[
|z_i|^{-k}\to 0
\]
as $k\to\infty$. This gives another explanation for why the autocovariance function of a causal AR($p$) process decays exponentially.
\end{remark}

\subsection{Asymptotic Decay of the ACF}

For a causal AR($p$) process, the autocovariance and autocorrelation functions decay exponentially. Up to a constant factor,
\[
\gamma_k
\approx
\left(\min\{|z_1|,\ldots,|z_p|\}\right)^{-k},
\qquad k\to\infty.
\]
The root closest to the unit circle controls the slowest rate of decay.

For example, if an AR(3) polynomial has roots
\[
z_1=2,
\qquad
z_2=3,
\qquad
z_3=10,
\]
then the closest root to the unit circle is $z_1=2$. Therefore,
\[
\gamma_k
\approx C\cdot 2^{-k}
\]
for large $k$, where $C$ is a constant depending on the initial autocovariances and model parameters.

\begin{remark}
The ACF of a causal AR($p$) process generally never becomes exactly zero after a finite lag. Instead, it decays toward zero. This contrasts with an MA($q$) process, whose ACF is exactly zero after lag $q$.
\end{remark}

\section{More on Autocovariance Functions}

\subsection{ACVF for AR($p$)}

Suppose $z_1,\ldots,z_p$ are the distinct roots of the AR polynomial
\[
\Phi(z)=1-\phi_1z-\phi_2z^2-\cdots-\phi_pz^p,
\]
and suppose the AR($p$) process is causal. Last time, we claimed that there exist constants $A_1,\ldots,A_p$ such that the solution to the recursion
\[
\gamma_k
=
\phi_1\gamma_{k-1}
+
\phi_2\gamma_{k-2}
+\cdots+
\phi_p\gamma_{k-p},
\qquad k\geq p,
\]
is
\[
\gamma_k
=
A_1z_1^{-k}+A_2z_2^{-k}+\cdots+A_pz_p^{-k},
\qquad k\geq 0.
\]

We verify that this expression satisfies the recursion. Substitute
\[
\gamma_k
=
A_1z_1^{-k}+\cdots+A_pz_p^{-k}
\]
into the right-hand side:
\begin{align*}
&\phi_1\gamma_{k-1}
+
\phi_2\gamma_{k-2}
+\cdots+
\phi_p\gamma_{k-p} \\
&=
\phi_1\left(A_1z_1^{-(k-1)}+\cdots+A_pz_p^{-(k-1)}\right)
+
\cdots
+
\phi_p\left(A_1z_1^{-(k-p)}+\cdots+A_pz_p^{-(k-p)}\right) \\
&=
A_1z_1^{-k}
\left(\phi_1z_1+\phi_2z_1^2+\cdots+\phi_pz_1^p\right)
+
\cdots
+
A_pz_p^{-k}
\left(\phi_1z_p+\phi_2z_p^2+\cdots+\phi_pz_p^p\right).
\end{align*}
Since each $z_i$ is a root of $\Phi(z)$,
\[
1-\phi_1z_i-\phi_2z_i^2-\cdots-\phi_pz_i^p=0,
\]
so
\[
\phi_1z_i+\phi_2z_i^2+\cdots+\phi_pz_i^p=1.
\]
Therefore,
\[
\phi_1\gamma_{k-1}
+
\cdots+
\phi_p\gamma_{k-p}
=
A_1z_1^{-k}+\cdots+A_pz_p^{-k}
=
\gamma_k.
\]
Thus the proposed expression satisfies the recursion.

The constants $A_1,\ldots,A_p$ are determined from the initial autocovariances
\[
\gamma_0,\gamma_1,\ldots,\gamma_{p-1}.
\]
Specifically, they solve the linear system
\[
\begin{bmatrix}
1 & 1 & \cdots & 1 \\
z_1^{-1} & z_2^{-1} & \cdots & z_p^{-1} \\
z_1^{-2} & z_2^{-2} & \cdots & z_p^{-2} \\
\vdots & \vdots & \ddots & \vdots \\
z_1^{-(p-1)} & z_2^{-(p-1)} & \cdots & z_p^{-(p-1)}
\end{bmatrix}
\begin{bmatrix}
A_1\\
A_2\\
\vdots\\
A_p
\end{bmatrix}
=
\begin{bmatrix}
\gamma_0\\
\gamma_1\\
\vdots\\
\gamma_{p-1}
\end{bmatrix}.
\]

\subsection{ACVF for AR(2)}

Consider the AR(2) process
\[
Y_t=\phi_1Y_{t-1}+\phi_2Y_{t-2}+\varepsilon_t,
\]
where $\phi_2\neq 0$. The AR polynomial is
\[
\Phi(z)=1-\phi_1z-\phi_2z^2.
\]
Solving $\Phi(z)=0$ gives
\[
z_{1,2}
=
\frac{-\phi_1\pm\sqrt{\phi_1^2+4\phi_2}}{2\phi_2}.
\]

There are three cases for the roots:

\begin{enumerate}
\item Two distinct real roots.
\item One repeated real root.
\item Two complex conjugate roots.
\end{enumerate}

\subsubsection{Case 1: Distinct Real Roots}

If $z_1,z_2\in\R$ and $z_1\neq z_2$, then
\[
\gamma_k
=
A_1z_1^{-k}+A_2z_2^{-k},
\qquad k\geq 0.
\]
The constants $A_1,A_2$ are determined by $\gamma_0$ and $\gamma_1$:
\[
A_1+A_2=\gamma_0,
\]
\[
A_1z_1^{-1}+A_2z_2^{-1}=\gamma_1.
\]

\subsubsection{Case 2: Repeated Real Root}

If $z_1=z_2=z\in\R$, then the solution has the form
\[
\gamma_k
=
(A_1+A_2k)z^{-k},
\qquad k\geq 0.
\]

\begin{exercise}
Verify that
\[
\gamma_k=(A_1+A_2k)z^{-k}
\]
satisfies the AR(2) autocovariance recursion.
\end{exercise}

\begin{proof}
The AR(2) autocovariance recursion is
\[
\gamma_k=\phi_1\gamma_{k-1}+\phi_2\gamma_{k-2}.
\]
A repeated root $z$ of
\[
\Phi(z)=1-\phi_1z-\phi_2z^2
\]
means
\[
\Phi(z)=0
\]
and
\[
\Phi'(z)=0.
\]
That is,
\[
1-\phi_1z-\phi_2z^2=0
\]
and
\[
-\phi_1-2\phi_2z=0.
\]
The second equation gives
\[
\phi_1=-2\phi_2z.
\]
Substituting into the first equation,
\[
1+2\phi_2z^2-\phi_2z^2=0,
\]
so
\[
\phi_2=-\frac{1}{z^2},
\qquad
\phi_1=\frac{2}{z}.
\]

Now let
\[
\gamma_k=(A_1+A_2k)z^{-k}.
\]
Then
\[
\gamma_{k-1}=(A_1+A_2(k-1))z^{-(k-1)}
\]
and
\[
\gamma_{k-2}=(A_1+A_2(k-2))z^{-(k-2)}.
\]
Therefore,
\begin{align*}
\phi_1\gamma_{k-1}+\phi_2\gamma_{k-2}
&=
\frac{2}{z}(A_1+A_2(k-1))z^{-(k-1)}
-\frac{1}{z^2}(A_1+A_2(k-2))z^{-(k-2)} \\
&=
2(A_1+A_2k-A_2)z^{-k}
-
(A_1+A_2k-2A_2)z^{-k} \\
&=
(A_1+A_2k)z^{-k} \\
&=
\gamma_k.
\end{align*}
Thus the proposed form satisfies the recursion.
\end{proof}

If the causality condition holds, then $|z|>1$, so $z^{-k}\to 0$ exponentially as $k\to\infty$. The factor $k$ grows only polynomially, so
\[
(A_1+A_2k)z^{-k}\to 0.
\]

\subsubsection{Case 3: Complex Conjugate Roots}

Suppose
\[
z_1,z_2\notin\R,
\qquad
z_2=\overline{z_1}.
\]
Since $\phi_1,\phi_2\in\R$, the roots must occur as a complex conjugate pair. In this case, the autocovariance function can be written as
\[
\gamma_k
=
A_1z_1^{-k}+A_2z_2^{-k},
\qquad k\geq 0,
\]
where $A_2=\overline{A_1}$ so that $\gamma_k$ is real.

The autocorrelation function has a damped sinusoidal form:
\[
\rho_k
=
R^k\frac{\sin(k\theta+\Phi)}{\sin(\Phi)},
\]
where
\[
R=\frac{1}{|z_1|}.
\]
Under causality, $|z_1|>1$, so $|R|<1$, which implies exponential decay.

For an AR(2), the complex-root case occurs when
\[
\phi_1^2+4\phi_2<0.
\]
Equivalently,
\[
\phi_2<-\frac{\phi_1^2}{4}.
\]
The frequency term satisfies
\[
\cos(\theta)
=
\frac{\phi_1}{2\sqrt{-\phi_2}}.
\]

\begin{exercise}
Show that when $\phi_2\to 0$, the AR(2) formula recovers the AR(1)-type behavior.
\end{exercise}

\begin{proof}
When $\phi_2=0$, the AR(2) equation reduces to
\[
Y_t=\phi_1Y_{t-1}+\varepsilon_t,
\]
which is an AR(1) process. The AR polynomial becomes
\[
\Phi(z)=1-\phi_1z.
\]
Its root is
\[
z=\frac{1}{\phi_1}.
\]
Therefore,
\[
z^{-k}=\phi_1^k.
\]
The autocorrelation function is
\[
\rho_k=\phi_1^k,
\qquad k\geq 0,
\]
which is exactly the AR(1) autocorrelation function.
\end{proof}

\subsection*{Other Remarks on ACFs}

For AR($p$), if the AR polynomial has one root $z_1$ with multiplicity $r$, and all other roots are distinct, then the ACF includes terms of the form
\[
(A_1+A_2k+\cdots+A_rk^{r-1})z_1^{-k}.
\]
The remaining distinct roots contribute terms of the form
\[
A_jz_j^{-k}.
\]

More generally, the ACF of an AR($p$) process is a linear combination of exponential terms, possibly multiplied by polynomials in $k$ when roots are repeated. If roots are complex, the ACF may display damped oscillations.

For an MA($q$) process, the autocovariance function satisfies
\[
\gamma_k=0
\qquad
\text{for all } |k|>q.
\]
Thus AR processes have autocorrelations that decay, while MA processes have autocorrelations that cut off after a finite lag.

\section{Invertibility}

\begin{definition}
A general linear process $\{Y_t\}$ is invertible if there exist coefficients $\pi_j$ such that
\[
\sum_{j=0}^{\infty}|\pi_j|<\infty
\]
and
\[
\varepsilon_t
=
\sum_{j=0}^{\infty}\pi_jY_{t-j}
=
\pi_0Y_t+\pi_1Y_{t-1}+\pi_2Y_{t-2}+\cdots.
\]
\end{definition}

Invertibility means that the innovations can be recovered from the observed time series using a causal linear filter. The representation above resembles an AR($\infty$) representation.

This is analogous to causality. A causal GLP writes
\[
Y_t
=
\sum_{j=0}^{\infty}\psi_j\varepsilon_{t-j},
\]
while invertibility writes
\[
\varepsilon_t
=
\sum_{j=0}^{\infty}\pi_jY_{t-j}.
\]

\subsection*{Invertibility of MA($q$)}

Consider an MA($q$) process
\[
Y_t
=
\varepsilon_t-\theta_1\varepsilon_{t-1}
-\cdots
-\theta_q\varepsilon_{t-q}.
\]
Using the backshift operator,
\[
Y_t
=
(1-\theta_1B-\cdots-\theta_qB^q)\varepsilon_t.
\]
Define the MA polynomial
\[
\Theta(z)=1-\theta_1z-\cdots-\theta_qz^q.
\]
Then
\[
Y_t=\Theta(B)\varepsilon_t.
\]
To invert this equation, we formally write
\[
\varepsilon_t=\Theta(B)^{-1}Y_t.
\]
Therefore, invertibility requires that
\[
\Theta(B)^{-1}
=
\sum_{j=0}^{\infty}\pi_jB^j
\]
with
\[
\sum_{j=0}^{\infty}|\pi_j|<\infty.
\]

\begin{theorem}[Invertibility Condition for MA($q$)]
An MA($q$) process is invertible if and only if all roots of the MA polynomial
\[
\Theta(z)=1-\theta_1z-\cdots-\theta_qz^q
\]
lie outside the unit circle.
\end{theorem}

\begin{proof}
Assume the roots of $\Theta(z)$ are $\xi_1,\ldots,\xi_q$, with $|\xi_i|>1$ for every $i$. Then
\[
\Theta(z)
=
\left(1-\frac{z}{\xi_1}\right)
\cdots
\left(1-\frac{z}{\xi_q}\right).
\]
Thus
\[
\Theta(B)^{-1}
=
\left(1-\frac{B}{\xi_1}\right)^{-1}
\cdots
\left(1-\frac{B}{\xi_q}\right)^{-1}.
\]
Since $|1/\xi_i|<1$, each factor has the expansion
\[
\left(1-\frac{B}{\xi_i}\right)^{-1}
=
\sum_{j=0}^{\infty}\xi_i^{-j}B^j.
\]
Multiplying these convergent series gives
\[
\Theta(B)^{-1}
=
\sum_{j=0}^{\infty}\pi_jB^j
\]
with
\[
\sum_{j=0}^{\infty}|\pi_j|<\infty.
\]
Therefore,
\[
\varepsilon_t
=
\Theta(B)^{-1}Y_t
=
\sum_{j=0}^{\infty}\pi_jY_{t-j},
\]
so the MA($q$) process is invertible.
\end{proof}

\begin{exercise}
Show that the MA(1) process is invertible if and only if $|\theta|<1$.
\end{exercise}

\begin{proof}
For an MA(1) process,
\[
Y_t=\varepsilon_t-\theta\varepsilon_{t-1}.
\]
The MA polynomial is
\[
\Theta(z)=1-\theta z.
\]
Its root is
\[
z=\frac{1}{\theta}.
\]
The invertibility condition requires
\[
\left|\frac{1}{\theta}\right|>1,
\]
which is equivalent to
\[
|\theta|<1.
\]
Thus the MA(1) process is invertible if and only if $|\theta|<1$.
\end{proof}

\section{Mixed Models ARMA($p,q$)}

\begin{definition}
A time series $\{Y_t\}$ is an ARMA($p,q$) process if
\[
Y_t-\phi_1Y_{t-1}-\cdots-\phi_pY_{t-p}
=
\varepsilon_t-\theta_1\varepsilon_{t-1}-\cdots-\theta_q\varepsilon_{t-q}.
\]
\end{definition}

Using the AR and MA polynomials,
\[
\Phi(z)=1-\phi_1z-\cdots-\phi_pz^p,
\]
and
\[
\Theta(z)=1-\theta_1z-\cdots-\theta_qz^q,
\]
the model can be written compactly as
\[
\Phi(B)Y_t=\Theta(B)\varepsilon_t.
\]

Special cases include:
\[
\mathrm{ARMA}(p,0)=\mathrm{AR}(p),
\]
and
\[
\mathrm{ARMA}(0,q)=\mathrm{MA}(q).
\]

For an ARMA model:
\begin{itemize}
\item Causality is determined by the roots of $\Phi(z)$.
\item Invertibility is determined by the roots of $\Theta(z)$.
\end{itemize}
Specifically, the model is causal if all roots of $\Phi(z)$ are outside the unit circle, and it is invertible if all roots of $\Theta(z)$ are outside the unit circle.

\subsection*{A Simple Example: ARMA(1,1)}

Consider the ARMA(1,1) model
\[
Y_t-\phi Y_{t-1}
=
\varepsilon_t-\theta\varepsilon_{t-1}.
\]
Assume $\phi\neq\theta$.

If $\phi=\theta$, then
\[
(1-\phi B)Y_t=(1-\phi B)\varepsilon_t.
\]
Canceling the common factor gives
\[
Y_t=\varepsilon_t,
\]
so the model reduces to white noise.

Assume now that $|\phi|<1$, so the ARMA(1,1) process is causal. Then
\[
Y_t
=
(1-\phi B)^{-1}(1-\theta B)\varepsilon_t.
\]
Using
\[
(1-\phi B)^{-1}
=
\sum_{j=0}^{\infty}\phi^jB^j,
\]
we obtain
\[
Y_t
=
\left(\sum_{j=0}^{\infty}\phi^jB^j\right)(1-\theta B)\varepsilon_t.
\]
Expanding,
\[
Y_t
=
\left[
1+(\phi-\theta)B+\phi(\phi-\theta)B^2+\phi^2(\phi-\theta)B^3+\cdots
\right]\varepsilon_t.
\]
Therefore,
\[
Y_t
=
\varepsilon_t+\sum_{j=1}^{\infty}\phi^{j-1}(\phi-\theta)\varepsilon_{t-j}.
\]
Thus the GLP coefficients are
\[
\psi_0=1,
\qquad
\psi_j=\phi^{j-1}(\phi-\theta),
\quad j\geq 1.
\]

\begin{exercise}
Verify that
\[
\sum_{j=0}^{\infty}|\psi_j|<\infty.
\]
\end{exercise}

\begin{proof}
Using the coefficients above,
\[
\sum_{j=0}^{\infty}|\psi_j|
=
1+\sum_{j=1}^{\infty}|\phi|^{j-1}|\phi-\theta|.
\]
Since $|\phi|<1$,
\[
\sum_{j=1}^{\infty}|\phi|^{j-1}
=
\frac{1}{1-|\phi|}.
\]
Therefore,
\[
\sum_{j=0}^{\infty}|\psi_j|
=
1+\frac{|\phi-\theta|}{1-|\phi|}
<\infty.
\]
\end{proof}

If $\theta=0$, then the GLP representation becomes
\[
Y_t
=
\varepsilon_t+\sum_{j=1}^{\infty}\phi^j\varepsilon_{t-j}
=
\sum_{j=0}^{\infty}\phi^j\varepsilon_{t-j},
\]
which recovers the AR(1) representation.

\section{More on ARMA(1,1)}

\subsection*{ACVF of ARMA(1,1)}

Recall the ARMA(1,1) model
\[
Y_t-\phi Y_{t-1}
=
\varepsilon_t-\theta\varepsilon_{t-1}.
\]
When $|\phi|<1$, the causal GLP representation is
\[
Y_t
=
\varepsilon_t+\sum_{j=1}^{\infty}\phi^{j-1}(\phi-\theta)\varepsilon_{t-j}.
\]

We can derive the autocovariances using Yule--Walker style equations. For $k\geq 0$, multiply both sides of
\[
Y_t-\phi Y_{t-1}
=
\varepsilon_t-\theta\varepsilon_{t-1}
\]
by $Y_{t-k}$ and take expectations:
\[
\E[Y_tY_{t-k}]
-
\phi\E[Y_{t-1}Y_{t-k}]
=
\E[\varepsilon_tY_{t-k}]
-
\theta\E[\varepsilon_{t-1}Y_{t-k}].
\]
Using stationarity,
\[
\gamma_k-\phi\gamma_{k-1}
=
\E[\varepsilon_tY_{t-k}]
-
\theta\E[\varepsilon_{t-1}Y_{t-k}].
\]

For $k\geq 2$, causality implies
\[
\E[\varepsilon_tY_{t-k}]=0
\]
and
\[
\E[\varepsilon_{t-1}Y_{t-k}]=0.
\]
Therefore,
\[
\gamma_k-\phi\gamma_{k-1}=0,
\qquad k\geq 2.
\]
Hence
\[
\gamma_k=\phi\gamma_{k-1},
\qquad k\geq 2.
\]

Now compute the first two equations. From the GLP representation,
\[
Y_t
=
\varepsilon_t+(\phi-\theta)\varepsilon_{t-1}
+\phi(\phi-\theta)\varepsilon_{t-2}
+\cdots.
\]
Thus
\[
\E[\varepsilon_tY_t]=\sigma_\varepsilon^2.
\]
Also,
\[
Y_t=\phi Y_{t-1}+\varepsilon_t-\theta\varepsilon_{t-1}.
\]
Therefore,
\[
\E[\varepsilon_{t-1}Y_t]
=
\phi\E[\varepsilon_{t-1}Y_{t-1}]
+
\E[\varepsilon_{t-1}\varepsilon_t]
-
\theta\E[\varepsilon_{t-1}^2].
\]
Since
\[
\E[\varepsilon_{t-1}Y_{t-1}]=\sigma_\varepsilon^2,
\]
and
\[
\E[\varepsilon_{t-1}\varepsilon_t]=0,
\]
we get
\[
\E[\varepsilon_{t-1}Y_t]
=
(\phi-\theta)\sigma_\varepsilon^2.
\]

For $k=0$,
\[
\gamma_0-\phi\gamma_1
=
\E[\varepsilon_tY_t]
-
\theta\E[\varepsilon_{t-1}Y_t].
\]
Thus
\[
\gamma_0-\phi\gamma_1
=
\sigma_\varepsilon^2
-
\theta(\phi-\theta)\sigma_\varepsilon^2.
\]
That is,
\[
\gamma_0-\phi\gamma_1
=
(1-\theta\phi+\theta^2)\sigma_\varepsilon^2.
\]

For $k=1$,
\[
\gamma_1-\phi\gamma_0
=
\E[\varepsilon_tY_{t-1}]
-
\theta\E[\varepsilon_{t-1}Y_{t-1}].
\]
By causality,
\[
\E[\varepsilon_tY_{t-1}]=0,
\]
and
\[
\E[\varepsilon_{t-1}Y_{t-1}]=\sigma_\varepsilon^2.
\]
Therefore,
\[
\gamma_1-\phi\gamma_0
=
-\theta\sigma_\varepsilon^2.
\]

The first two Yule--Walker equations are
\[
\begin{cases}
\gamma_0-\phi\gamma_1=(1-\theta\phi+\theta^2)\sigma_\varepsilon^2,\\
\gamma_1-\phi\gamma_0=-\theta\sigma_\varepsilon^2.
\end{cases}
\]
From the second equation,
\[
\gamma_1=\phi\gamma_0-\theta\sigma_\varepsilon^2.
\]
Substitute this into the first equation:
\begin{align*}
\gamma_0-\phi(\phi\gamma_0-\theta\sigma_\varepsilon^2)
&=
(1-\theta\phi+\theta^2)\sigma_\varepsilon^2 \\
(1-\phi^2)\gamma_0+\phi\theta\sigma_\varepsilon^2
&=
(1-\theta\phi+\theta^2)\sigma_\varepsilon^2.
\end{align*}
Hence
\[
(1-\phi^2)\gamma_0
=
(1-2\theta\phi+\theta^2)\sigma_\varepsilon^2,
\]
so
\[
\gamma_0
=
\frac{1-2\theta\phi+\theta^2}{1-\phi^2}\sigma_\varepsilon^2.
\]

Then
\begin{align*}
\gamma_1
&=
\phi\gamma_0-\theta\sigma_\varepsilon^2 \\
&=
\phi
\left(
\frac{1-2\theta\phi+\theta^2}{1-\phi^2}
\right)\sigma_\varepsilon^2
-
\theta\sigma_\varepsilon^2 \\
&=
\frac{\phi(1-2\theta\phi+\theta^2)-\theta(1-\phi^2)}{1-\phi^2}\sigma_\varepsilon^2 \\
&=
\frac{\phi-\theta-\theta\phi^2+\phi\theta^2}{1-\phi^2}\sigma_\varepsilon^2 \\
&=
\frac{(\phi-\theta)(1-\theta\phi)}{1-\phi^2}\sigma_\varepsilon^2.
\end{align*}

For $k\geq 2$,
\[
\gamma_k=\phi\gamma_{k-1}.
\]
Therefore,
\[
\gamma_k
=
\phi^{k-1}\gamma_1
=
\phi^{k-1}
\frac{(\phi-\theta)(1-\theta\phi)}{1-\phi^2}
\sigma_\varepsilon^2,
\qquad k\geq 1.
\]
Thus $\gamma_k\to 0$ exponentially fast when $|\phi|<1$.

\subsection{Invertible Representation for ARMA(1,1)}

Invertibility depends on the MA polynomial
\[
\Theta(z)=1-\theta z.
\]
The root is $z=1/\theta$. Therefore, the ARMA(1,1) process is invertible if and only if
\[
|\theta|<1.
\]
When invertibility holds,
\[
\varepsilon_t
=
\Theta(B)^{-1}\Phi(B)Y_t
=
(1-\theta B)^{-1}(1-\phi B)Y_t.
\]
Since
\[
(1-\theta B)^{-1}
=
\sum_{j=0}^{\infty}\theta^jB^j,
\]
we get
\[
\varepsilon_t
=
\left(\sum_{j=0}^{\infty}\theta^jB^j\right)(1-\phi B)Y_t.
\]
Expanding,
\[
\varepsilon_t
=
Y_t+(\theta-\phi)Y_{t-1}
+\theta(\theta-\phi)Y_{t-2}
+\theta^2(\theta-\phi)Y_{t-3}
+\cdots.
\]
Equivalently,
\[
\varepsilon_t
=
Y_t+\sum_{j=1}^{\infty}\theta^{j-1}(\theta-\phi)Y_{t-j}.
\]

\begin{exercise}
Verify that
\[
\sum_{j=0}^{\infty}|\pi_j|<\infty
\]
for the invertible ARMA(1,1) representation.
\end{exercise}

\begin{proof}
The invertible representation has coefficients
\[
\pi_0=1,
\qquad
\pi_j=\theta^{j-1}(\theta-\phi),
\quad j\geq 1.
\]
Thus
\[
\sum_{j=0}^{\infty}|\pi_j|
=
1+\sum_{j=1}^{\infty}|\theta|^{j-1}|\theta-\phi|.
\]
If $|\theta|<1$, then
\[
\sum_{j=1}^{\infty}|\theta|^{j-1}
=
\frac{1}{1-|\theta|}.
\]
Therefore,
\[
\sum_{j=0}^{\infty}|\pi_j|
=
1+\frac{|\theta-\phi|}{1-|\theta|}
<\infty.
\]
\end{proof}

\section{General ARMA($p,q$)}

\subsection{Causal GLP Representation}

Consider the ARMA($p,q$) model
\[
Y_t-\phi_1Y_{t-1}-\cdots-\phi_pY_{t-p}
=
\varepsilon_t-\theta_1\varepsilon_{t-1}-\cdots-\theta_q\varepsilon_{t-q}.
\]
Equivalently,
\[
\Phi(B)Y_t=\Theta(B)\varepsilon_t.
\]
Assume all roots of $\Phi(z)$ lie outside the unit circle. Then the process is causal and has a GLP representation
\[
Y_t=\sum_{j=0}^{\infty}\psi_j\varepsilon_{t-j}.
\]

To find the coefficients $\{\psi_j\}$, substitute
\[
Y_t=\psi_0\varepsilon_t+\psi_1\varepsilon_{t-1}+\psi_2\varepsilon_{t-2}+\cdots
\]
into
\[
Y_t-\phi_1Y_{t-1}-\cdots-\phi_pY_{t-p}
=
\varepsilon_t-\theta_1\varepsilon_{t-1}-\cdots-\theta_q\varepsilon_{t-q}.
\]
Matching coefficients of $\varepsilon_{t-j}$ gives
\[
\psi_0=1,
\]
\[
\psi_1-\phi_1\psi_0=-\theta_1,
\]
\[
\psi_2-\phi_1\psi_1-\phi_2\psi_0=-\theta_2,
\]
and in general,
\[
\psi_j-\phi_1\psi_{j-1}-\cdots-\phi_p\psi_{j-p}=-\theta_j,
\]
where we define
\[
\phi_i=0 \quad \text{for } i>p,
\qquad
\theta_j=0 \quad \text{for } j>q,
\]
and
\[
\psi_j=0 \quad \text{for } j<0.
\]
Thus the recursion is
\[
\psi_j
=
\phi_1\psi_{j-1}
+\cdots+
\phi_p\psi_{j-p}
-\theta_j,
\qquad j\geq 1.
\]
For large $j$ such that $j>p$ and $j>q$, this becomes
\[
\psi_j
=
\phi_1\psi_{j-1}
+\cdots+
\phi_p\psi_{j-p}.
\]

\subsection{Yule--Walker Approach for the ACVF}

For the ARMA($p,q$) model,
\[
\Phi(B)Y_t=\Theta(B)\varepsilon_t.
\]
For lags $k>q$, the right-hand side is uncorrelated with $Y_{t-k}$ under causality. Therefore, multiplying the ARMA equation by $Y_{t-k}$ and taking expectations gives
\[
\E[Y_tY_{t-k}]
-\phi_1\E[Y_{t-1}Y_{t-k}]
-\cdots
-\phi_p\E[Y_{t-p}Y_{t-k}]
=0,
\qquad k>q.
\]
Using stationarity,
\[
\gamma_k-\phi_1\gamma_{k-1}-\cdots-\phi_p\gamma_{k-p}=0,
\qquad k>q.
\]
Thus, for large enough lags, the autocovariance function of an ARMA($p,q$) process satisfies the same recursion as the autocovariance function of an AR($p$) process.

If the roots of $\Phi(z)$ are distinct, then for sufficiently large $k$,
\[
\gamma_k
=
A_1z_1^{-k}+A_2z_2^{-k}+\cdots+A_pz_p^{-k}.
\]
The constants are determined by the initial autocovariances.

\begin{remark}
The MA part affects the initial autocovariances, but the eventual decay pattern of the ACF is governed by the AR roots.
\end{remark}

\section{Summary of Basic ARMA Behavior}

\begin{xltabular}{\textwidth}{l l l l l}
\toprule
Process & Causal & Invertible & Weakly stationary & ACF behavior \\
\midrule
MA(1)
& Always
& $|\theta|<1$
& Always
& $\rho_1\neq 0$, $\rho_k=0$ for $k\geq 2$ \\

MA($q$)
& Always
& Roots of $\Theta(z)$ outside unit circle
& Always
& $\rho_k=0$ for $k>q$ \\

AR(1)
& $|\phi|<1$
& Always
& $|\phi|\neq 1$
& $\rho_k=\phi^k$ for $k\geq 0$ \\

AR($p$)
& Roots of $\Phi(z)$ outside unit circle
& Always
& No roots of $\Phi(z)$ on unit circle
& Exponential decay, possibly damped oscillation \\

ARMA(1,1)
& $|\phi|<1$
& $|\theta|<1$
& $|\phi|\neq 1$
& Exponential decay after lag $1$ \\

ARMA($p,q$)
& Roots of $\Phi(z)$ outside unit circle
& Roots of $\Theta(z)$ outside unit circle
& No roots of $\Phi(z)$ on unit circle
& AR-type decay after lag $q$ \\
\bottomrule
\end{xltabular}

\section{Trends}

So far, we have considered stationary time series models such as AR($p$), MA($q$), ARMA($p,q$), and general linear processes. These models are useful for modeling the stochastic component of a time series.

For an observed time series $\{Y_t\}$, it is often useful to decompose the series as
\[
\underbrace{Y_t}_{\text{observed}}
=
\underbrace{\mu_t}_{\text{deterministic}}
+
\underbrace{X_t}_{\text{stochastic}}.
\]
Here, $\{X_t\}$ is a stationary stochastic component, and $\{\mu_t\}$ is a deterministic component. The deterministic component may reflect a trend, seasonality, or a combination of both.

The basic strategy is:

\begin{enumerate}
\item Estimate the deterministic component $\mu_t$ by $\widehat{\mu}_t$.
\item Define the residual or detrended series
\[
\widehat{X}_t=Y_t-\widehat{\mu}_t.
\]
\item Fit a stationary time series model, such as AR($p$), MA($q$), or ARMA($p,q$), to $\{\widehat{X}_t\}$.
\end{enumerate}

\section{Estimating the Sample Mean of a Time Series}

The simplest trend model is a constant mean:
\[
\mu_t=\mu.
\]
In this case,
\[
Y_t=\mu+X_t,
\]
where $\{X_t\}$ is mean-zero and stationary. A natural estimator of $\mu$ is the sample mean
\[
\overline{Y}
=
\frac{1}{n}\sum_{t=1}^nY_t.
\]

Under suitable dependence conditions,
\[
\overline{Y}\convd \Normal(\mu,\Var(\overline{Y})).
\]
Equivalently,
\[
\frac{\overline{Y}-\mu}{\sqrt{\Var(\overline{Y})}}
\convd
\Normal(0,1).
\]

\begin{remark}
This conclusion requires dependence conditions. The classical central limit theorem assumes IID observations. For time series data, serial dependence changes the variance of the sample mean.
\end{remark}

For a stationary time series with autocorrelation function $\rho_k$, we previously derived
\[
\Var(\overline{Y})
=
\frac{\gamma_0}{n}
\left[
1+2\sum_{k=1}^n
\left(1-\frac{k}{n}\right)\rho_k
\right].
\]
If $\rho_k=0$ for all $k\neq 0$, then
\[
\Var(\overline{Y})
=
\frac{\gamma_0}{n},
\]
which is the IID variance formula.

If the process is $q$-dependent, then $\rho_k=0$ for $k>q$, so
\[
\Var(\overline{Y})
=
\frac{\gamma_0}{n}
\left[
1+2\sum_{k=1}^{q}
\left(1-\frac{k}{n}\right)\rho_k
\right].
\]

If
\[
\sum_{k=1}^{\infty}|\rho_k|<\infty,
\]
then for large $n$,
\[
\Var(\overline{Y})
\approx
\frac{\gamma_0}{n}
\left[
1+2\sum_{k=1}^{\infty}\rho_k
\right].
\]
The quantity
\[
\gamma_0\left[
1+2\sum_{k=1}^{\infty}\rho_k
\right]
=
\gamma_0+2\sum_{k=1}^{\infty}\gamma_k
\]
is the long-run variance.

\subsection*{Example: MA(1)}

Consider the MA(1) process
\[
Y_t=\varepsilon_t-\frac{1}{2}\varepsilon_{t-1}.
\]
This is $1$-dependent. For an MA(1) process,
\[
\rho_1
=
\frac{-\theta}{1+\theta^2}.
\]
With $\theta=1/2$,
\[
\rho_1
=
\frac{-1/2}{1+1/4}
=
-\frac{2}{5}.
\]
Also,
\[
\rho_k=0
\qquad
\text{for } k\geq 2.
\]
Therefore,
\begin{align*}
\Var(\overline{Y})
&=
\frac{\gamma_0}{n}
\left[
1+2\left(1-\frac{1}{n}\right)\rho_1
\right] \\
&=
\frac{\gamma_0}{n}
\left[
1+2\left(1-\frac{1}{n}\right)\left(-\frac{2}{5}\right)
\right] \\
&=
\frac{\gamma_0}{n}
\left[
1-\frac{4}{5}\left(1-\frac{1}{n}\right)
\right].
\end{align*}
For large $n$,
\[
\Var(\overline{Y})
\approx
\frac{\gamma_0}{n}
\left(1-\frac{4}{5}\right)
=
\frac{\gamma_0}{5n}.
\]
Thus the variance of the sample mean is approximately five times smaller than the IID formula $\gamma_0/n$.

Since the process is $1$-dependent, the $q$-dependent CLT applies. Therefore,
\[
\overline{Y}
\approx
\Normal\left(\mu,\frac{\gamma_0}{5n}\right)
\]
for large $n$. An approximate $95\%$ confidence interval for $\mu$ is
\[
\overline{Y}\pm 2\sqrt{\frac{\gamma_0}{5n}}.
\]

\subsection*{Example: AR(1)}

Consider a causal AR(1) process,
\[
Y_t-\phi Y_{t-1}=\varepsilon_t,
\qquad
|\phi|<1.
\]
For a causal AR(1),
\[
\rho_k=\phi^k,
\qquad k\geq 0.
\]
Therefore,
\[
\Var(\overline{Y})
=
\frac{\gamma_0}{n}
\left[
1+2\sum_{k=1}^{n}
\left(1-\frac{k}{n}\right)\phi^k
\right].
\]
For large $n$,
\[
\Var(\overline{Y})
\approx
\frac{\gamma_0}{n}
\left[
1+2\sum_{k=1}^{\infty}\phi^k
\right].
\]
Since
\[
\sum_{k=1}^{\infty}\phi^k
=
\frac{\phi}{1-\phi},
\]
we get
\[
\Var(\overline{Y})
\approx
\frac{\gamma_0}{n}
\left[
1+\frac{2\phi}{1-\phi}
\right]
=
\frac{\gamma_0}{n}
\left(
\frac{1+\phi}{1-\phi}
\right).
\]

\begin{exercise}
Verify the final simplification above.
\end{exercise}

\begin{proof}
Starting with
\[
1+\frac{2\phi}{1-\phi},
\]
put the terms over the common denominator $1-\phi$:
\[
1+\frac{2\phi}{1-\phi}
=
\frac{1-\phi}{1-\phi}
+
\frac{2\phi}{1-\phi}
=
\frac{1+\phi}{1-\phi}.
\]
\end{proof}

For example, if $\phi=0.9$, then
\[
\frac{1+\phi}{1-\phi}
=
\frac{1.9}{0.1}
=
19.
\]
Thus, for large $n$,
\[
\Var(\overline{Y})
\approx
\frac{19\gamma_0}{n},
\]
which is approximately $19$ times larger than the IID variance formula.

\begin{remark}
This example shows why positively correlated time series can have much less effective information than IID data. Positive serial dependence inflates the variance of the sample mean. Negative serial dependence can reduce it.
\end{remark}

\subsection*{Example: Random Walk}

Consider the random walk
\[
Y_1=\varepsilon_1,
\qquad
Y_t=Y_{t-1}+\varepsilon_t,
\qquad
\varepsilon_t \iid \Normal(0,\sigma^2).
\]
Equivalently,
\[
Y_t=\sum_{j=1}^t\varepsilon_j.
\]
The random walk is not stationary, so the stationary sample mean formulas above do not apply. Nevertheless, we can calculate the variance directly.

The sample mean is
\[
\overline{Y}
=
\frac{1}{n}\sum_{t=1}^nY_t
=
\frac{1}{n}\sum_{t=1}^n\sum_{j=1}^t\varepsilon_j.
\]
Reverse the order of summation. For a fixed innovation $\varepsilon_j$, it appears in $Y_j,Y_{j+1},\ldots,Y_n$, so it appears $n-j+1$ times. Therefore,
\[
\overline{Y}
=
\frac{1}{n}\sum_{j=1}^n(n-j+1)\varepsilon_j.
\]
Hence
\begin{align*}
\Var(\overline{Y})
&=
\frac{1}{n^2}
\Var\left(
\sum_{j=1}^n(n-j+1)\varepsilon_j
\right) \\
&=
\frac{1}{n^2}
\sum_{j=1}^n(n-j+1)^2\Var(\varepsilon_j) \\
&=
\frac{\sigma^2}{n^2}
\sum_{j=1}^n(n-j+1)^2.
\end{align*}
Since
\[
\sum_{j=1}^n(n-j+1)^2
=
\sum_{m=1}^n m^2
=
\frac{n(n+1)(2n+1)}{6},
\]
we get
\[
\Var(\overline{Y})
=
\frac{\sigma^2}{n^2}
\frac{n(n+1)(2n+1)}{6}.
\]
As $n\to\infty$,
\[
\Var(\overline{Y})
\sim
\frac{\sigma^2}{n^2}\cdot \frac{2n^3}{6}
=
\frac{\sigma^2 n}{3}.
\]
Thus the variance of the sample mean grows with $n$. As the sample size increases, we become less certain about the mean level of a random walk.

\section{Different Models for Trend and Stochastic Components}

There are several common ways to combine deterministic and stochastic components.

\begin{itemize}
\item Additive trend model:
\[
Y_t=\mu_t+X_t.
\]

\item Additive trend and seasonality model:
\[
Y_t=T_t+S_t+X_t.
\]

\item Multiplicative model:
\[
Y_t=\mu_tX_t,
\qquad
\text{or}
\qquad
Y_t=T_tS_tX_t.
\]
Taking logarithms transforms this into an additive model:
\[
\log Y_t=\log T_t+\log S_t+\log X_t.
\]

\item Mixture of additive and multiplicative models:
\[
Y_t=T_tS_t+X_t,
\qquad
\text{or}
\qquad
Y_t=(T_t+S_t)X_t.
\]
\end{itemize}

\section{Regression Methods for Trends}

Suppose
\[
Y_t=\mu_t+X_t,
\]
where $\mu_t$ is deterministic and $\{X_t\}$ is a stationary stochastic component. The regression approach estimates $\mu_t$ using the observed time series, then defines
\[
\widehat{X}_t=Y_t-\widehat{\mu}_t.
\]
The residual series $\{\widehat{X}_t\}$ is then treated as an estimate of the stationary component.

\subsection{Linear Trend}

Assume the deterministic trend is linear:
\[
\mu_t=\beta_0+\beta_1t.
\]
Given observations $Y_1,\ldots,Y_n$, estimate $\beta_0,\beta_1$ by minimizing
\[
Q(\beta_0,\beta_1)
=
\sum_{t=1}^n(Y_t-\beta_0-\beta_1t)^2.
\]
The normal equations are obtained by differentiating:
\[
\frac{\partial Q}{\partial \beta_0}
=
-2\sum_{t=1}^n(Y_t-\beta_0-\beta_1t)=0,
\]
and
\[
\frac{\partial Q}{\partial \beta_1}
=
-2\sum_{t=1}^n t(Y_t-\beta_0-\beta_1t)=0.
\]
The resulting least squares estimators are
\[
\widehat{\beta}_1
=
\frac{\sum_{t=1}^n(t-\overline{t})(Y_t-\overline{Y})}
{\sum_{t=1}^n(t-\overline{t})^2},
\qquad
\widehat{\beta}_0
=
\overline{Y}-\widehat{\beta}_1\overline{t},
\]
where
\[
\overline{t}=\frac{1}{n}\sum_{t=1}^nt,
\qquad
\overline{Y}=\frac{1}{n}\sum_{t=1}^nY_t.
\]
The estimated trend is
\[
\widehat{\mu}_t=\widehat{\beta}_0+\widehat{\beta}_1t.
\]

\subsection{Other Regression Trend Models}

A quadratic trend can be written as
\[
\mu_t=\beta_0+\beta_1t+\beta_2t^2.
\]
This is still a linear regression model because it is linear in the unknown coefficients:
\[
\mu_t=
\begin{bmatrix}
1 & t & t^2
\end{bmatrix}
\begin{bmatrix}
\beta_0\\
\beta_1\\
\beta_2
\end{bmatrix}.
\]

A cosine trend can be written as
\[
\mu_t
=
\beta_0
+
\beta_1\cos\left(\frac{2\pi t}{f}\right)
+
\beta_2\sin\left(\frac{2\pi t}{f}\right),
\]
where $f$ is the frequency.

For seasonal or cyclical means, suppose $t$ is measured monthly and the mean repeats every year. Then
\[
\mu_1=\mu_{13}=\mu_{25}=\cdots
\]
for January,
\[
\mu_2=\mu_{14}=\mu_{26}=\cdots
\]
for February, and so on. This can be estimated using a dummy-variable regression:
\[
Y_t
=
\beta_1X_{\text{Jan},t}
+
\beta_2X_{\text{Feb},t}
+
\cdots
+
\beta_{12}X_{\text{Dec},t}
+
\varepsilon_t,
\]
where
\[
X_{\text{Jan},t}
=
\begin{cases}
1, & t \text{ is January},\\
0, & \text{otherwise}.
\end{cases}
\]
After fitting the regression, the estimated seasonal means are
\[
\widehat{\mu}_1,\ldots,\widehat{\mu}_{12}
=
\widehat{\beta}_1,\ldots,\widehat{\beta}_{12}.
\]

\section{More on Regression for Time Series}

Recall that regression methods are used to estimate $\mu_t$ in the additive model
\[
Y_t=\mu_t+X_t.
\]
After estimating $\mu_t$ by $\widehat{\mu}_t$, we work with
\[
\widehat{X}_t=Y_t-\widehat{\mu}_t.
\]

In ordinary regression, the error vector is often modeled as
\[
\varepsilon\sim\Normal(0,\sigma^2I).
\]
Then
\[
\widehat{\beta}
=
(X^\top X)^{-1}X^\top Y
\]
has variance
\[
\Var(\widehat{\beta})
=
\sigma^2(X^\top X)^{-1}.
\]
For time series regression, the stochastic component $\{X_t\}$ is usually not IID. In matrix form, we may instead have
\[
\varepsilon\sim\Normal(0,V),
\]
where $V$ is not diagonal. Then
\[
\Var(\widehat{\beta})
=
(X^\top X)^{-1}X^\top VX(X^\top X)^{-1}.
\]
Thus the usual regression standard errors are generally not reliable unless serial dependence is accounted for.

\begin{theorem}
If the deterministic trend is polynomial, trigonometric, seasonal, or a finite linear combination of these basis functions, then the ordinary least squares estimator of the trend coefficients is unbiased for large samples under suitable stationarity assumptions on the stochastic component.
\end{theorem}

\begin{proof}
Write the model as
\[
Y=X\beta+\varepsilon,
\]
where $X$ is the deterministic design matrix and $\varepsilon$ is the stochastic time series component with
\[
\E[\varepsilon]=0.
\]
The ordinary least squares estimator is
\[
\widehat{\beta}=(X^\top X)^{-1}X^\top Y.
\]
Substituting $Y=X\beta+\varepsilon$ gives
\[
\widehat{\beta}
=
(X^\top X)^{-1}X^\top(X\beta+\varepsilon)
=
\beta+(X^\top X)^{-1}X^\top\varepsilon.
\]
Taking expectations,
\[
\E[\widehat{\beta}]
=
\beta+(X^\top X)^{-1}X^\top\E[\varepsilon]
=
\beta.
\]
Thus the estimator is unbiased whenever the stochastic component has mean zero and the design matrix has full column rank.
\end{proof}

\section{Sample Autocorrelation Function}

After detrending a series, we use the residuals
\[
\widehat{X}_t=Y_t-\widehat{\mu}_t
\]
to study the dependence structure of the stochastic component.

For a stationary time series, the theoretical autocorrelation function is
\[
\rho_k=\frac{\gamma_k}{\gamma_0}.
\]
The sample autocorrelation function estimates $\rho_k$ from data.

\begin{definition}
Given observations $Y_1,\ldots,Y_n$, the sample autocorrelation at lag $k$ is
\[
r_k
=
\frac{
\sum_{t=k+1}^n(Y_t-\overline{Y})(Y_{t-k}-\overline{Y})
}{
\sum_{t=1}^n(Y_t-\overline{Y})^2
}.
\]
\end{definition}

Equivalently,
\[
r_k
=
\frac{
\frac{1}{n}\sum_{t=k+1}^n(Y_t-\overline{Y})(Y_{t-k}-\overline{Y})
}{
\frac{1}{n}\sum_{t=1}^n(Y_t-\overline{Y})^2
}.
\]
For large $n$ and fixed $k$, this is approximately
\[
r_k
\approx
\frac{\widehat{\Cov}(Y_t,Y_{t-k})}{\widehat{\Var}(Y_t)}.
\]

\begin{remark}
The construction of the sample ACF produces a sample autocorrelation matrix that is nonnegative definite. This property is important for valid covariance modeling.
\end{remark}

\begin{theorem}
If $Y_t\iid \Normal(0,\sigma^2)$, then for large $n$ and fixed $k$,
\[
r_k\approx \Normal\left(0,\frac{1}{n}\right).
\]
\end{theorem}

This result motivates the common approximate white-noise bounds
\[
\pm \frac{2}{\sqrt{n}}.
\]
Sample autocorrelations outside these bounds provide evidence against IID white noise.

\section{A Model for Nonstationary Time Series: ARIMA($p,d,q$)}

\subsection{Differencing Operators}

\begin{definition}
The differencing operator $\nabla$ is defined by
\[
\nabla Y_t
=
(1-B)Y_t
=
Y_t-Y_{t-1}.
\]
\end{definition}

The second difference is
\begin{align*}
\nabla^2Y_t
&=
\nabla(Y_t-Y_{t-1}) \\
&=
(Y_t-Y_{t-1})-(Y_{t-1}-Y_{t-2}) \\
&=
Y_t-2Y_{t-1}+Y_{t-2}.
\end{align*}
Equivalently,
\[
\nabla^2Y_t=(1-B)^2Y_t.
\]
More generally,
\[
\nabla^dY_t=(1-B)^dY_t.
\]

For seasonal models, the lag-$s$ differencing operator is
\[
\nabla_sY_t
=
Y_t-Y_{t-s}
=
(1-B^s)Y_t.
\]
A combined seasonal and ordinary difference can be written as
\[
\nabla_s\nabla^dY_t
=
(1-B^s)(1-B)^dY_t.
\]

\subsection*{Examples of Differencing}

\begin{example}
Consider a linear trend plus stationary noise:
\[
Y_t=\beta_0+\beta_1t+X_t,
\]
where $\{X_t\}$ is stationary. Then
\begin{align*}
\nabla Y_t
&=
Y_t-Y_{t-1} \\
&=
(\beta_0+\beta_1t+X_t)
-
(\beta_0+\beta_1(t-1)+X_{t-1}) \\
&=
\beta_1+X_t-X_{t-1}.
\end{align*}
The differenced process is stationary because it is a constant plus a linear combination of stationary variables.
\end{example}

\begin{exercise}
Why is $Y_t$ not stationary in the previous example?
\end{exercise}

\begin{proof}
The mean function is
\[
\E[Y_t]
=
\beta_0+\beta_1t+\E[X_t].
\]
If $\{X_t\}$ is stationary, then $\E[X_t]=\mu_X$ is constant. Hence
\[
\E[Y_t]
=
\beta_0+\mu_X+\beta_1t.
\]
When $\beta_1\neq 0$, the mean depends on $t$, so $Y_t$ is not weakly stationary.
\end{proof}

\begin{example}
Consider the random walk
\[
Y_1=\varepsilon_1,
\qquad
Y_t=Y_{t-1}+\varepsilon_t.
\]
Then
\[
\nabla Y_t
=
Y_t-Y_{t-1}
=
\varepsilon_t.
\]
Thus the first difference of a random walk is white noise, which is stationary.
\end{example}

\subsection{ARIMA($p,d,q$)}

In the random walk example,
\[
\nabla Y_t=\varepsilon_t.
\]
Thus the differenced series is an ARMA$(0,0)$ process.

More generally, if
\[
W_t=\nabla^dY_t
\]
is an ARMA($p,q$) process, then $\{Y_t\}$ is called an ARIMA($p,d,q$) process.

\begin{definition}
A time series $\{Y_t\}$ is ARIMA($p,d,q$) if
\[
\nabla^dY_t
\]
is an ARMA($p,q$) process.
\end{definition}

Equivalently,
\[
\Phi(B)(1-B)^dY_t=\Theta(B)\varepsilon_t.
\]

\subsection*{Random Walk Plus Noise}

Consider
\[
Y_t=X_t+\eta_t,
\]
where
\[
X_t=\sum_{j=1}^t\varepsilon_j
\]
is a random walk, $\{\eta_t\}$ is an independent white noise process, and
\[
\varepsilon_t\iid \Normal(0,\sigma_\varepsilon^2),
\qquad
\eta_t\iid \Normal(0,\sigma_\eta^2).
\]
Then
\[
\Var(Y_t)=\Var(X_t)+\Var(\eta_t)
=
t\sigma_\varepsilon^2+\sigma_\eta^2,
\]
so the process is not stationary.

Now take first differences:
\begin{align*}
W_t
&=
\nabla Y_t \\
&=
(X_t+\eta_t)-(X_{t-1}+\eta_{t-1}) \\
&=
(X_t-X_{t-1})+\eta_t-\eta_{t-1} \\
&=
\varepsilon_t+\eta_t-\eta_{t-1}.
\end{align*}
The autocovariance function of $W_t$ satisfies
\[
\gamma_k=0
\qquad
\text{for } |k|\geq 2.
\]
Thus $\{W_t\}$ has the same second-order structure as an MA(1) process. Therefore,
\[
W_t\sim \mathrm{MA}(1),
\]
and
\[
Y_t\sim \mathrm{ARIMA}(0,1,1).
\]

\begin{remark}
The notation $W_t\sim \mathrm{MA}(1)$ here means that $W_t$ has the autocovariance structure of an MA(1) process. Under Gaussian assumptions, this second-order matching corresponds to equality in distribution for the process.
\end{remark}

\section{More on ARIMA($p,d,q$)}

Last time, we introduced ARIMA$(p,d,q)$ as a model for nonstationary time series. If
\[
Y_t\sim \mathrm{ARIMA}(p,d,q),
\]
then taking differences $d$ times gives a stationary time series
\[
W_t=\nabla^dY_t
\]
such that
\[
W_t\sim \mathrm{ARMA}(p,q).
\]

\subsection*{Review of the Random Walk Plus Noise Example}

Consider the random walk plus noise model
\[
Y_t=X_t+\eta_t
=
\sum_{j=1}^t\varepsilon_j+\eta_t,
\]
where
\[
\eta_t\iid \Normal(0,\sigma_\eta^2),
\qquad
\varepsilon_t\iid \Normal(0,\sigma_\varepsilon^2),
\]
and $\{\eta_t\}$ is independent of $\{\varepsilon_t\}$.

Taking first differences gives
\[
W_t
=
\nabla Y_t
=
Y_t-Y_{t-1}
=
\varepsilon_t+\eta_t-\eta_{t-1}.
\]
Now compute the moments of $W_t$. First,
\[
\E[W_t]
=
\E[\varepsilon_t+\eta_t-\eta_{t-1}]
=
0.
\]
Next,
\[
\Var(W_t)
=
\Var(\varepsilon_t+\eta_t-\eta_{t-1})
=
\sigma_\varepsilon^2+2\sigma_\eta^2,
\]
by independence.

For lag $1$,
\begin{align*}
\gamma_1
&=
\Cov(W_t,W_{t-1}) \\
&=
\Cov(\varepsilon_t+\eta_t-\eta_{t-1},
\varepsilon_{t-1}+\eta_{t-1}-\eta_{t-2}) \\
&=
-\Var(\eta_{t-1}) \\
&=
-\sigma_\eta^2.
\end{align*}
For $|k|\geq 2$, there are no overlapping innovation terms, so
\[
\gamma_k=0,
\qquad |k|\geq 2.
\]
Thus $\{W_t\}$ is stationary and has the autocovariance structure of an MA(1) process.

Therefore, there exists an uncorrelated innovation sequence $\{\widetilde{\varepsilon}_t\}$ and a constant $\theta$ such that
\[
W_t
=
\widetilde{\varepsilon}_t-\theta\widetilde{\varepsilon}_{t-1}
\sim \mathrm{MA}(1).
\]
Hence
\[
Y_t\sim \mathrm{ARIMA}(0,1,1).
\]

\subsection{AR and MA Polynomials for ARIMA($p,d,q$)}

If
\[
Y_t\sim \mathrm{ARIMA}(p,d,q),
\]
then
\[
W_t=\nabla^dY_t
\]
satisfies
\[
W_t\sim \mathrm{ARMA}(p,q).
\]
Thus
\[
\Phi(B)W_t=\Theta(B)\varepsilon_t.
\]
Since
\[
W_t=(1-B)^dY_t,
\]
we have
\[
\Phi(B)(1-B)^dY_t=\Theta(B)\varepsilon_t.
\]
Therefore, the AR polynomial of the nonstationary ARIMA model is
\[
\Phi^\ast(z)
=
\Phi(z)(1-z)^d.
\]
This polynomial has the root $z=1$ repeated $d$ times. The remaining roots are the roots of $\Phi(z)$.

If $W_t$ is causal, then the roots of $\Phi(z)$ are outside the unit circle. The roots at $z=1$ are exactly what create nonstationarity in the original process $\{Y_t\}$.

\subsection{Overdifferencing}

Usually, $d=1$ or $d=2$. If $d$ is too large, the series is overdifferenced. Overdifferencing can:

\begin{itemize}
\item lead to models that are more complicated than necessary,
\item introduce artificial negative autocorrelation,
\item lead to noninvertible moving average components.
\end{itemize}

For example, consider the random walk
\[
Y_t=\sum_{i=1}^t\varepsilon_i.
\]
Then
\[
\nabla Y_t
=
Y_t-Y_{t-1}
=
\varepsilon_t.
\]
Thus the first difference is white noise and can be modeled as ARMA$(0,0)$.

If we difference one more time, then
\[
\nabla^2Y_t
=
\nabla\varepsilon_t
=
\varepsilon_t-\varepsilon_{t-1}.
\]
This is an MA(1) process with parameter $\theta=1$ under the convention
\[
W_t=\varepsilon_t-\theta\varepsilon_{t-1}.
\]
Since the MA polynomial is
\[
\Theta(z)=1-z,
\]
the root is $z=1$, which lies on the unit circle. Therefore, the overdifferenced process is noninvertible.

\section*{ARFIMA($p,d,q$)}

ARFIMA stands for autoregressive fractionally integrated moving average. The letter FI means fractionally integrated.

For a real number
\[
d\in(0,0.5),
\]
define the fractional differencing operator by
\[
\nabla^d
=
(1-B)^d
=
\sum_{i=0}^{\infty}b_iB^i,
\]
where the coefficients $\{b_i\}$ are defined by the binomial expansion:
\[
b_i
=
(-1)^i\binom{d}{i}
=
(-1)^i
\frac{d(d-1)\cdots(d-i+1)}{i!}.
\]
Thus
\[
(1-B)^d
=
1-dB+\frac{d(d-1)}{2}B^2-\frac{d(d-1)(d-2)}{6}B^3+\cdots.
\]

Fractional differencing is used to model long-range dependence. Unlike ARMA processes, whose autocorrelations usually decay exponentially, fractionally integrated processes can have autocorrelations that decay polynomially.

\section{GLP-Like Representation of ARIMA($p,d,q$)}

If
\[
Y_t\sim \mathrm{ARIMA}(p,d,q),
\]
then
\[
Y_t
\]
is nonstationary, but
\[
W_t=\nabla^dY_t
\]
is stationary ARMA$(p,q)$. Hence
\[
\Phi(B)W_t=\Theta(B)\varepsilon_t.
\]
Substituting $W_t=(1-B)^dY_t$ gives
\[
\Phi(B)(1-B)^dY_t=\Theta(B)\varepsilon_t.
\]

For a nonstationary process, we cannot obtain a true stationary GLP representation. However, it is still useful to derive GLP-like expansions to understand the autocovariance behavior and the effect of integration.

\subsection{The Corresponding Nonstationary ARMA$(p+d,q)$ Model}

Suppose $d=1$. Then
\[
W_t=Y_t-Y_{t-1}.
\]
If
\[
W_t\sim \mathrm{ARMA}(p,q),
\]
then
\[
W_t-\phi_1W_{t-1}-\cdots-\phi_pW_{t-p}
=
\varepsilon_t-\theta_1\varepsilon_{t-1}-\cdots-\theta_q\varepsilon_{t-q}.
\]
Substitute $W_t=Y_t-Y_{t-1}$:
\begin{align*}
&(Y_t-Y_{t-1})
-\phi_1(Y_{t-1}-Y_{t-2})
-\cdots
-\phi_p(Y_{t-p}-Y_{t-p-1}) \\
&=
\varepsilon_t-\theta_1\varepsilon_{t-1}-\cdots-\theta_q\varepsilon_{t-q}.
\end{align*}
Collecting terms,
\[
Y_t-(1+\phi_1)Y_{t-1}
-(\phi_2-\phi_1)Y_{t-2}
-\cdots
-(\phi_p-\phi_{p-1})Y_{t-p}
+\phi_pY_{t-p-1}
=
\varepsilon_t-\theta_1\varepsilon_{t-1}-\cdots-\theta_q\varepsilon_{t-q}.
\]
This resembles an ARMA$(p+1,q)$ equation, but it is nonstationary because the AR polynomial contains the unit-root factor $(1-z)$:
\[
\Phi^\ast(z)
=
(1-z)\Phi(z).
\]

\subsection*{GLP-Like Representation of ARIMA$(0,1,1)$}

Suppose
\[
Y_t\sim \mathrm{ARIMA}(0,1,1).
\]
Then
\[
Y_t-Y_{t-1}
=
\varepsilon_t-\theta\varepsilon_{t-1}.
\]
Thus
\[
Y_t
=
Y_{t-1}
+
\varepsilon_t
-
\theta\varepsilon_{t-1}.
\]
Iterating backward,
\begin{align*}
Y_t
&=
Y_{t-1}+\varepsilon_t-\theta\varepsilon_{t-1} \\
&=
Y_{t-2}+\varepsilon_{t-1}-\theta\varepsilon_{t-2}
+\varepsilon_t-\theta\varepsilon_{t-1} \\
&=
Y_{t-2}
+\varepsilon_t
+
(1-\theta)\varepsilon_{t-1}
-\theta\varepsilon_{t-2}.
\end{align*}
Continuing,
\[
Y_t
=
Y_{t-m}
+
\varepsilon_t
+
(1-\theta)\sum_{j=1}^{m-1}\varepsilon_{t-j}
-\theta\varepsilon_{t-m}.
\]
If the process is initialized at $0$ and $m$ is large, this motivates the informal representation
\[
Y_t
\approx
\sum_{j=1}^{\infty}(1-\theta)\varepsilon_{t-j}
+
\varepsilon_t.
\]
This is not a valid GLP because
\[
\sum_{j=1}^{\infty}|1-\theta|
\]
diverges unless $\theta=1$.

One can show that
\[
\Var(Y_t)
=
\left[
1+\theta^2+(1-\theta)^2(t-1)
\right]\sigma^2,
\]
which grows linearly in $t$ unless $\theta=1$.

For moderate $k$ and large $t$, the autocorrelation satisfies approximately
\[
\rho_{t,t-k}
\approx 1.
\]
Thus an ARIMA$(0,1,1)$ process has random-walk-like persistence.

\subsection*{GLP-Like Representation of ARIMA$(1,1,0)$}

Suppose
\[
Y_t\sim \mathrm{ARIMA}(1,1,0),
\]
so
\[
Y_t-Y_{t-1}
=
\phi(Y_{t-1}-Y_{t-2})+\varepsilon_t.
\]
Then
\[
Y_t
=
(1+\phi)Y_{t-1}
-
\phi Y_{t-2}
+
\varepsilon_t.
\]
This is a nonstationary AR(2) equation.

Now suppose there is a GLP-like representation
\[
Y_t
=
\psi_0\varepsilon_t+\psi_1\varepsilon_{t-1}+\psi_2\varepsilon_{t-2}+\cdots.
\]
Substitute into the nonstationary AR(2) equation:
\begin{align*}
\psi_0\varepsilon_t+\psi_1\varepsilon_{t-1}+\psi_2\varepsilon_{t-2}+\cdots
&=
(1+\phi)
(\psi_0\varepsilon_{t-1}+\psi_1\varepsilon_{t-2}+\psi_2\varepsilon_{t-3}+\cdots) \\
&\quad
-\phi
(\psi_0\varepsilon_{t-2}+\psi_1\varepsilon_{t-3}+\psi_2\varepsilon_{t-4}+\cdots)
+\varepsilon_t.
\end{align*}
Matching coefficients gives
\[
\psi_0=1,
\]
\[
\psi_1=1+\phi,
\]
and for $k\geq 2$,
\[
\psi_k=(1+\phi)\psi_{k-1}-\phi\psi_{k-2}.
\]
This recursion has solution
\[
\psi_k
=
1+\phi+\cdots+\phi^k
=
\frac{1-\phi^{k+1}}{1-\phi},
\qquad k\geq 0,
\]
when $\phi\neq 1$.

\begin{exercise}
Verify the formula
\[
\psi_k
=
1+\phi+\cdots+\phi^k.
\]
\end{exercise}

\begin{proof}
We use induction. For $k=0$,
\[
\psi_0=1,
\]
which agrees with the formula. For $k=1$,
\[
\psi_1=1+\phi,
\]
which also agrees.

Assume for some $k\geq 2$ that
\[
\psi_{k-1}=1+\phi+\cdots+\phi^{k-1}
\]
and
\[
\psi_{k-2}=1+\phi+\cdots+\phi^{k-2}.
\]
Then
\begin{align*}
\psi_k
&=
(1+\phi)\psi_{k-1}-\phi\psi_{k-2} \\
&=
(1+\phi)(1+\phi+\cdots+\phi^{k-1})
-
\phi(1+\phi+\cdots+\phi^{k-2}) \\
&=
(1+\phi+\cdots+\phi^{k-1})
+
(\phi+\phi^2+\cdots+\phi^k)
-
(\phi+\phi^2+\cdots+\phi^{k-1}) \\
&=
1+\phi+\cdots+\phi^k.
\end{align*}
Thus the formula holds for all $k\geq 0$.
\end{proof}

As expected, this is not a valid stationary GLP representation because
\[
\sum_{k=0}^{\infty}|\psi_k|
\]
does not converge.

\section{Transformations of Time Series}

Suppose a time series $\{Y_t\}$ satisfies
\[
\E[Y_t]=\mu_t
\]
and
\[
\Var(Y_t)\approx \mu_t^2\sigma^2.
\]
In this case, the variance grows approximately proportionally to the square of the mean. A common variance-stabilizing transformation is the logarithm:
\[
\widetilde{Y}_t=\log Y_t.
\]

Using a first-order Taylor expansion around $\mu_t$,
\[
\log Y_t
\approx
\log \mu_t+\frac{1}{\mu_t}(Y_t-\mu_t).
\]
Therefore,
\[
\Var(\log Y_t)
\approx
\Var\left(\frac{1}{\mu_t}(Y_t-\mu_t)\right)
=
\frac{1}{\mu_t^2}\Var(Y_t).
\]
If
\[
\Var(Y_t)\approx \mu_t^2\sigma^2,
\]
then
\[
\Var(\log Y_t)\approx \sigma^2,
\]
which is approximately constant.

Thus, under these conditions, the logarithm is a variance-stabilizing transformation.

\subsection*{More Transformations}

The goal of transformation is to make a time series more stationary or closer to a normal process. If possible, a transformation should stabilize the variance and reduce skewness.

One common family is the Box--Cox transformation:
\[
g(y)
=
\begin{cases}
\dfrac{y^\lambda-1}{\lambda}, & \lambda\neq 0,\\
\log y, & \lambda=0.
\end{cases}
\]

\begin{exercise}
Show that
\[
\lim_{\lambda\to 0}\frac{y^\lambda-1}{\lambda}
=
\log y.
\]
\end{exercise}

\begin{proof}
Write
\[
y^\lambda=e^{\lambda\log y}.
\]
Then
\[
\frac{y^\lambda-1}{\lambda}
=
\frac{e^{\lambda\log y}-1}{\lambda}.
\]
Using the expansion
\[
e^{\lambda\log y}
=
1+\lambda\log y+o(\lambda),
\]
we get
\[
\frac{e^{\lambda\log y}-1}{\lambda}
=
\frac{\lambda\log y+o(\lambda)}{\lambda}
\to
\log y.
\]
\end{proof}

Another common transformation is the log-difference:
\[
\nabla \log Y_t
=
\log Y_t-\log Y_{t-1}.
\]
Suppose
\[
Y_t=Y_{t-1}(1+X_t).
\]
Then
\[
\log Y_t
=
\log Y_{t-1}+\log(1+X_t),
\]
so
\[
\nabla \log Y_t
=
\log(1+X_t).
\]
When $X_t$ is small,
\[
\log(1+X_t)\approx X_t.
\]
Therefore, log-differences are often interpreted as approximate percentage changes.

\begin{remark}
If $X_t>0$, then $\log(1+X_t)<X_t$. If $X_t<0$, then $\log(1+X_t)>X_t$ as long as $1+X_t>0$.
\end{remark}

\section*{Roadmap for Time Series Modeling}

Suppose we observe
\[
Y_1,\ldots,Y_n.
\]
A typical time series modeling workflow is:

\begin{enumerate}
\item Model specification: determine an appropriate model class and order, such as ARIMA$(p,d,q)$, using the sample ACF, sample PACF, transformations, differencing, and domain knowledge.
\item Parameter estimation: fit the selected model using methods such as maximum likelihood, least squares, or method of moments.
\item Diagnostics: analyze residuals to assess whether the fitted model captures the dependence structure.
\item Forecasting: use the fitted model to predict future values and quantify uncertainty.
\end{enumerate}

\section{Bartlett's Theorem}

Let
\[
Y_1,\ldots,Y_n
\]
be an observed time series. The sample autocorrelation at lag $k$ is
\[
\widehat{\rho}_k
=
r_k
=
\frac{
\sum_{t=k+1}^n(Y_t-\overline{Y})(Y_{t-k}-\overline{Y})
}{
\sum_{t=1}^n(Y_t-\overline{Y})^2
}.
\]

The theoretical autocorrelation $\rho_k$ is a fixed property of the data-generating process. The sample autocorrelation $r_k$ is random because it is computed from one realized sample path.

\begin{theorem}[Bartlett's Theorem]
For fixed $m$, the joint sampling distribution of
\[
(r_1,\ldots,r_m)
\]
is approximately multivariate normal:
\[
\bm{r}
\approx
\Normal\left(\bm{\rho},\frac{1}{n}C\right),
\]
where
\[
\bm{r}
=
\begin{bmatrix}
r_1\\
\vdots\\
r_m
\end{bmatrix},
\qquad
\bm{\rho}
=
\begin{bmatrix}
\rho_1\\
\vdots\\
\rho_m
\end{bmatrix},
\]
and $C$ is an $m\times m$ covariance matrix depending on the true autocorrelation function.
\end{theorem}

The diagonal entries of $C$ are
\[
c_{ii}
=
\sum_{k=-\infty}^{\infty}
\left(
\rho_k^2+\rho_{k+i}\rho_{k-i}
-4\rho_i\rho_k\rho_{k-i}
+2\rho_i^2\rho_k^2
\right).
\]
Thus, for large $n$,
\[
r_i\approx \Normal\left(\rho_i,\frac{c_{ii}}{n}\right).
\]

\subsection*{Example: White Noise}

Suppose
\[
Y_t\sim WN(0,\sigma^2).
\]
Then
\[
\rho_0=1,
\qquad
\rho_k=0
\quad
\text{for } |k|\geq 1.
\]
For $i\geq 1$,
\[
c_{ii}
=
\sum_{k=-\infty}^{\infty}
\left(
\rho_k^2+\rho_{k+i}\rho_{k-i}
\right),
\]
because $\rho_i=0$. The only nonzero contribution to
\[
\sum_{k=-\infty}^{\infty}\rho_k^2
\]
comes from $k=0$, so it equals $1$. Also,
\[
\sum_{k=-\infty}^{\infty}\rho_{k+i}\rho_{k-i}=0
\]
for $i\geq 1$. Therefore,
\[
c_{ii}=1.
\]
Hence
\[
r_i\approx \Normal\left(0,\frac{1}{n}\right),
\qquad i\geq 1.
\]
This gives the approximate white-noise bounds
\[
-\frac{2}{\sqrt{n}}
\leq
r_i
\leq
\frac{2}{\sqrt{n}}.
\]

\subsection{Example: AR(1)}

Suppose $\{Y_t\}$ follows a causal AR(1) process. Then
\[
\rho_k=\phi^{|k|}.
\]
Using Bartlett's theorem, one can show that for large $n$,
\[
\Var(r_i)
\approx
\frac{1}{n}
\left[
\frac{1+\phi^2}{1-\phi^2}
-
2i\phi^{2i}
\right].
\]
When $i=1$,
\[
\Var(r_1)
\approx
\frac{1-\phi^2}{n}.
\]
Thus, if $|\phi|$ is close to $1$, then $r_1$ is a relatively precise estimator of $\rho_1$.

For large $i$,
\[
\Var(r_i)
\approx
\frac{1}{n}
\frac{1+\phi^2}{1-\phi^2}.
\]
If $|\phi|$ is close to $1$, this variance can be large, meaning that high-lag sample autocorrelations can be noisy.

\subsection{Example: MA(1)}

Suppose $\{Y_t\}$ follows an MA(1) process. Then
\[
\rho_0=1,
\qquad
\rho_1=\rho_{-1}\neq 0,
\qquad
\rho_k=0
\quad
\text{for } |k|\geq 2.
\]
For $c_{11}$,
\[
c_{11}
=
\sum_{k=-\infty}^{\infty}
\left(
\rho_k^2+\rho_{k+1}\rho_{k-1}
-4\rho_1\rho_k\rho_{k-1}
+2\rho_1^2\rho_k^2
\right).
\]
Only finitely many terms are nonzero. Expanding those terms gives
\[
c_{11}
=
1-3\rho_1^2+4\rho_1^4.
\]
Therefore,
\[
r_1
\approx
\Normal\left(
\rho_1,
\frac{1-3\rho_1^2+4\rho_1^4}{n}
\right).
\]

For $i\geq 2$,
\[
c_{ii}=1+2\rho_1^2.
\]
Therefore,
\[
r_i
\approx
\Normal\left(
0,
\frac{1+2\rho_1^2}{n}
\right),
\qquad i\geq 2.
\]

\begin{exercise}
Verify that
\[
c_{ii}=1+2\rho_1^2
\]
for $i\geq 2$ in the MA(1) setting.
\end{exercise}

\begin{proof}
For an MA(1), the only nonzero autocorrelations are
\[
\rho_0=1,
\qquad
\rho_1=\rho_{-1},
\]
and all other $\rho_k$ are zero.

For $i\geq 2$, $\rho_i=0$. Hence the Bartlett diagonal term reduces to
\[
c_{ii}
=
\sum_{k=-\infty}^{\infty}
\left(
\rho_k^2+\rho_{k+i}\rho_{k-i}
\right).
\]
The first sum is
\[
\sum_{k=-\infty}^{\infty}\rho_k^2
=
\rho_{-1}^2+\rho_0^2+\rho_1^2
=
1+2\rho_1^2.
\]
For the second sum, $\rho_{k+i}\rho_{k-i}$ can be nonzero only if both $k+i$ and $k-i$ belong to $\{-1,0,1\}$. When $i\geq 2$, this cannot occur simultaneously. Therefore,
\[
\sum_{k=-\infty}^{\infty}\rho_{k+i}\rho_{k-i}=0.
\]
Thus
\[
c_{ii}=1+2\rho_1^2.
\]
\end{proof}

\subsection*{Example: MA($q$)}

For an MA($q$) process, the ACF cuts off after lag $q$:
\[
\rho_k=0
\qquad
\text{for } |k|>q.
\]
For lags $i\geq q+1$,
\[
r_i
\approx
\Normal\left(
0,
\frac{1+2\sum_{j=1}^q\rho_j^2}{n}
\right).
\]

\section{Hypothesis Testing for MA($q$)}

To test whether a series is MA($q$), use the fact that an MA($q$) process has autocorrelation zero after lag $q$.

For example, to test whether a series is MA(1), consider
\[
H_0:\text{ the series is MA(1)}
\]
versus
\[
H_a:\text{ the series is not MA(1)}.
\]
Under $H_0$, for any $i\geq 2$,
\[
r_i
\approx
\Normal\left(
0,
\frac{1+2\rho_1^2}{n}
\right).
\]
Thus, an approximate $95\%$ acceptance region is
\[
|r_i|
\leq
\frac{2}{\sqrt{n}}\sqrt{1+2\rho_1^2},
\qquad i\geq 2.
\]
The corresponding rejection region is
\[
|r_i|
>
\frac{2}{\sqrt{n}}\sqrt{1+2\rho_1^2}.
\]

If any lag $i\geq 2$ violates this bound, then we reject the MA(1) model. Since $\rho_1$ is unknown, we often replace it with the sample autocorrelation $r_1$.

If MA(1) is rejected, we can test MA(2):
\[
H_0:\text{ the series is MA(2)}
\]
versus
\[
H_a:\text{ the series is not MA(2)}.
\]
For MA(2), the rejection region for lags $i\geq 3$ is
\[
|r_i|
>
\frac{2}{\sqrt{n}}\sqrt{1+2r_1^2+2r_2^2}.
\]
More generally, for MA($q$), test lags $i\geq q+1$ using
\[
|r_i|
>
\frac{2}{\sqrt{n}}
\sqrt{
1+2\sum_{j=1}^{q}r_j^2
}.
\]

\begin{example}
Suppose $n=100$, and the sample autocorrelations are
\[
r_1=0.5,
\qquad
r_2=0.4,
\qquad
r_3=0.4,
\qquad
r_4=0.3.
\]
First test MA(0), which is white noise. The white-noise rejection region is
\[
|r_i|>\frac{2}{\sqrt{n}}=0.2.
\]
Since $r_1,r_2,r_3,r_4$ all exceed $0.2$ in absolute value, reject white noise.

Next test MA(1). The rejection threshold is
\[
\frac{2}{\sqrt{n}}\sqrt{1+2r_1^2}
=
0.2\sqrt{1+2(0.5)^2}
=
0.2\sqrt{1.5}
\approx 0.245.
\]
For MA(1), we test lags $i\geq 2$. Since
\[
|r_2|=0.4>0.245,
\]
we reject MA(1). We would then continue to test MA(2), MA(3), and so on until a plausible order is found.
\end{example}

\section*{Brief Introduction to the Partial Autocorrelation Function}

The sample ACF is useful for detecting MA behavior because an MA($q$) process has autocorrelation zero after lag $q$. For AR processes, the ACF usually decays gradually rather than cutting off. The partial autocorrelation function is designed to help identify AR order.

\begin{definition}
The partial autocorrelation function at lag $k$ is denoted by $\phi_{kk}$ and can be defined as the conditional correlation between $Y_t$ and $Y_{t-k}$ given the intermediate variables:
\[
\phi_{kk}
=
\Corr(Y_t,Y_{t-k}\mid Y_{t-1},Y_{t-2},\ldots,Y_{t-k+1}).
\]
\end{definition}

For $k=0$,
\[
\phi_{00}=1.
\]
For $k=1$, there are no intermediate variables, so
\[
\phi_{11}
=
\Corr(Y_t,Y_{t-1})
=
\rho_1.
\]

For an AR($p$) process, the PACF can be nonzero through lag $p$, but satisfies
\[
\phi_{kk}=0
\qquad
\text{for } k\geq p+1.
\]
Thus the PACF cuts off for AR processes, while the ACF cuts off for MA processes.

\begin{remark}
This ACF/PACF behavior gives a practical model identification rule:
\begin{itemize}
\item MA($q$): ACF cuts off after lag $q$, PACF decays.
\item AR($p$): PACF cuts off after lag $p$, ACF decays.
\item ARMA($p,q$): both ACF and PACF usually decay.
\end{itemize}
\end{remark}

\section{Partial Autocorrelation Function}

\subsection{Definition Using Conditional Correlation}

\begin{definition}
The partial autocorrelation function at lag $k$ is
\[
\phi_{kk}
=
\Corr(Y_t,Y_{t-k}\mid Y_{t-1},Y_{t-2},\ldots,Y_{t-k+1}).
\]
\end{definition}

Thus $\phi_{kk}$ measures the correlation between $Y_t$ and $Y_{t-k}$ after removing the linear information contained in the intermediate variables
\[
Y_{t-1},Y_{t-2},\ldots,Y_{t-k+1}.
\]

For an AR(1) process
\[
Y_t=\phi Y_{t-1}+\varepsilon_t,
\]
we have
\[
\phi_{11}=\rho_1=\phi
\]
and
\[
\phi_{kk}=0,
\qquad k\geq 2.
\]

\subsection{Definition Using Regression Residuals}

\begin{definition}
For $k\geq 2$, define
\[
\phi_{kk}
=
\Corr(\mathrm{Res}_t,\mathrm{Res}_{t-k}),
\]
where $\mathrm{Res}_t$ is the residual from the linear regression of $Y_t$ on
\[
Y_{t-1},\ldots,Y_{t-k+1},
\]
and $\mathrm{Res}_{t-k}$ is the residual from the linear regression of $Y_{t-k}$ on
\[
Y_{t-1},\ldots,Y_{t-k+1}.
\]
\end{definition}

For example, consider the AR(1) process
\[
Y_t=\phi Y_{t-1}+\varepsilon_t.
\]
To compute $\phi_{22}$, regress $Y_t$ on $Y_{t-1}$ and regress $Y_{t-2}$ on $Y_{t-1}$.

First, find $a$ minimizing
\[
\E[(Y_t-aY_{t-1})^2].
\]
Differentiate:
\[
\frac{\partial}{\partial a}\E[(Y_t-aY_{t-1})^2]
=
-2\E[(Y_t-aY_{t-1})Y_{t-1}].
\]
Setting this equal to zero gives
\[
\E[(Y_t-aY_{t-1})Y_{t-1}]=0.
\]
Thus
\[
\gamma_1-a\gamma_0=0,
\]
so
\[
a=\frac{\gamma_1}{\gamma_0}=\rho_1=\phi.
\]
Therefore,
\[
\mathrm{Res}_t=Y_t-\phi Y_{t-1}=\varepsilon_t.
\]

Similarly, find $b$ minimizing
\[
\E[(Y_{t-2}-bY_{t-1})^2].
\]
The normal equation is
\[
\E[(Y_{t-2}-bY_{t-1})Y_{t-1}]=0.
\]
Hence
\[
\gamma_1-b\gamma_0=0,
\]
so
\[
b=\frac{\gamma_1}{\gamma_0}=\rho_1=\phi.
\]
Therefore,
\[
\mathrm{Res}_{t-2}=Y_{t-2}-\phi Y_{t-1}.
\]
Using the AR(1) equation shifted one step,
\[
Y_{t-1}=\phi Y_{t-2}+\varepsilon_{t-1},
\]
we get
\[
Y_{t-2}-\phi Y_{t-1}
=
Y_{t-2}-\phi(\phi Y_{t-2}+\varepsilon_{t-1})
=
(1-\phi^2)Y_{t-2}-\phi\varepsilon_{t-1}.
\]
This residual depends only on innovations up to time $t-1$, while $\varepsilon_t$ is independent of those innovations. Hence
\[
\phi_{22}
=
\Corr(\mathrm{Res}_t,\mathrm{Res}_{t-2})
=
0.
\]
The same idea gives
\[
\phi_{kk}=0,
\qquad k\geq 2,
\]
for AR(1).

\subsection*{ACF and PACF Behavior}

\begin{xltabular}{\textwidth}{l l l l}
\toprule
 & MA($q$) & AR($p$) & ARMA($p,q$) \\
\midrule
ACF & Cuts off after $q$ & Exponential decay & Exponential decay \\
PACF & Exponential decay & Cuts off after $p$ & Exponential decay \\
\bottomrule
\end{xltabular}

\subsection{Definition Using Linear Prediction}

A computationally convenient definition of PACF is based on fitting an AR($k$) approximation. Consider the best linear prediction problem
\[
Y_t
=
\phi_{k1}Y_{t-1}
+
\phi_{k2}Y_{t-2}
+\cdots+
\phi_{kk}Y_{t-k}
+
e_t.
\]
The coefficients are chosen to minimize
\[
\E\left[
(Y_t-\phi_{k1}Y_{t-1}-\cdots-\phi_{kk}Y_{t-k})^2
\right].
\]

\begin{definition}
For a stationary time series, the PACF at lag $k$ is the last coefficient $\phi_{kk}$ in the best linear AR($k$) approximation to $Y_t$.
\end{definition}

The normal equations for this least squares problem are
\[
\Gamma_k\bm{\phi}_k=\bm{\gamma}_k,
\]
where
\[
\Gamma_k
=
\begin{bmatrix}
\gamma_0 & \gamma_1 & \cdots & \gamma_{k-1}\\
\gamma_1 & \gamma_0 & \cdots & \gamma_{k-2}\\
\vdots & \vdots & \ddots & \vdots\\
\gamma_{k-1} & \gamma_{k-2} & \cdots & \gamma_0
\end{bmatrix},
\qquad
\bm{\phi}_k
=
\begin{bmatrix}
\phi_{k1}\\
\phi_{k2}\\
\vdots\\
\phi_{kk}
\end{bmatrix},
\qquad
\bm{\gamma}_k
=
\begin{bmatrix}
\gamma_1\\
\gamma_2\\
\vdots\\
\gamma_k
\end{bmatrix}.
\]
Equivalently, dividing through by $\gamma_0$,
\[
R_k\bm{\phi}_k=\bm{\rho}_k,
\]
where
\[
R_k
=
\begin{bmatrix}
1 & \rho_1 & \cdots & \rho_{k-1}\\
\rho_1 & 1 & \cdots & \rho_{k-2}\\
\vdots & \vdots & \ddots & \vdots\\
\rho_{k-1} & \rho_{k-2} & \cdots & 1
\end{bmatrix},
\qquad
\bm{\rho}_k
=
\begin{bmatrix}
\rho_1\\
\rho_2\\
\vdots\\
\rho_k
\end{bmatrix}.
\]
Thus
\[
\bm{\phi}_k=R_k^{-1}\bm{\rho}_k,
\]
and the PACF $\phi_{kk}$ is the last entry of $R_k^{-1}\bm{\rho}_k$.

For an AR($p$) process, the true model is
\[
Y_t=\phi_1Y_{t-1}+\cdots+\phi_pY_{t-p}+\varepsilon_t.
\]
If $k=p$, then
\[
\phi_{pp}=\phi_p.
\]
If $k>p$, then the best AR($k$) approximation has coefficients
\[
(\phi_1,\ldots,\phi_p,0,\ldots,0),
\]
so
\[
\phi_{kk}=0,
\qquad k>p.
\]

\section{Durbin--Levinson Recursion}

Although the PACF can be computed by directly inverting $R_k$, this becomes expensive for large $k$. The Durbin--Levinson recursion computes PACF values recursively by exploiting the Toeplitz structure of the autocorrelation matrix.

\begin{definition}
Define
\[
\phi_{00}=1.
\]
For $l\geq 0$, the Durbin--Levinson recursion is
\[
\phi_{l+1,l+1}
=
\frac{
\rho_{l+1}-\sum_{j=1}^{l}\phi_{lj}\rho_{l+1-j}
}{
1-\sum_{j=1}^{l}\phi_{lj}\rho_j
},
\]
and, for $1\leq j\leq l$,
\[
\phi_{l+1,j}
=
\phi_{lj}
-
\phi_{l+1,l+1}\phi_{l,l+1-j}.
\]
\end{definition}

\begin{example}
For an AR(1) process
\[
Y_t=cY_{t-1}+\varepsilon_t,
\]
the autocorrelation function is
\[
\rho_k=c^k.
\]
For $l=0$,
\[
\phi_{11}
=
\rho_1
=
c.
\]
For $l=1$,
\[
\phi_{22}
=
\frac{\rho_2-\phi_{11}\rho_1}{1-\phi_{11}\rho_1}
=
\frac{c^2-c^2}{1-c^2}
=
0.
\]
Also,
\[
\phi_{21}
=
\phi_{11}-\phi_{22}\phi_{11}
=
c.
\]
For $l=2$,
\[
\phi_{33}
=
\frac{\rho_3-(\phi_{21}\rho_2+\phi_{22}\rho_1)}
{1-(\phi_{21}\rho_1+\phi_{22}\rho_2)}
=
\frac{c^3-(c\cdot c^2+0\cdot c)}
{1-(c\cdot c+0\cdot c^2)}
=
0.
\]
Thus the PACF cuts off after lag $1$, as expected.
\end{example}

\begin{exercise}
Find $\phi_{44}$ for an AR(1) process using the Durbin--Levinson recursion.
\end{exercise}

\begin{proof}
From the previous calculations for an AR(1),
\[
\phi_{33}=0,
\qquad
\phi_{31}=c,
\qquad
\phi_{32}=0.
\]
Using the recursion with $l=3$,
\[
\phi_{44}
=
\frac{
\rho_4-(\phi_{31}\rho_3+\phi_{32}\rho_2+\phi_{33}\rho_1)
}{
1-(\phi_{31}\rho_1+\phi_{32}\rho_2+\phi_{33}\rho_3)
}.
\]
Since $\rho_k=c^k$,
\[
\phi_{44}
=
\frac{
c^4-(c\cdot c^3+0\cdot c^2+0\cdot c)
}{
1-(c\cdot c+0\cdot c^2+0\cdot c^3)
}
=
\frac{c^4-c^4}{1-c^2}
=
0.
\]
\end{proof}

\section{Sample PACF}

There are two common ways to compute the sample PACF.

\subsection{Method 1: Replacing the ACF in Durbin--Levinson}

Since the Durbin--Levinson recursion expresses the theoretical PACF in terms of the theoretical ACF, we can estimate PACF values by replacing
\[
\rho_k
\]
with the sample ACF
\[
r_k.
\]
This gives sample PACF estimates
\[
\widehat{\phi}_{kk}.
\]

\subsection{Method 2: Solving the Empirical Linear System}

Alternatively, replace the theoretical autocorrelation matrix by the sample autocorrelation matrix:
\[
\widehat{R}_k
=
\begin{bmatrix}
1 & r_1 & \cdots & r_{k-1}\\
r_1 & 1 & \cdots & r_{k-2}\\
\vdots & \vdots & \ddots & \vdots\\
r_{k-1} & r_{k-2} & \cdots & 1
\end{bmatrix},
\qquad
\widehat{\bm{\rho}}_k
=
\begin{bmatrix}
r_1\\
r_2\\
\vdots\\
r_k
\end{bmatrix}.
\]
Then
\[
\widehat{R}_k\widehat{\bm{\phi}}_k
=
\widehat{\bm{\rho}}_k.
\]
Solving gives
\[
\widehat{\bm{\phi}}_k
=
\widehat{R}_k^{-1}\widehat{\bm{\rho}}_k,
\]
and the sample PACF $\widehat{\phi}_{kk}$ is the last entry.

\begin{remark}
The matrix $\widehat{R}_k$ constructed from the sample ACF is invertible under the usual sample ACF construction. This is why using the sample ACF in the PACF equations is numerically stable.
\end{remark}

For an AR($p$) process, the large-sample distribution of the sample PACF satisfies
\[
\widehat{\phi}_{kk}
\approx
\Normal\left(0,\frac{1}{n}\right),
\qquad k\geq p+1.
\]
Therefore, to test whether a series is AR($p$), we can use the approximate rejection rule
\[
|\widehat{\phi}_{kk}|>\frac{2}{\sqrt{n}},
\qquad k\geq p+1.
\]

\section{Extended Autocorrelation Function}

\begin{definition}
The extended autocorrelation function, abbreviated EACF, is a diagnostic tool for identifying the order $(p,q)$ of an ARMA($p,q$) model.
\end{definition}

The idea is to fit AR filters of different orders and then examine whether the residuals behave like an MA process. If the true model is ARMA($p,q$), then after applying the correct AR filter, the remaining residual series should have MA($q$)-type behavior.

The EACF algorithm can be summarized as follows.

\begin{algorithm}[H]
\caption{EACF Order Identification}
\begin{algorithmic}[1]
\For{$p=0,1,\ldots,p_{\max}$}
    \State Fit an AR($p$) model to $\{Y_t\}$ and obtain residuals $\widehat{W}_t$.
    \For{$q=0,1,\ldots,q_{\max}$}
        \State Test whether the residual series $\widehat{W}_t$ is consistent with an MA($q$) process.
        \State Record ``O'' if the MA($q$) null is not rejected, and ``X'' otherwise.
    \EndFor
\EndFor
\State Look for a triangular region of ``O'' entries.
\end{algorithmic}
\end{algorithm}

For an ARMA($p,q$) process, overfitting the AR part or the MA part can still produce acceptable residual behavior. Hence the EACF table often contains a triangular region of zeros or ``O'' symbols. The upper-left corner of this triangle is used as a candidate estimate of $(p,q)$.

\begin{remark}
In R, the EACF can be computed using functions from the \texttt{TSA} package.
\end{remark}

\section{Unit Root Test: Augmented Dickey--Fuller Test}

A unit root test is used to assess whether a time series is nonstationary because its AR polynomial has a root at $1$.

The augmented Dickey--Fuller test is a common unit root test. It tests
\[
H_0:\text{ the AR polynomial has a unit root}
\]
against
\[
H_a:\text{ the AR polynomial does not have a unit root}.
\]
Equivalently,
\[
H_0:\{Y_t\}\text{ is nonstationary}
\]
against
\[
H_a:\{Y_t\}\text{ is stationary}.
\]

\begin{remark}
In R, the ADF test can be computed using
\[
\texttt{adf.test()}
\]
from the \texttt{tseries} package. A small $p$-value gives evidence against the unit-root null and supports stationarity. A large $p$-value does not provide evidence against nonstationarity.
\end{remark}

\chapter{Parameter Estimation}

After using tools such as the sample ACF, sample PACF, EACF, transformations, and unit root tests, we arrive at a few candidate models. The next task is to estimate the unknown parameters:
\[
\phi_1,\ldots,\phi_p,
\qquad
\theta_1,\ldots,\theta_q,
\qquad
\sigma_\varepsilon^2.
\]

Common estimation methods include:
\begin{itemize}
\item method of moments,
\item least squares,
\item maximum likelihood.
\end{itemize}

\section{Method of Moments}

The method of moments estimates parameters by matching theoretical moments to sample moments:
\[
\text{theoretical moment}
=
\text{sample moment}.
\]

Suppose the theoretical $k$th moment is
\[
\mu_k=\E[Y^k],
\]
and the sample $k$th moment is
\[
m_k=\frac{1}{n}\sum_{i=1}^nY_i^k.
\]
Then the method of moments solves
\[
\mu_k=m_k
\]
for the unknown parameters.

\subsection{IID Normal Example}

Suppose
\[
Y_1,\ldots,Y_n\iid \Normal(\mu,\sigma^2).
\]
Then
\[
\mu_1=\E[Y]=\mu,
\]
and
\[
\mu_2=\E[Y^2]=\mu^2+\sigma^2.
\]
The sample moments are
\[
m_1=\frac{1}{n}\sum_{i=1}^nY_i=\overline{Y},
\]
and
\[
m_2=\frac{1}{n}\sum_{i=1}^nY_i^2.
\]
Matching moments gives
\[
\mu=m_1
\]
and
\[
\mu^2+\sigma^2=m_2.
\]
Therefore,
\[
\widehat{\mu}_{\mathrm{MOM}}=\overline{Y},
\]
and
\[
\widehat{\sigma}_{\mathrm{MOM}}^2
=
m_2-m_1^2
=
\frac{1}{n}\sum_{i=1}^nY_i^2-\overline{Y}^{\,2}.
\]
Equivalently,
\[
\widehat{\sigma}_{\mathrm{MOM}}^2
=
\frac{1}{n}\sum_{i=1}^n(Y_i-\overline{Y})^2.
\]

\begin{exercise}
Verify the final equality above.
\end{exercise}

\begin{proof}
Expand:
\[
\frac{1}{n}\sum_{i=1}^n(Y_i-\overline{Y})^2
=
\frac{1}{n}\sum_{i=1}^n(Y_i^2-2Y_i\overline{Y}+\overline{Y}^{\,2}).
\]
Thus
\[
\frac{1}{n}\sum_{i=1}^n(Y_i-\overline{Y})^2
=
\frac{1}{n}\sum_{i=1}^nY_i^2
-
2\overline{Y}\frac{1}{n}\sum_{i=1}^nY_i
+
\overline{Y}^{\,2}.
\]
Since
\[
\frac{1}{n}\sum_{i=1}^nY_i=\overline{Y},
\]
we get
\[
\frac{1}{n}\sum_{i=1}^n(Y_i-\overline{Y})^2
=
\frac{1}{n}\sum_{i=1}^nY_i^2
-
2\overline{Y}^{\,2}
+
\overline{Y}^{\,2}
=
\frac{1}{n}\sum_{i=1}^nY_i^2-\overline{Y}^{\,2}.
\]
\end{proof}

\subsection{Generalized Method of Moments for Time Series}

For time series models, we often match theoretical autocorrelations to sample autocorrelations:
\[
\text{theoretical ACF}
=
\text{sample ACF}.
\]
This is a version of generalized method of moments.

\subsection*{Example: AR(1)}

Suppose $\{Y_t\}$ is a mean-zero AR(1) process:
\[
Y_t-\phi Y_{t-1}=\varepsilon_t.
\]
The theoretical autocorrelation is
\[
\rho_k=\phi^k.
\]
In particular,
\[
\rho_1=\phi.
\]
Matching $\rho_1$ to the sample autocorrelation $r_1$ gives
\[
\widehat{\phi}_{\mathrm{MOM}}=r_1.
\]
Using the formula for $r_1$,
\[
\widehat{\phi}_{\mathrm{MOM}}
=
\frac{\sum_{t=2}^n(Y_t-\overline{Y})(Y_{t-1}-\overline{Y})}
{\sum_{t=1}^n(Y_t-\overline{Y})^2}.
\]

\subsection*{Example: AR(2)}

Suppose $\{Y_t\}$ follows a mean-zero AR(2) process:
\[
Y_t-\phi_1Y_{t-1}-\phi_2Y_{t-2}=\varepsilon_t.
\]
The first two Yule--Walker equations for autocorrelations are
\[
\rho_1=\phi_1+\phi_2\rho_1,
\]
and
\[
\rho_2=\phi_1\rho_1+\phi_2.
\]
From the first equation,
\[
\phi_1=\rho_1(1-\phi_2).
\]
Substitute into the second equation:
\[
\rho_2=\rho_1^2(1-\phi_2)+\phi_2.
\]
Thus
\[
\rho_2-\rho_1^2
=
\phi_2(1-\rho_1^2),
\]
so
\[
\phi_2
=
\frac{\rho_2-\rho_1^2}{1-\rho_1^2}.
\]
Then
\[
\phi_1
=
\rho_1(1-\phi_2)
=
\rho_1
\left(
1-\frac{\rho_2-\rho_1^2}{1-\rho_1^2}
\right)
=
\frac{\rho_1(1-\rho_2)}{1-\rho_1^2}.
\]

Matching theoretical autocorrelations to sample autocorrelations gives
\[
\widehat{\phi}_{2,\mathrm{MOM}}
=
\frac{r_2-r_1^2}{1-r_1^2},
\]
and
\[
\widehat{\phi}_{1,\mathrm{MOM}}
=
\frac{r_1(1-r_2)}{1-r_1^2}.
\]

\begin{exercise}
Verify the AR(2) method-of-moments formulas.
\end{exercise}

\begin{proof}
Starting from
\[
\rho_1=\phi_1+\phi_2\rho_1,
\]
we have
\[
\phi_1=\rho_1-\phi_2\rho_1=\rho_1(1-\phi_2).
\]
Substitute this into
\[
\rho_2=\phi_1\rho_1+\phi_2.
\]
Then
\[
\rho_2=\rho_1^2(1-\phi_2)+\phi_2
=
\rho_1^2-\rho_1^2\phi_2+\phi_2.
\]
Thus
\[
\rho_2-\rho_1^2
=
\phi_2(1-\rho_1^2),
\]
so
\[
\phi_2=
\frac{\rho_2-\rho_1^2}{1-\rho_1^2}.
\]
Then
\[
\phi_1
=
\rho_1(1-\phi_2)
=
\rho_1
\left[
1-
\frac{\rho_2-\rho_1^2}{1-\rho_1^2}
\right]
=
\rho_1
\left[
\frac{1-\rho_1^2-\rho_2+\rho_1^2}{1-\rho_1^2}
\right]
=
\frac{\rho_1(1-\rho_2)}{1-\rho_1^2}.
\]
Replacing $\rho_1,\rho_2$ by $r_1,r_2$ gives the method-of-moments estimators.
\end{proof}

\subsection{Example: AR($p$)}

For a mean-zero AR($p$) process,
\[
Y_t-\phi_1Y_{t-1}-\cdots-\phi_pY_{t-p}=\varepsilon_t.
\]
The Yule--Walker autocorrelation equations are
\[
\rho_k
=
\phi_1\rho_{k-1}
+
\phi_2\rho_{k-2}
+\cdots+
\phi_p\rho_{k-p},
\qquad k=1,\ldots,p,
\]
with $\rho_{-j}=\rho_j$ and $\rho_0=1$.

In matrix form,
\[
R_p\bm{\phi}=\bm{\rho}_p,
\]
where
\[
R_p
=
\begin{bmatrix}
1 & \rho_1 & \cdots & \rho_{p-1}\\
\rho_1 & 1 & \cdots & \rho_{p-2}\\
\vdots & \vdots & \ddots & \vdots\\
\rho_{p-1} & \rho_{p-2} & \cdots & 1
\end{bmatrix},
\qquad
\bm{\phi}
=
\begin{bmatrix}
\phi_1\\
\phi_2\\
\vdots\\
\phi_p
\end{bmatrix},
\qquad
\bm{\rho}_p
=
\begin{bmatrix}
\rho_1\\
\rho_2\\
\vdots\\
\rho_p
\end{bmatrix}.
\]
Replacing theoretical autocorrelations by sample autocorrelations gives
\[
\widehat{R}_p\widehat{\bm{\phi}}_{\mathrm{MOM}}
=
\widehat{\bm{\rho}}_p,
\]
so
\[
\widehat{\bm{\phi}}_{\mathrm{MOM}}
=
\widehat{R}_p^{-1}\widehat{\bm{\rho}}_p.
\]

\begin{remark}
The matrix $\widehat{R}_p$ constructed from the sample ACF is invertible by the same property used when defining the sample PACF.
\end{remark}

\section{Method of Moments for Other Models and Parameters}

Recall that in method of moments or generalized method of moments, we solve equations of the form
\[
\text{theoretical moment}=\text{sample moment},
\]
or
\[
\text{theoretical ACF}=\text{sample ACF}.
\]
The theoretical moments and autocorrelations are functions of unknown parameters, while the sample moments and sample autocorrelations are computed from the observed data $Y_1,\ldots,Y_n$.

\subsection*{Method of Moments for MA(1)}

Suppose $\{Y_t\}$ follows an MA(1) model
\[
Y_t=\varepsilon_t-\theta\varepsilon_{t-1}.
\]
For an MA(1), the lag-one autocorrelation is
\[
\rho_1=\frac{-\theta}{1+\theta^2}.
\]
The method of moments sets
\[
\rho_1=r_1.
\]
Thus
\[
\frac{-\theta}{1+\theta^2}=r_1.
\]
Multiplying through gives
\[
r_1\theta^2+\theta+r_1=0.
\]
If $r_1=0$, then
\[
\widehat{\theta}_{\mathrm{MOM}}=0.
\]
If $r_1\neq 0$, the quadratic formula gives
\[
\theta
=
\frac{-1\pm\sqrt{1-4r_1^2}}{2r_1}.
\]
Therefore,
\[
\widehat{\theta}_{1,2}
=
\frac{-1\pm\sqrt{1-4r_1^2}}{2r_1}.
\]

The method-of-moments estimator does not always exist. For an MA(1), the theoretical ACF satisfies
\[
|\rho_1|\leq \frac{1}{2}.
\]
Thus if
\[
|r_1|>\frac{1}{2},
\]
then the quadratic equation has no real solution.

When
\[
0<|r_1|<\frac{1}{2},
\]
there are two real solutions. One corresponds to an invertible MA(1), and the other corresponds to a noninvertible MA(1). We choose the solution satisfying
\[
|\theta|<1.
\]
This gives
\[
\widehat{\theta}_{\mathrm{MOM}}
=
\frac{-1+\sqrt{1-4r_1^2}}{2r_1}.
\]

\begin{remark}
For MA($q$), method-of-moments equations are usually nonlinear and may have multiple solutions. The invertibility condition is used to choose the admissible solution.
\end{remark}

\subsection{Method of Moments Estimate for $\mu$}

If the time series has constant mean $\mu$, then the method-of-moments estimator is
\[
\widehat{\mu}_{\mathrm{MOM}}
=
\frac{1}{n}\sum_{t=1}^nY_t
=
\overline{Y}.
\]

\subsection{Method of Moments Estimate for $\sigma_\varepsilon^2$}

Suppose we want to estimate the innovation variance $\sigma_\varepsilon^2$. The general idea is:

\begin{enumerate}
\item Express $\gamma_0$ in terms of the model parameters and $\sigma_\varepsilon^2$.
\item Solve for $\sigma_\varepsilon^2$ in terms of $\gamma_0$ and the other parameters.
\item Replace theoretical quantities with their sample or estimated versions.
\end{enumerate}

For an AR($p$) process,
\[
Y_t-\phi_1Y_{t-1}-\cdots-\phi_pY_{t-p}=\varepsilon_t.
\]
The $0$th Yule--Walker equation is
\[
\gamma_0-\phi_1\gamma_1-\phi_2\gamma_2-\cdots-\phi_p\gamma_p
=
\sigma_\varepsilon^2.
\]
Therefore,
\[
\sigma_\varepsilon^2
=
\gamma_0(1-\phi_1\rho_1-\phi_2\rho_2-\cdots-\phi_p\rho_p).
\]
Replacing the unknown quantities by estimates gives
\[
\widehat{\sigma}_{\varepsilon,\mathrm{MOM}}^2
=
s^2
\left(
1-\widehat{\phi}_1r_1-\widehat{\phi}_2r_2-\cdots-\widehat{\phi}_pr_p
\right),
\]
where
\[
s^2
=
\frac{1}{n-1}\sum_{t=1}^n(Y_t-\overline{Y})^2.
\]

\section{Least Squares Estimation}

Least squares estimation constructs an objective function as a sum of squared residuals and estimates parameters by minimizing that sum.

\subsection{Conditional Least Squares for AR(1)}

Consider an AR(1) process with mean $\mu$:
\[
Y_t-\mu=\phi(Y_{t-1}-\mu)+\varepsilon_t,
\qquad
\varepsilon_t\iid \Normal(0,\sigma_\varepsilon^2).
\]
Given observations $Y_1,\ldots,Y_n$, the conditional least squares objective is
\[
S_c(\phi,\mu)
=
\sum_{t=2}^n
\left[
(Y_t-\mu)-\phi(Y_{t-1}-\mu)
\right]^2.
\]
The subscript $c$ indicates that this is conditional on the initial observation $Y_1$.

The conditional least squares estimator is
\[
(\widehat{\phi}_{\mathrm{LS}},\widehat{\mu}_{\mathrm{LS}})
=
\argmin_{\phi,\mu}S_c(\phi,\mu).
\]

To minimize with respect to $\mu$, differentiate:
\[
\frac{\partial S_c}{\partial \mu}
=
2\sum_{t=2}^n
\left[
(Y_t-\mu)-\phi(Y_{t-1}-\mu)
\right](-1+\phi).
\]
For a stationary AR(1), $\phi\neq 1$. Setting this derivative equal to zero gives
\[
\sum_{t=2}^n
\left[
(Y_t-\mu)-\phi(Y_{t-1}-\mu)
\right]
=
0.
\]
Therefore,
\[
\sum_{t=2}^nY_t-(n-1)\mu
=
\phi\left(
\sum_{t=2}^nY_{t-1}-(n-1)\mu
\right).
\]
Solving for $\mu$ gives
\[
\widehat{\mu}_{\mathrm{LS}}
=
\frac{
\sum_{t=2}^nY_t-\phi\sum_{t=2}^nY_{t-1}
}{
(n-1)(1-\phi)
}.
\]
Since
\[
\sum_{t=2}^nY_{t-1}
=
\sum_{t=1}^{n-1}Y_t,
\]
we can also write
\[
\widehat{\mu}_{\mathrm{LS}}
=
\frac{
\sum_{t=2}^nY_t-\phi\sum_{t=1}^{n-1}Y_t
}{
(n-1)(1-\phi)
}.
\]
For large $n$, ignoring endpoint effects,
\[
\widehat{\mu}_{\mathrm{LS}}\approx \overline{Y}.
\]

Now differentiate with respect to $\phi$:
\[
\frac{\partial S_c}{\partial \phi}
=
-2\sum_{t=2}^n
\left[
(Y_t-\mu)-\phi(Y_{t-1}-\mu)
\right](Y_{t-1}-\mu).
\]
Setting this derivative equal to zero gives
\[
\sum_{t=2}^n
\left[
(Y_t-\mu)-\phi(Y_{t-1}-\mu)
\right](Y_{t-1}-\mu)=0.
\]
If we use the large-sample approximation $\widehat{\mu}_{\mathrm{LS}}\approx\overline{Y}$, then
\[
\widehat{\phi}_{\mathrm{LS}}
\approx
\frac{
\sum_{t=2}^n(Y_t-\overline{Y})(Y_{t-1}-\overline{Y})
}{
\sum_{t=2}^n(Y_{t-1}-\overline{Y})^2
}.
\]
For large $n$, this is approximately
\[
\widehat{\phi}_{\mathrm{LS}}\approx r_1.
\]
Thus the conditional least squares estimator for AR(1) is asymptotically close to the method-of-moments estimator.

\subsection*{Conditional Least Squares for AR(2)}

Consider an AR(2) process with mean $\mu$:
\[
Y_t-\mu
=
\phi_1(Y_{t-1}-\mu)+\phi_2(Y_{t-2}-\mu)+\varepsilon_t.
\]
Given observations $Y_1,\ldots,Y_n$, define
\[
S_c(\phi_1,\phi_2,\mu)
=
\sum_{t=3}^n
\left[
(Y_t-\mu)
-\phi_1(Y_{t-1}-\mu)
-\phi_2(Y_{t-2}-\mu)
\right]^2.
\]
For large $n$, the conditional least squares estimates satisfy approximately
\[
\widehat{\mu}_{\mathrm{LS}}\approx\overline{Y},
\]
and
\[
\widehat{\phi}_{1,\mathrm{LS}}\approx\widehat{\phi}_{1,\mathrm{MOM}},
\qquad
\widehat{\phi}_{2,\mathrm{LS}}\approx\widehat{\phi}_{2,\mathrm{MOM}}.
\]

\subsection{Conditional Least Squares for MA(1)}

Consider an MA(1) process:
\[
Y_t=\varepsilon_t-\theta\varepsilon_{t-1}.
\]
For MA models, the residuals $\varepsilon_t$ cannot be written as a finite function of the observed data unless we impose an initial condition. Under invertibility, we can write
\[
\varepsilon_t
=
Y_t+\theta Y_{t-1}+\theta^2Y_{t-2}+\cdots.
\]
Given observations $Y_1,\ldots,Y_n$, set $\varepsilon_0=0$ and recursively define
\[
\varepsilon_t(\theta)
=
Y_t+\theta\varepsilon_{t-1}(\theta).
\]
Then
\[
\varepsilon_1(\theta)=Y_1,
\]
\[
\varepsilon_2(\theta)=Y_2+\theta Y_1,
\]
\[
\varepsilon_3(\theta)=Y_3+\theta Y_2+\theta^2Y_1,
\]
and in general,
\[
\varepsilon_t(\theta)
=
Y_t+\theta Y_{t-1}+\theta^2Y_{t-2}+\cdots+\theta^{t-1}Y_1.
\]
The conditional least squares objective is
\[
S_c(\theta)
=
\sum_{t=1}^n
\varepsilon_t(\theta)^2
=
\sum_{t=1}^n
\left(
Y_t+\theta Y_{t-1}+\theta^2Y_{t-2}+\cdots+\theta^{t-1}Y_1
\right)^2.
\]
The conditional least squares estimate is
\[
\widehat{\theta}_{\mathrm{LS}}
=
\argmin_\theta S_c(\theta).
\]
This objective is nonlinear in $\theta$, so it is usually minimized numerically.

\subsection{Conditional Least Squares for ARMA($p,q$)}

For an ARMA($p,q$) process,
\[
\Phi(B)Y_t=\Theta(B)\varepsilon_t.
\]
Given parameters, residuals can be computed recursively after choosing initial residual values. A common conditional least squares approach sets
\[
\varepsilon_0=\varepsilon_{-1}=\cdots=\varepsilon_{-q+1}=0.
\]
Then define residuals recursively and minimize
\[
S_c(\phi_1,\ldots,\phi_p,\theta_1,\ldots,\theta_q)
=
\sum_t \varepsilon_t^2.
\]

For example, for ARMA(1,1),
\[
Y_t-\phi Y_{t-1}=\varepsilon_t-\theta\varepsilon_{t-1}.
\]
Setting $\varepsilon_1=0$ for initialization, we compute recursively:
\[
\varepsilon_2
=
\theta\varepsilon_1+Y_2-\phi Y_1
=
Y_2-\phi Y_1,
\]
\[
\varepsilon_3
=
\theta\varepsilon_2+Y_3-\phi Y_2,
\]
and in general,
\[
\varepsilon_t
=
\theta\varepsilon_{t-1}+Y_t-\phi Y_{t-1}.
\]
Then the objective is
\[
S_c(\phi,\theta)
=
\sum_{t=2}^n \varepsilon_t^2.
\]

\section{Maximum Likelihood Estimation}

Maximum likelihood uses the joint distribution of the observed data. The likelihood function is
\[
L(\text{parameters}\mid Y_1,\ldots,Y_n)
=
f(Y_1,\ldots,Y_n\mid \text{parameters}).
\]
The maximum likelihood estimator maximizes $L$, or equivalently maximizes the log-likelihood
\[
\ell=\log L.
\]

Assume
\[
\varepsilon_t\iid \Normal(0,\sigma_\varepsilon^2).
\]
Then the density of a single innovation is
\[
f(\varepsilon_t)
=
(2\pi\sigma_\varepsilon^2)^{-1/2}
\exp\left(
-\frac{\varepsilon_t^2}{2\sigma_\varepsilon^2}
\right).
\]
For independent innovations,
\[
f(\varepsilon_1,\ldots,\varepsilon_n)
=
(2\pi\sigma_\varepsilon^2)^{-n/2}
\exp\left(
-\frac{1}{2\sigma_\varepsilon^2}
\sum_{t=1}^n\varepsilon_t^2
\right).
\]

\subsection{MLE for AR(1)}

Consider
\[
Y_t-\mu=\phi(Y_{t-1}-\mu)+\varepsilon_t,
\qquad
\varepsilon_t\iid \Normal(0,\sigma_\varepsilon^2).
\]
For a stationary AR(1),
\[
Y_1\sim \Normal\left(\mu,\frac{\sigma_\varepsilon^2}{1-\phi^2}\right).
\]
Hence
\[
f(Y_1)
=
\left(2\pi\frac{\sigma_\varepsilon^2}{1-\phi^2}\right)^{-1/2}
\exp\left[
-\frac{(1-\phi^2)(Y_1-\mu)^2}{2\sigma_\varepsilon^2}
\right].
\]

For $t\geq 2$,
\[
Y_t\mid Y_{t-1}
\sim
\Normal\left(
\mu+\phi(Y_{t-1}-\mu),
\sigma_\varepsilon^2
\right).
\]
Thus
\[
f(Y_2,\ldots,Y_n\mid Y_1)
=
\prod_{t=2}^n
(2\pi\sigma_\varepsilon^2)^{-1/2}
\exp\left[
-\frac{
\left((Y_t-\mu)-\phi(Y_{t-1}-\mu)\right)^2
}{2\sigma_\varepsilon^2}
\right].
\]
Combining the marginal density of $Y_1$ and the conditional densities gives
\[
L(\mu,\phi,\sigma_\varepsilon^2)
=
(2\pi\sigma_\varepsilon^2)^{-n/2}
(1-\phi^2)^{1/2}
\exp\left[
-\frac{1}{2\sigma_\varepsilon^2}
S(\mu,\phi)
\right],
\]
where
\[
S(\mu,\phi)
=
S_c(\mu,\phi)
+
(1-\phi^2)(Y_1-\mu)^2,
\]
and
\[
S_c(\mu,\phi)
=
\sum_{t=2}^n
\left[
(Y_t-\mu)-\phi(Y_{t-1}-\mu)
\right]^2.
\]

Therefore,
\[
\ell(\mu,\phi,\sigma_\varepsilon^2)
=
-\frac{n}{2}\log(2\pi\sigma_\varepsilon^2)
+
\frac{1}{2}\log(1-\phi^2)
-
\frac{1}{2\sigma_\varepsilon^2}S(\mu,\phi).
\]

Differentiating with respect to $\sigma_\varepsilon^2$,
\[
\frac{\partial \ell}{\partial \sigma_\varepsilon^2}
=
-\frac{n}{2\sigma_\varepsilon^2}
+
\frac{S(\mu,\phi)}{2(\sigma_\varepsilon^2)^2}.
\]
Setting this equal to zero gives
\[
\widehat{\sigma}_{\varepsilon,\mathrm{MLE}}^2
=
\frac{1}{n}S(\widehat{\mu}_{\mathrm{MLE}},\widehat{\phi}_{\mathrm{MLE}}).
\]

\begin{remark}
The MLE of $\sigma_\varepsilon^2$ is biased in finite samples. A common approximately unbiased alternative is
\[
\widehat{\sigma}_\varepsilon^2
=
\frac{1}{n-2}S(\widehat{\mu},\widehat{\phi}),
\]
because two parameters, $\mu$ and $\phi$, were estimated.
\end{remark}

The partial derivative with respect to $\mu$ is
\[
\frac{\partial \ell}{\partial \mu}
=
-\frac{1}{2\sigma_\varepsilon^2}
\frac{\partial S(\mu,\phi)}{\partial \mu}.
\]
Thus the likelihood equation for $\mu$ is
\[
\frac{\partial S(\mu,\phi)}{\partial \mu}=0.
\]

The partial derivative with respect to $\phi$ is
\[
\frac{\partial \ell}{\partial \phi}
=
\frac{\partial}{\partial \phi}
\left[
\frac{1}{2}\log(1-\phi^2)
-
\frac{1}{2\sigma_\varepsilon^2}S(\mu,\phi)
\right].
\]
In general, solving
\[
\frac{\partial \ell}{\partial \phi}=0
\]
requires numerical optimization.

\subsection{Unconditional Sum of Squares}

Maximum likelihood for ARMA models can be computationally difficult. Since the AR(1) likelihood above depends on $S(\mu,\phi)$, one simpler alternative is to minimize
\[
S(\mu,\phi)
=
S_c(\mu,\phi)
+
(1-\phi^2)(Y_1-\mu)^2.
\]
This method is called unconditional sum of squares.

The unconditional sum of squares estimator is
\[
(\widehat{\mu},\widehat{\phi})
=
\argmin_{\mu,\phi}S(\mu,\phi).
\]
It differs from exact MLE because MLE also includes the term
\[
\frac{1}{2}\log(1-\phi^2)
\]
in the log-likelihood. However, for large sample sizes, the two methods are often close.

\section{Comparison of Estimation Methods}

The main estimation methods are:

\begin{itemize}
\item maximum likelihood estimation: maximizes the likelihood or log-likelihood,
\item unconditional sum of squares: minimizes the unconditional sum of squares,
\item conditional sum of squares: minimizes the conditional sum of squares.
\end{itemize}

For large samples, these methods often give similar estimates. In small or moderate samples, maximum likelihood is usually preferred because it uses more of the distributional information.

\begin{remark}
Conditional least squares can be viewed as a simplified version of likelihood estimation that conditions on initial observations or sets initial residuals to zero.
\end{remark}

In general,
\[
\argmin S_c(\mu,\phi,\theta)
\neq
\argmax L(\mu,\phi,\theta),
\]
especially for ARMA models near nonstationarity or noninvertibility.

Maximum likelihood is conceptually better because it uses the full distributional model, but it is computationally harder and usually requires numerical optimization.

\section*{Asymptotic Theory}

In large samples, estimators such as MLE, conditional least squares, and unconditional least squares are typically:

\begin{itemize}
\item consistent,
\item asymptotically unbiased,
\item asymptotically normal.
\end{itemize}

For an AR(1) model,
\[
\Var(\widehat{\phi})
\approx
\frac{1-\phi^2}{n}.
\]
Therefore,
\[
\widehat{\phi}
\approx
\Normal\left(\phi,\frac{1-\phi^2}{n}\right).
\]

For an AR(2) model, approximately,
\[
\Var(\widehat{\phi}_1)
\approx
\Var(\widehat{\phi}_2)
\approx
\frac{1-\phi_2^2}{n},
\]
and
\[
\Corr(\widehat{\phi}_1,\widehat{\phi}_2)
\approx
-\frac{\phi_1}{1-\phi_2}.
\]

For an MA(1) model,
\[
\widehat{\theta}
\approx
\Normal\left(\theta,\frac{1-\theta^2}{n}\right).
\]

For an ARMA(1,1) model,
\[
\Var(\widehat{\phi})
\approx
\frac{1-\phi^2}{n}
\left(
\frac{1-\phi\theta}{\phi-\theta}
\right)^2,
\]
and
\[
\Var(\widehat{\theta})
\approx
\frac{1-\theta^2}{n}
\left(
\frac{1-\phi\theta}{\phi-\theta}
\right)^2.
\]

\begin{remark}
If $\phi\approx\theta$ in an ARMA(1,1) model, then the variances above become very large. This occurs because the AR and MA terms nearly cancel:
\[
Y_t-\phi Y_{t-1}
=
\varepsilon_t-\theta\varepsilon_{t-1},
\]
and when $\phi=\theta$, the model reduces to
\[
Y_t=\varepsilon_t.
\]
Near cancellation makes the individual AR and MA parameters difficult to estimate separately.
\end{remark}

\section{Overfitting as a Tool}

Overfitting can be used diagnostically. Suppose the true process is AR(1):
\[
Y_t-\phi Y_{t-1}=\varepsilon_t.
\]
If we fit the correct AR(1) model, then
\[
\widehat{\phi}
\approx
\Normal\left(
\phi,
\frac{1-\phi^2}{n}
\right).
\]
If we instead fit an AR(2) model,
\[
Y_t-\phi_1Y_{t-1}-\phi_2Y_{t-2}=\varepsilon_t,
\]
then the true AR(1) process corresponds to
\[
\phi_1=\phi,
\qquad
\phi_2=0.
\]
The estimated coefficient $\widehat{\phi}_2$ should therefore be statistically insignificant.

The variance of $\widehat{\phi}_1$ in the overfitted AR(2) model is approximately
\[
\Var(\widehat{\phi}_1)
\approx
\frac{1}{n},
\]
whereas in the correct AR(1) model,
\[
\Var(\widehat{\phi})
\approx
\frac{1-\phi^2}{n}.
\]
Since
\[
\frac{1}{n}
\geq
\frac{1-\phi^2}{n},
\]
overfitting increases the standard error of the main coefficient.

For example, suppose $n=100$ and the true AR(1) coefficient is $\phi=0.7$. Under the correct AR(1) model,
\[
\widehat{\phi}
\approx
\Normal\left(0.7,\frac{1-0.7^2}{100}\right)
=
\Normal\left(0.7,\frac{0.51}{100}\right).
\]
An approximate $95\%$ interval is
\[
\widehat{\phi}\pm 2\sqrt{\frac{0.51}{100}}
=
\widehat{\phi}\pm 0.143.
\]
If we overfit AR(2), then
\[
\widehat{\phi}_1
\approx
\Normal\left(0.7,\frac{1}{100}\right),
\]
so an approximate $95\%$ interval is
\[
\widehat{\phi}_1\pm 2\sqrt{\frac{1}{100}}
=
\widehat{\phi}_1\pm 0.2.
\]
Thus the standard error increases when the model is unnecessarily large.

\section{Model Diagnostics}

So far, the time series workflow has included model specification and parameter estimation. After fitting a model, the next step is diagnostic checking.

\subsection{Residual Analysis}

Suppose we estimate model parameters using the observed data $Y_1,\ldots,Y_n$. Using the estimated parameters, we compute fitted values $\widehat{Y}_t$ and residuals
\[
\widehat{\varepsilon}_t
=
Y_t-\widehat{Y}_t.
\]

For an AR($p$) model, the fitted value is
\[
\widehat{Y}_t
=
\widehat{\mu}
+
\widehat{\phi}_1(Y_{t-1}-\widehat{\mu})
+\cdots+
\widehat{\phi}_p(Y_{t-p}-\widehat{\mu}),
\]
so the residual is
\[
\widehat{\varepsilon}_t
=
Y_t-\widehat{Y}_t.
\]

For an MA($q$) or ARMA($p,q$) model, the residuals are computed recursively from the fitted model. For example, an ARMA model has residuals of the form
\[
\widehat{\varepsilon}_t
=
Y_t
-
\widehat{\phi}_1Y_{t-1}
-\widehat{\phi}_2Y_{t-2}
-\cdots
+
\widehat{\theta}_1\widehat{\varepsilon}_{t-1}
+
\widehat{\theta}_2\widehat{\varepsilon}_{t-2}
+\cdots.
\]
The estimated coefficients are used in place of the unknown true parameters.

A good fitted model should produce residuals that behave approximately like white noise. In particular, residuals should be approximately independent or at least uncorrelated, have mean zero, have constant variance, and often be approximately normal when Gaussian inference is being used.

\subsection{Diagnostic Tools}

The first tool is the sample ACF of the residuals. Let
\[
\widehat{r}_k
=
\Corr(\widehat{\varepsilon}_t,\widehat{\varepsilon}_{t-k})
\]
denote the sample autocorrelation of the residuals. If the fitted model is adequate, then the residual autocorrelations should be close to zero.

\begin{remark}
The notation $\widehat{r}_k$ is used here for the sample ACF of residuals. This is different from $r_k$, which usually denotes the sample ACF of the original series.
\end{remark}

The second tool is a plot of standardized residuals:
\[
\frac{\widehat{\varepsilon}_t}{\widehat{\sigma}_\varepsilon}.
\]
This plot is used to check for outliers, changing variance, and remaining time-varying structure.

The third tool is a QQ plot or normality test, such as the Shapiro--Wilk test. These tools assess whether the residuals are approximately normal.

The fourth tool is the Ljung--Box test, which tests whether several residual autocorrelations are jointly zero.

\begin{definition}
The Ljung--Box statistic is
\[
Q_{\mathrm{LB}}
=
n(n+2)
\left(
\frac{\widehat{r}_1^2}{n-1}
+
\frac{\widehat{r}_2^2}{n-2}
+
\cdots+
\frac{\widehat{r}_K^2}{n-K}
\right),
\]
where $K$ is the number of lags being tested.
\end{definition}

The hypotheses are
\[
H_0:\mathrm{ACF}_k(\widehat{\varepsilon}_t)=0
\text{ for all } k>0
\]
against
\[
H_a:\mathrm{ACF}_k(\widehat{\varepsilon}_t)\neq 0
\text{ for some } k>0.
\]
Under the null hypothesis, the residuals are consistent with white noise. Large values of $Q_{\mathrm{LB}}$ provide evidence against the fitted model.

Under $H_0$, the Ljung--Box statistic is approximately chi-square distributed. If the fitted model has $p+q$ estimated ARMA parameters, then a common approximation is
\[
Q_{\mathrm{LB}}
\approx
\chi^2_{K-p-q}.
\]

\begin{itemize}
\item If the $p$-value is small, reject $H_0$. This suggests the residuals still contain serial dependence, so the model should be adjusted.
\item If the $p$-value is large, do not reject $H_0$. This suggests the residuals are approximately uncorrelated.
\end{itemize}

\subsection{Overfitting}

Overfitting can be used as a diagnostic tool. Suppose AR(2) is the correct model, but we fit AR(3). We can still compare AR(2) and AR(3). If AR(2) is adequate, then in the AR(3) fit:

\begin{enumerate}
\item the estimated third coefficient should not be statistically significant;
\item the estimated first two coefficients should be similar to those obtained from the AR(2) fit;
\item residual diagnostics should not improve substantially.
\end{enumerate}

Similarly, we can overfit ARMA(2,1) and compare it to AR(2).

\begin{remark}
Do not compare ARMA$(3,1)$ directly to AR$(2)$ as a simple overfitting diagnostic, because it changes both the AR and MA structure. More generally, fitting ARMA$(p+1,q+1)$ to an ARMA$(p,q)$ model creates redundancy.
\end{remark}

Indeed,
\[
\Phi(B)Y_t=\Theta(B)\varepsilon_t.
\]
For any constant $c$,
\[
(1-cB)\Phi(B)Y_t=(1-cB)\Theta(B)\varepsilon_t.
\]
This is an ARMA$(p+1,q+1)$ representation of the same process. Therefore, overfitting both the AR and MA orders simultaneously can create parameter redundancy or unidentifiability.

\chapter{Model Prediction}
\section{Forecasting}

If we observe
\[
Y_1,\ldots,Y_t,
\]
we want to predict future observations
\[
Y_{t+1},Y_{t+2},\ldots.
\]

The key object in forecasting is conditional expectation. Given observations $Y_1,\ldots,Y_t$, the optimal mean squared error predictor of $Y_{t+h}$ is
\[
\widehat{Y}_t(h)
=
\E[Y_{t+h}\mid Y_1,\ldots,Y_t].
\]

\begin{theorem}
For any random variable $X$, the constant $c$ minimizing
\[
\E[(X-c)^2]
\]
is
\[
c=\E[X].
\]
More generally, the random variable $g(Y)$ minimizing
\[
\E[(X-g(Y))^2]
\]
among all functions of $Y$ is
\[
g(Y)=\E[X\mid Y].
\]
\end{theorem}

\begin{proof}
For the unconditional case,
\[
\E[(X-c)^2]
=
\E[(X-\E[X]+\E[X]-c)^2].
\]
Expanding,
\[
\E[(X-c)^2]
=
\E[(X-\E[X])^2]
+
2(\E[X]-c)\E[X-\E[X]]
+
(\E[X]-c)^2.
\]
The middle term is zero, so
\[
\E[(X-c)^2]
=
\Var(X)+(\E[X]-c)^2.
\]
This is minimized when $c=\E[X]$.

The conditional version follows by applying the same argument conditionally on $Y$.
\end{proof}

\subsection{Basic Notation}

We use the following notation:

\begin{itemize}
\item $t$ is the forecast origin.
\item $h$ is the lead time.
\item $\widehat{Y}_t(h)$ is the prediction of $Y_{t+h}$ made at time $t$.
\item $e_t(h)$ is the forecast error:
\[
e_t(h)=Y_{t+h}-\widehat{Y}_t(h).
\]
\end{itemize}

If the process is normal, an approximate $(1-\alpha)100\%$ prediction interval for $Y_{t+h}$ is
\[
\widehat{Y}_t(h)
\pm
z_{\alpha/2}\sqrt{\Var(e_t(h))}.
\]

\subsection*{Conditional Expectation}

The forecast is
\[
\widehat{Y}_t(h)
=
\E[Y_{t+h}\mid Y_1,\ldots,Y_t].
\]
It satisfies
\[
\widehat{Y}_t(h)
=
\argmin_c
\E[(Y_{t+h}-c)^2\mid Y_1,\ldots,Y_t].
\]

Some useful properties of conditional expectation are:

\begin{itemize}
\item If $X$ is known given the conditioning information, then
\[
\E[g(X)\mid X]=g(X).
\]

\item If $X$ and $Y$ are independent, then
\[
\E[X\mid Y]=\E[X].
\]

\item If $a$ and $b$ are constants, then
\[
\E[aX+bY\mid \mathcal{F}]
=
a\E[X\mid \mathcal{F}]
+
b\E[Y\mid \mathcal{F}].
\]
\end{itemize}

\subsection{Example: Trend Plus Noise}

Consider
\[
Y_t=\mu_t+X_t,
\qquad
X_t\iid \Normal(0,\sigma^2),
\]
where $\mu_t$ is deterministic. Then
\begin{align*}
\widehat{Y}_t(h)
&=
\E[Y_{t+h}\mid Y_1,\ldots,Y_t] \\
&=
\E[\mu_{t+h}+X_{t+h}\mid Y_1,\ldots,Y_t] \\
&=
\mu_{t+h}
+
\E[X_{t+h}\mid Y_1,\ldots,Y_t].
\end{align*}
Since $X_{t+h}$ is independent of the past and has mean zero,
\[
\widehat{Y}_t(h)=\mu_{t+h}.
\]
The forecast error is
\[
e_t(h)=Y_{t+h}-\widehat{Y}_t(h)=X_{t+h}.
\]
Therefore,
\[
\E[e_t(h)]=0,
\qquad
\Var(e_t(h))=\sigma^2.
\]
If the noise is normal, an approximate $95\%$ prediction interval is
\[
\widehat{Y}_t(h)\pm 2\sigma
=
\mu_{t+h}\pm 2\sigma.
\]
The interval width does not depend on $h$ because the noise terms are IID.

\subsection*{Example: AR(1)}

Consider a causal stationary AR(1) process with mean $\mu$:
\[
Y_t-\mu=\phi(Y_{t-1}-\mu)+\varepsilon_t,
\qquad
\varepsilon_t\iid \Normal(0,\sigma_\varepsilon^2).
\]
For one-step forecasting,
\begin{align*}
\widehat{Y}_t(1)
&=
\E[Y_{t+1}\mid Y_1,\ldots,Y_t] \\
&=
\E[\mu+\phi(Y_t-\mu)+\varepsilon_{t+1}\mid Y_1,\ldots,Y_t] \\
&=
\mu+\phi(Y_t-\mu).
\end{align*}
For $h\geq 2$,
\[
Y_{t+h}
=
\mu+\phi(Y_{t+h-1}-\mu)+\varepsilon_{t+h}.
\]
Taking conditional expectations gives
\[
\widehat{Y}_t(h)
=
\mu+\phi(\widehat{Y}_t(h-1)-\mu).
\]
Thus
\[
\widehat{Y}_t(h)-\mu
=
\phi(\widehat{Y}_t(h-1)-\mu).
\]
Iterating,
\[
\widehat{Y}_t(h)
=
\mu+\phi^h(Y_t-\mu).
\]
As $h\to\infty$,
\[
\widehat{Y}_t(h)\to \mu
\]
because $|\phi|<1$.

Now derive the forecast error. Using the GLP representation,
\[
Y_t=\mu+\sum_{j=0}^{\infty}\phi^j\varepsilon_{t-j}.
\]
Then
\[
Y_{t+h}
=
\mu+\sum_{j=0}^{\infty}\phi^j\varepsilon_{t+h-j}.
\]
The terms involving future innovations are
\[
\varepsilon_{t+h},
\phi\varepsilon_{t+h-1},
\ldots,
\phi^{h-1}\varepsilon_{t+1}.
\]
The remaining terms are functions of the information available at time $t$. Therefore,
\[
e_t(h)
=
Y_{t+h}-\widehat{Y}_t(h)
=
\sum_{j=0}^{h-1}\phi^j\varepsilon_{t+h-j}.
\]
Hence
\[
\E[e_t(h)]=0,
\]
and
\[
\Var(e_t(h))
=
\sigma_\varepsilon^2
\sum_{j=0}^{h-1}\phi^{2j}
=
\sigma_\varepsilon^2
\frac{1-\phi^{2h}}{1-\phi^2}.
\]
As $h\to\infty$,
\[
\Var(e_t(h))
\to
\frac{\sigma_\varepsilon^2}{1-\phi^2}
=
\gamma_0.
\]
If the process is normal, an approximate $95\%$ prediction interval is
\[
\widehat{Y}_t(h)
\pm
2\sigma_\varepsilon
\sqrt{
\sum_{j=0}^{h-1}\phi^{2j}
}.
\]

\subsection*{Example: MA(1)}

Consider an invertible MA(1) process with mean $\mu$:
\[
Y_t=\mu+\varepsilon_t-\theta\varepsilon_{t-1}.
\]
For one-step forecasting,
\[
\widehat{Y}_t(1)
=
\E[Y_{t+1}\mid Y_1,\ldots,Y_t]
=
\E[\mu+\varepsilon_{t+1}-\theta\varepsilon_t\mid Y_1,\ldots,Y_t].
\]
Under invertibility, $\varepsilon_t$ can be recovered from the observed past. Therefore,
\[
\widehat{Y}_t(1)=\mu-\theta\varepsilon_t.
\]
The one-step forecast error is
\[
e_t(1)
=
Y_{t+1}-\widehat{Y}_t(1)
=
\varepsilon_{t+1}.
\]
Thus
\[
\Var(e_t(1))=\sigma_\varepsilon^2.
\]

For $h\geq 2$,
\[
Y_{t+h}=\mu+\varepsilon_{t+h}-\theta\varepsilon_{t+h-1}.
\]
Both innovations on the right-hand side are future innovations relative to time $t$, so their conditional expectations are zero. Hence
\[
\widehat{Y}_t(h)=\mu,
\qquad h\geq 2.
\]

\begin{exercise}
Verify that $\widehat{Y}_t(h)=\mu$ for $h\geq 2$.
\end{exercise}

\begin{proof}
For $h\geq 2$,
\[
Y_{t+h}
=
\mu+\varepsilon_{t+h}-\theta\varepsilon_{t+h-1}.
\]
Both $\varepsilon_{t+h}$ and $\varepsilon_{t+h-1}$ occur after time $t$, so they are independent of $Y_1,\ldots,Y_t$ and have mean zero. Therefore,
\[
\E[\varepsilon_{t+h}\mid Y_1,\ldots,Y_t]=0,
\]
and
\[
\E[\varepsilon_{t+h-1}\mid Y_1,\ldots,Y_t]=0.
\]
Thus
\[
\widehat{Y}_t(h)
=
\E[Y_{t+h}\mid Y_1,\ldots,Y_t]
=
\mu.
\]
\end{proof}

For $h\geq 2$, the forecast error is
\[
e_t(h)
=
Y_{t+h}-\mu
=
\varepsilon_{t+h}-\theta\varepsilon_{t+h-1}.
\]
Thus
\[
\Var(e_t(h))
=
(1+\theta^2)\sigma_\varepsilon^2
=
\gamma_0.
\]

\subsection*{Example: Random Walk with Drift}

Consider
\[
Y_t=Y_{t-1}+\theta_0+\varepsilon_t.
\]
For one-step forecasting,
\[
\widehat{Y}_t(1)
=
\E[Y_{t+1}\mid Y_1,\ldots,Y_t]
=
Y_t+\theta_0.
\]
In general,
\[
\widehat{Y}_t(h)=Y_t+h\theta_0.
\]

\begin{exercise}
Show that
\[
\widehat{Y}_t(h)=Y_t+h\theta_0.
\]
\end{exercise}

\begin{proof}
Iterating the model forward,
\[
Y_{t+h}
=
Y_t+h\theta_0+\sum_{j=1}^{h}\varepsilon_{t+j}.
\]
Taking conditional expectation given $Y_1,\ldots,Y_t$,
\[
\E[Y_{t+h}\mid Y_1,\ldots,Y_t]
=
Y_t+h\theta_0
+
\sum_{j=1}^{h}
\E[\varepsilon_{t+j}\mid Y_1,\ldots,Y_t].
\]
The future innovations have conditional mean zero, so
\[
\widehat{Y}_t(h)=Y_t+h\theta_0.
\]
\end{proof}

The forecast error is
\[
e_t(h)
=
\sum_{j=1}^{h}\varepsilon_{t+j}.
\]
Therefore,
\[
\E[e_t(h)]=0,
\qquad
\Var(e_t(h))=h\sigma_\varepsilon^2.
\]
Thus the forecast variance grows without bound as $h\to\infty$.

\subsection{Example: ARMA(1,1) with Mean}

Consider an invertible ARMA(1,1) model
\[
Y_t=\phi Y_{t-1}+\varepsilon_t-\theta\varepsilon_{t-1}+\theta_0.
\]
If the process is stationary, its mean is
\[
\mu=\frac{\theta_0}{1-\phi}.
\]
For one-step forecasting,
\begin{align*}
\widehat{Y}_t(1)
&=
\E[\phi Y_t+\varepsilon_{t+1}-\theta\varepsilon_t+\theta_0\mid Y_1,\ldots,Y_t] \\
&=
\phi Y_t-\theta\varepsilon_t+\theta_0.
\end{align*}

For $h\geq 2$, future innovations have conditional mean zero, and the recursion gives
\[
\widehat{Y}_t(h)
=
\phi\widehat{Y}_t(h-1)+\theta_0.
\]
Solving this recursion,
\[
\widehat{Y}_t(h)
=
\phi^hY_t-\phi^{h-1}\theta\varepsilon_t
+
\frac{1-\phi^h}{1-\phi}\theta_0.
\]
As $h\to\infty$, if $|\phi|<1$,
\[
\widehat{Y}_t(h)\to \frac{\theta_0}{1-\phi}=\mu.
\]

\section*{General Forecasting Observations}

The previous examples illustrate common behavior:

\begin{itemize}
\item For nonstationary random walks, forecast variance grows without bound.
\item For stationary ARMA processes, forecasts converge to the mean.
\item For causal stationary ARMA processes, forecast error variance increases with $h$ but converges to the process variance $\gamma_0$.
\item For an MA($q$) process, forecasts equal the mean after more than $q$ steps.
\end{itemize}

\section{Prediction for ARMA($p,q$)}

Assume the ARMA($p,q$) model is causal and invertible. The model is
\[
Y_{t+h}
=
\phi_1Y_{t+h-1}
+\cdots+
\phi_pY_{t+h-p}
+
\varepsilon_{t+h}
-\theta_1\varepsilon_{t+h-1}
-\cdots-
\theta_q\varepsilon_{t+h-q}.
\]
Taking conditional expectations given $Y_1,\ldots,Y_t$ gives
\[
\widehat{Y}_t(h)
=
\phi_1\widehat{Y}_t(h-1)
+\cdots+
\phi_p\widehat{Y}_t(h-p)
+
\E[\varepsilon_{t+h}-\theta_1\varepsilon_{t+h-1}-\cdots-\theta_q\varepsilon_{t+h-q}\mid Y_1,\ldots,Y_t].
\]
For $h>q$, all innovation terms in the conditional expectation are future innovations, so
\[
\E[\varepsilon_{t+h}-\theta_1\varepsilon_{t+h-1}-\cdots-\theta_q\varepsilon_{t+h-q}\mid Y_1,\ldots,Y_t]=0.
\]
Thus, for $h>q$,
\[
\widehat{Y}_t(h)
=
\phi_1\widehat{Y}_t(h-1)
+\cdots+
\phi_p\widehat{Y}_t(h-p).
\]
Therefore, long-horizon ARMA forecasts satisfy the same recursion as the AR part.

If the AR roots are distinct, then for large $h$,
\[
\widehat{Y}_t(h)
=
C_1\left(\frac{1}{z_1}\right)^h
+\cdots+
C_p\left(\frac{1}{z_p}\right)^h.
\]
Under causality, $|z_i|>1$, so
\[
\widehat{Y}_t(h)\to 0
\]
for a mean-zero process. If the process has mean $\mu$, then
\[
\widehat{Y}_t(h)\to \mu.
\]

\subsection{Forecast Error for ARMA($p,q$)}

Using the causal GLP representation,
\[
Y_t=\sum_{j=0}^{\infty}\psi_j\varepsilon_{t-j}.
\]
Then
\[
Y_{t+h}
=
\sum_{j=0}^{\infty}\psi_j\varepsilon_{t+h-j}.
\]
Separate this into future and past parts:
\[
Y_{t+h}
=
\underbrace{\sum_{j=0}^{h-1}\psi_j\varepsilon_{t+h-j}}_{I_t(h)}
+
\underbrace{\sum_{j=h}^{\infty}\psi_j\varepsilon_{t+h-j}}_{C_t(h)}.
\]
The first part contains future innovations relative to time $t$, while the second part contains innovations dated at or before time $t$.

Since the process is invertible, the past innovations are functions of $Y_1,\ldots,Y_t$. Therefore,
\[
\widehat{Y}_t(h)
=
\E[Y_{t+h}\mid Y_1,\ldots,Y_t]
=
C_t(h).
\]
The forecast error is
\[
e_t(h)
=
Y_{t+h}-\widehat{Y}_t(h)
=
I_t(h)
=
\sum_{j=0}^{h-1}\psi_j\varepsilon_{t+h-j}.
\]
Therefore,
\[
\E[e_t(h)]=0,
\]
and
\[
\Var(e_t(h))
=
\sigma_\varepsilon^2
\sum_{j=0}^{h-1}\psi_j^2.
\]
Since the process is stationary,
\[
\sum_{j=0}^{\infty}\psi_j^2<\infty.
\]
Thus,
\[
\Var(e_t(h))
\to
\sigma_\varepsilon^2
\sum_{j=0}^{\infty}\psi_j^2
=
\gamma_0.
\]

\section{Forecasting ARIMA($p,d,q$)}

For ARIMA($p,d,q$),
\[
\Phi(B)(1-B)^dY_t=\Theta(B)\varepsilon_t.
\]
Equivalently,
\[
\widetilde{\Phi}(B)Y_t=\Theta(B)\varepsilon_t,
\]
where
\[
\widetilde{\Phi}(z)=\Phi(z)(1-z)^d.
\]
The modified AR polynomial has $d$ unit roots. These unit roots cause nonstationarity.

\begin{example}
Consider ARIMA$(1,1,1)$:
\[
Y_t-Y_{t-1}=W_t,
\]
where
\[
W_t=\phi W_{t-1}+\varepsilon_t-\theta\varepsilon_{t-1}.
\]
Then
\[
Y_t=(1+\phi)Y_{t-1}-\phi Y_{t-2}+\varepsilon_t-\theta\varepsilon_{t-1}.
\]
The forecast recursion is
\[
\widehat{Y}_t(1)
=
(1+\phi)Y_t-\phi Y_{t-1}-\theta\varepsilon_t,
\]
and for $h\geq 2$,
\[
\widehat{Y}_t(h)
=
(1+\phi)\widehat{Y}_t(h-1)-\phi\widehat{Y}_t(h-2).
\]
\end{example}

For ARIMA models with $d\geq 1$, the forecast error variance usually diverges as $h\to\infty$. In a GLP-like representation,
\[
e_t(h)
=
\sum_{j=0}^{h-1}\psi_j\varepsilon_{t+h-j},
\]
but the coefficients typically do not have a square-summable infinite limit. Thus
\[
\Var(e_t(h))
=
\sigma_\varepsilon^2\sum_{j=0}^{h-1}\psi_j^2
\to \infty
\]
for integrated processes.

\section{Prediction Intervals}

For either GLP or GLP-like forecasting representations, the forecast error often has the form
\[
e_t(h)
=
\sum_{j=0}^{h-1}\psi_j\varepsilon_{t+h-j}.
\]
If
\[
\varepsilon_t\iid \Normal(0,\sigma_\varepsilon^2),
\]
then
\[
e_t(h)
\sim
\Normal\left(
0,
\sigma_\varepsilon^2\sum_{j=0}^{h-1}\psi_j^2
\right).
\]
Therefore,
\[
\Prob\left(
-z_{\alpha/2}
\leq
\frac{e_t(h)}
{\sigma_\varepsilon\sqrt{\sum_{j=0}^{h-1}\psi_j^2}}
\leq
z_{\alpha/2}
\right)
=
1-\alpha.
\]
Since
\[
e_t(h)=Y_{t+h}-\widehat{Y}_t(h),
\]
we obtain
\[
\Prob\left(
\widehat{Y}_t(h)
-
z_{\alpha/2}\sqrt{\Var(e_t(h))}
\leq
Y_{t+h}
\leq
\widehat{Y}_t(h)
+
z_{\alpha/2}\sqrt{\Var(e_t(h))}
\right)
=
1-\alpha.
\]
Thus a $(1-\alpha)100\%$ prediction interval is
\[
\widehat{Y}_t(h)
\pm
z_{\alpha/2}\sqrt{\Var(e_t(h))}.
\]

\begin{remark}
This is a prediction interval for the future random variable $Y_{t+h}$. It is not a confidence interval for a fixed unknown parameter.
\end{remark}

\section{Exponentially Weighted Moving Average}

EWMA stands for exponentially weighted moving average. It is a simple forecasting method that is useful when a quick forecast is needed.

The EWMA forecast is defined recursively by
\[
\widehat{Y}_t(1)
=
(1-\theta)Y_t+\theta\widehat{Y}_{t-1}(1).
\]
Thus the new forecast is a weighted average of the current observed value $Y_t$ and the previous forecast $\widehat{Y}_{t-1}(1)$.

Equivalently,
\[
\widehat{Y}_t(1)
=
\widehat{Y}_{t-1}(1)
+
(1-\theta)\left(Y_t-\widehat{Y}_{t-1}(1)\right).
\]
This says that the forecast is updated by adding a fraction of the most recent forecast error.

\subsection{EWMA from an ARIMA$(0,1,1)$ Model}

Consider the ARIMA$(0,1,1)$ model
\[
Y_t-Y_{t-1}=W_t,
\qquad
W_t=\varepsilon_t-\theta\varepsilon_{t-1}.
\]
Equivalently,
\[
Y_t=Y_{t-1}+\varepsilon_t-\theta\varepsilon_{t-1}.
\]
For an invertible model, the innovations have an invertible representation in terms of past observations:
\[
\varepsilon_t
=
\pi_0Y_t+\pi_1Y_{t-1}+\pi_2Y_{t-2}+\cdots.
\]
Then
\[
Y_{t+1}
=
Y_t+\varepsilon_{t+1}-\theta\varepsilon_t.
\]
Taking conditional expectation given the observed past,
\[
\widehat{Y}_t(1)
=
\E[Y_{t+1}\mid Y_t,Y_{t-1},\ldots]
=
Y_t-\theta\varepsilon_t.
\]
Substituting the invertible representation gives
\[
\widehat{Y}_t(1)
=
Y_t-\theta\sum_{j=0}^{\infty}\pi_jY_{t-j}.
\]

For ARIMA$(0,1,1)$, one can show that
\[
\pi_0=1,
\qquad
\pi_j=(\theta-1)\theta^{j-1},
\quad j\geq 1.
\]
Therefore,
\begin{align*}
\widehat{Y}_t(1)
&=
Y_t-\theta Y_t-\theta\sum_{j=1}^{\infty}(\theta-1)\theta^{j-1}Y_{t-j} \\
&=
(1-\theta)Y_t+\theta(1-\theta)\sum_{j=1}^{\infty}\theta^{j-1}Y_{t-j} \\
&=
(1-\theta)Y_t+\theta\left[(1-\theta)Y_{t-1}+\theta(1-\theta)Y_{t-2}+\cdots\right].
\end{align*}
But the expression in brackets is the previous EWMA forecast $\widehat{Y}_{t-1}(1)$. Hence
\[
\widehat{Y}_t(1)
=
(1-\theta)Y_t+\theta\widehat{Y}_{t-1}(1).
\]

\begin{remark}
The smoothing parameter is often chosen ad hoc in simple EWMA forecasting, but the ARIMA$(0,1,1)$ derivation gives it a model-based interpretation.
\end{remark}

\section{Seasonal ARMA Models}

Seasonal ARMA models extend ARMA models by allowing dependence at seasonal lags. For monthly data, the seasonal period is often $s=12$; for quarterly data, it is often $s=4$.

\subsection{Seasonal MA}

Consider the seasonal MA model
\[
Y_t=\varepsilon_t-\Theta\varepsilon_{t-12}.
\]
This is a seasonal MA model of order $1$ with seasonal period $12$, denoted by
\[
\mathrm{MA}(1)_{12}.
\]

The mean is
\[
\E[Y_t]=0.
\]
The variance is
\[
\gamma_0
=
\Var(Y_t)
=
(1+\Theta^2)\sigma^2.
\]
The lag-$12$ autocovariance is
\[
\gamma_{12}
=
\Cov(Y_t,Y_{t-12})
=
\Cov(\varepsilon_t-\Theta\varepsilon_{t-12},
\varepsilon_{t-12}-\Theta\varepsilon_{t-24})
=
-\Theta\sigma^2.
\]
All other autocovariances are zero except by symmetry:
\[
\gamma_k=0
\qquad
\text{if } k\neq 0,\pm 12.
\]
Therefore,
\[
\rho_{12}
=
\frac{-\Theta}{1+\Theta^2},
\qquad
\rho_k=0
\quad
\text{if } k\neq 0,\pm 12.
\]

\begin{exercise}
Verify the ACVF and ACF above.
\end{exercise}

\begin{proof}
Since
\[
Y_t=\varepsilon_t-\Theta\varepsilon_{t-12},
\]
we have
\[
\Var(Y_t)
=
\Var(\varepsilon_t)+\Theta^2\Var(\varepsilon_{t-12})
=
(1+\Theta^2)\sigma^2.
\]
For lag $12$,
\begin{align*}
\Cov(Y_t,Y_{t-12})
&=
\Cov(\varepsilon_t-\Theta\varepsilon_{t-12},
\varepsilon_{t-12}-\Theta\varepsilon_{t-24}) \\
&=
-\Theta\Cov(\varepsilon_{t-12},\varepsilon_{t-12}) \\
&=
-\Theta\sigma^2.
\end{align*}
For all other nonzero lags, the innovation terms do not overlap, so the covariance is zero. Dividing by $\gamma_0=(1+\Theta^2)\sigma^2$ gives the ACF.
\end{proof}

More generally, a seasonal MA$(Q)$ model with seasonal period $s$ is
\[
Y_t
=
\varepsilon_t-\Theta_1\varepsilon_{t-s}
-\Theta_2\varepsilon_{t-2s}
-\cdots
-\Theta_Q\varepsilon_{t-Qs}.
\]
Its MA polynomial is
\[
\Theta(B^s)
=
1-\Theta_1B^s-\Theta_2B^{2s}-\cdots-\Theta_QB^{Qs}.
\]

\subsection{Seasonal AR}

Consider the seasonal AR model
\[
Y_t=\Phi Y_{t-12}+\varepsilon_t.
\]
This is denoted by
\[
\mathrm{AR}(1)_{12}.
\]
The causality condition is
\[
|\Phi|<1.
\]
The autocorrelation function satisfies
\[
\rho_{12k}=\Phi^k,
\qquad k\geq 0,
\]
and
\[
\rho_m=0
\qquad
\text{if } m \neq 12k.
\]

\begin{exercise}
Verify the ACF above.
\end{exercise}

\begin{proof}
The model separates into $12$ independent AR(1)-type subsequences. For each fixed month $r\in\{1,\ldots,12\}$, the subsequence
\[
Y_r,Y_{r+12},Y_{r+24},\ldots
\]
satisfies an AR(1) recursion with coefficient $\Phi$. Therefore, correlations at lags that are multiples of $12$ are
\[
\rho_{12k}=\Phi^k.
\]
If a lag is not a multiple of $12$, then the two observations belong to different seasonal subsequences and share no innovations, so their covariance is zero.
\end{proof}

More generally, a seasonal AR$(P)$ model with seasonal period $s$ is
\[
Y_t=\Phi_1Y_{t-s}+\Phi_2Y_{t-2s}+\cdots+\Phi_PY_{t-Ps}+\varepsilon_t.
\]
Its AR polynomial is
\[
\Phi(B^s)
=
1-\Phi_1B^s-\Phi_2B^{2s}-\cdots-\Phi_PB^{Ps}.
\]

\section{Multiplicative Seasonal ARMA Models}

A multiplicative seasonal ARMA model combines a nonseasonal ARMA$(p,q)$ model and a seasonal ARMA$(P,Q)$ model by multiplying the corresponding AR and MA polynomials.

The notation is
\[
\mathrm{ARMA}(p,q)\times(P,Q)_s.
\]
The AR polynomial is
\[
\phi(B)\Phi(B^s),
\]
and the MA polynomial is
\[
\theta(B)\Theta(B^s),
\]
where
\[
\phi(B)
=
1-\phi_1B-\cdots-\phi_pB^p,
\]
\[
\Phi(B^s)
=
1-\Phi_1B^s-\cdots-\Phi_PB^{Ps},
\]
\[
\theta(B)
=
1-\theta_1B-\cdots-\theta_qB^q,
\]
and
\[
\Theta(B^s)
=
1-\Theta_1B^s-\cdots-\Theta_QB^{Qs}.
\]

\begin{example}
Consider
\[
\mathrm{ARMA}(0,1)\times(1,0)_{12}.
\]
The AR polynomial is
\[
1-\Phi B^{12},
\]
and the MA polynomial is
\[
1-\theta B.
\]
Thus
\[
(1-\Phi B^{12})Y_t=(1-\theta B)\varepsilon_t,
\]
or
\[
Y_t-\Phi Y_{t-12}
=
\varepsilon_t-\theta\varepsilon_{t-1}.
\]
\end{example}

For this model, assuming causality, the autocovariances satisfy
\[
\gamma_0
=
\frac{1+\theta^2}{1-\Phi^2}\sigma^2,
\]
\[
\gamma_1
=
\frac{-\theta}{1-\Phi^2}\sigma^2,
\]
and
\[
\gamma_{12}
=
\Phi\gamma_0.
\]
More generally,
\[
\gamma_{12k}
=
\Phi^k\gamma_0,
\qquad
\gamma_{12k+1}
=
\Phi^k\gamma_1,
\qquad
\gamma_m=0
\text{ otherwise}.
\]

\section{Multiplicative Seasonal ARIMA Models}

A multiplicative seasonal ARIMA model combines a nonseasonal ARIMA$(p,d,q)$ model with a seasonal ARIMA$(P,D,Q)_s$ model.

The notation is
\[
\mathrm{ARIMA}(p,d,q)\times(P,D,Q)_s.
\]
The ordinary differencing operator is
\[
\nabla Y_t
=
(1-B)Y_t
=
Y_t-Y_{t-1}.
\]
The seasonal differencing operator with period $s$ is
\[
\nabla_sY_t
=
(1-B^s)Y_t
=
Y_t-Y_{t-s}.
\]
Thus
\[
\nabla^d\nabla_s^DY_t
=
(1-B)^d(1-B^s)^DY_t.
\]

The full model is
\[
\phi(B)\Phi(B^s)(1-B)^d(1-B^s)^DY_t
=
\theta(B)\Theta(B^s)\varepsilon_t.
\]

\begin{example}
Consider
\[
Y_t=0.5Y_{t-1}+Y_{t-4}-0.5Y_{t-5}+\varepsilon_t-0.3\varepsilon_{t-1}.
\]
Rearranging,
\[
(1-0.5B-B^4+0.5B^5)Y_t
=
(1-0.3B)\varepsilon_t.
\]
Factor the left-hand side:
\[
1-0.5B-B^4+0.5B^5
=
(1-0.5B)(1-B^4).
\]
Therefore, this is an
\[
\mathrm{ARIMA}(1,0,1)\times(0,1,0)_4
\]
model.
\end{example}

\begin{example}
Consider
\[
Y_t=Y_{t-4}+\varepsilon_t-\theta_1\varepsilon_{t-1}-\theta_2\varepsilon_{t-2}.
\]
Then
\[
(1-B^4)Y_t=(1-\theta_1B-\theta_2B^2)\varepsilon_t.
\]
This is an
\[
\mathrm{ARIMA}(0,0,2)\times(0,1,0)_4
\]
model.
\end{example}

\begin{example}
Consider
\[
Y_t=Y_{t-1}+Y_{t-12}-Y_{t-13}
+\varepsilon_t-0.1\varepsilon_{t-1}-0.1\varepsilon_{t-12}
+0.01\varepsilon_{t-13}.
\]
The AR side factors as
\[
1-B-B^{12}+B^{13}
=
(1-B)(1-B^{12}),
\]
and the MA side factors as
\[
1-0.1B-0.1B^{12}+0.01B^{13}
=
(1-0.1B)(1-0.1B^{12}).
\]
Therefore, the model is
\[
\mathrm{ARIMA}(0,1,1)\times(0,1,1)_{12}.
\]
\end{example}

\section{Cross-Covariance and Cross-Correlation}

Previously, we used past values of a single time series $\{Y_t\}$ to predict future values of that same series. Now suppose we have two time series $\{X_t\}$ and $\{Y_t\}$, and we want to use past values of $\{X_t\}$ to help predict $\{Y_t\}$.

\subsection{Cross-Covariance Function}

\begin{definition}
The cross-covariance function between $\{X_t\}$ and $\{Y_t\}$ is
\[
\gamma_k(X,Y)
=
\Cov(X_t,Y_{t-k}).
\]
\end{definition}

The vector process $(X_t,Y_t)$ is jointly weakly stationary if:

\begin{enumerate}
\item $\E[X_t]=\mu_X$ and $\E[Y_t]=\mu_Y$ for all $t$.
\item $\Var(X_t)$ and $\Var(Y_t)$ are constant over time.
\item $\Cov(X_t,X_{t-k})$ and $\Cov(Y_t,Y_{t-k})$ depend only on $k$.
\item $\Cov(X_t,Y_{t-k})$ depends only on $k$.
\end{enumerate}

If joint stationarity holds, then
\[
\gamma_0(X,Y)=\Cov(X_t,Y_t),
\]
\[
\gamma_1(X,Y)=\Cov(X_t,Y_{t-1}),
\]
and
\[
\gamma_{-1}(X,Y)=\Cov(X_t,Y_{t+1}).
\]
Unlike the autocovariance function of a single series, the cross-covariance function is not necessarily symmetric:
\[
\gamma_k(X,Y)
\neq
\gamma_{-k}(X,Y)
\]
in general. Instead,
\[
\gamma_k(X,Y)=\gamma_{-k}(Y,X).
\]

\subsection{Cross-Correlation Function}

\begin{definition}
The cross-correlation function between two jointly stationary series is
\[
\rho_k(X,Y)
=
\frac{\gamma_k(X,Y)}
{\sqrt{\gamma_0(X)\gamma_0(Y)}}.
\]
\end{definition}

\begin{example}
Suppose
\[
X_t\iid \Normal(0,\sigma_X^2)
\]
and
\[
Y_t=\beta_0+\beta_1X_{t-d}+\varepsilon_t,
\qquad
\varepsilon_t\iid \Normal(0,\sigma_\varepsilon^2),
\]
where $\{X_t\}$ and $\{\varepsilon_t\}$ are independent. Then
\[
\gamma_k(X,Y)
=
\Cov(X_t,Y_{t-k})
=
\Cov(X_t,\beta_0+\beta_1X_{t-k-d}+\varepsilon_{t-k}).
\]
Only the term involving $X_t$ can produce nonzero covariance. Thus
\[
\gamma_k(X,Y)
=
\begin{cases}
\beta_1\sigma_X^2, & k=-d,\\
0, & k\neq -d.
\end{cases}
\]
Since
\[
\gamma_0(X)=\sigma_X^2,
\qquad
\gamma_0(Y)=\beta_1^2\sigma_X^2+\sigma_\varepsilon^2,
\]
we get
\[
\rho_k(X,Y)
=
\begin{cases}
\dfrac{\beta_1\sigma_X}{\sqrt{\beta_1^2\sigma_X^2+\sigma_\varepsilon^2}}, & k=-d,\\
0, & k\neq -d.
\end{cases}
\]
\end{example}

\begin{exercise}
Verify the cross-correlation function in the previous example.
\end{exercise}

\begin{proof}
The cross-covariance calculation gives
\[
\gamma_{-d}(X,Y)=\beta_1\sigma_X^2
\]
and all other cross-covariances are zero. Also,
\[
\Var(X_t)=\sigma_X^2,
\]
and by independence,
\[
\Var(Y_t)=\Var(\beta_1X_{t-d}+\varepsilon_t)
=
\beta_1^2\sigma_X^2+\sigma_\varepsilon^2.
\]
Therefore,
\[
\rho_{-d}(X,Y)
=
\frac{\beta_1\sigma_X^2}
{\sqrt{\sigma_X^2(\beta_1^2\sigma_X^2+\sigma_\varepsilon^2)}}
=
\frac{\beta_1\sigma_X}
{\sqrt{\beta_1^2\sigma_X^2+\sigma_\varepsilon^2}}.
\]
\end{proof}

\subsection{Bartlett's Theorem for Sample CCF}

Let $r_m(X,Y)$ denote the sample cross-correlation at lag $m$. For large $n$,
\[
r_m(X,Y)
\approx
\Normal\left(
\rho_m(X,Y),
\frac{1}{n}
\left[
1+2\sum_{k=1}^{\infty}\rho_k(X)\rho_k(Y)
\right]
\right).
\]

This may lead to spurious correlation. Even when the theoretical cross-correlation $\rho_m(X,Y)$ is zero, the sample cross-correlation can appear large if the variance term is large.

\begin{example}
Suppose
\[
X_t\sim \mathrm{AR}(1),
\qquad
Y_t\sim \mathrm{AR}(1),
\]
and the two processes are independent. Let their AR coefficients be $\phi_X$ and $\phi_Y$. Then the theoretical CCF is
\[
\rho_m(X,Y)=0
\]
for all $m$. However,
\[
\rho_k(X)=\phi_X^k,
\qquad
\rho_k(Y)=\phi_Y^k.
\]
Thus
\[
\Var(r_m(X,Y))
\approx
\frac{1}{n}
\left[
1+2\sum_{k=1}^{\infty}(\phi_X\phi_Y)^k
\right]
=
\frac{1}{n}
\left[
1+\frac{2\phi_X\phi_Y}{1-\phi_X\phi_Y}
\right].
\]
Equivalently,
\[
\Var(r_m(X,Y))
\approx
\frac{1+\phi_X\phi_Y}{n(1-\phi_X\phi_Y)}.
\]
If $\phi_X=\phi_Y=1/2$, then
\[
\Var(r_m(X,Y))
\approx
\frac{1+1/4}{n(1-1/4)}
=
\frac{5}{3n}.
\]
This is larger than the white-noise benchmark $1/n$, so the standard $\pm 2/\sqrt{n}$ rule is not reliable.
\end{example}

\section{Spurious Correlation and Prewhitening}

\subsection{Spurious Correlation}

The sampling distribution of the sample CCF can have variance larger than $1/n$:
\[
r_m(X,Y)
\approx
\Normal\left(
\rho_m(X,Y),
\frac{1}{n}
\left[
1+2\sum_{k=1}^{\infty}\rho_k(X)\rho_k(Y)
\right]
\right).
\]
If the variance is much larger than $1/n$, then the usual bounds
\[
\pm \frac{2}{\sqrt{n}}
\]
can falsely suggest cross-correlation. This is called spurious correlation.

If at least one of the two series is white noise, then the extra term disappears. For example, if
\[
\rho_k(X)=0
\qquad
\text{for all } k\geq 1,
\]
then
\[
\frac{1}{n}
\left[
1+2\sum_{k=1}^{\infty}\rho_k(X)\rho_k(Y)
\right]
=
\frac{1}{n}.
\]

\subsection{Prewhitening}

Suppose $\{X_t\}$ follows an ARMA$(p,q)$ model:
\[
\Phi(B)X_t=\Theta(B)\varepsilon_t.
\]
If the MA part is invertible, then
\[
\varepsilon_t
=
\Theta(B)^{-1}\Phi(B)X_t
=
\Pi(B)X_t,
\]
where
\[
\Pi(B)=\sum_{j=0}^{\infty}\pi_jB^j
\]
is called a prewhitening filter.

The idea of prewhitening is to apply $\Pi(B)$ to $X_t$ so that
\[
\Pi(B)X_t=\varepsilon_t,
\]
which is approximately white noise. Then we apply the same filter to $Y_t$:
\[
\Pi(B)Y_t=\widetilde{Y}_t.
\]
Now the dependence structure inside $X_t$ has been removed, and we can examine the CCF between $\varepsilon_t$ and $\widetilde{Y}_t$.

In practice, the procedure is:

\begin{enumerate}
\item Make $X_t$ stationary, if necessary.
\item Fit an ARMA model to $X_t$.
\item Use the fitted model to construct an approximate prewhitening filter $\widehat{\Pi}(B)$.
\item Compute
\[
\widehat{\Pi}(B)X_t=\widehat{\varepsilon}_t
\]
and
\[
\widehat{\Pi}(B)Y_t=\widetilde{Y}_t.
\]
\item Estimate the CCF between $\widehat{\varepsilon}_t$ and $\widetilde{Y}_t$.
\end{enumerate}

\begin{example}
Consider the model
\[
Y_t
=
\sum_{k=-\infty}^{\infty}\beta_kX_{t-k}+Z_t,
\]
where $\{Z_t\}$ is white noise independent of $\{X_t\}$. Suppose $\{X_t\}$ is not white noise. Since $X_t$ is serially correlated, the raw CCF between $X_t$ and $Y_t$ can be misleading.

Let $\Pi(B)$ be a prewhitening filter for $X_t$. Then
\[
\Pi(B)X_t=\widetilde{X}_t,
\]
where $\widetilde{X}_t$ is white noise. Applying the same filter to $Y_t$ gives
\[
\Pi(B)Y_t=\widetilde{Y}_t.
\]
Under the model,
\[
\widetilde{Y}_t
=
\sum_{k=-\infty}^{\infty}\beta_k\widetilde{X}_{t-k}
+
\widetilde{Z}_t.
\]
Therefore,
\[
\rho_k(\widetilde{X},\widetilde{Y})
\]
is informative about the coefficient $\beta_k$. If
\[
\rho_k(\widetilde{X},\widetilde{Y})
\]
is significantly nonzero, this suggests that $X_{t-k}$ helps predict $Y_t$.
\end{example}

\section{ARCH and GARCH Models}

ARCH/GARCH stands for autoregressive conditional heteroskedasticity and generalized autoregressive conditional heteroskedasticity. These models are designed for time series with changing conditional variance.

ARMA models describe conditional means. ARCH/GARCH models describe conditional variances.

For example, an AR(1) model can be written as
\[
Y_t=\phi Y_{t-1}+\varepsilon_t,
\qquad
\varepsilon_t\iid(0,\sigma^2).
\]
Then
\[
\E[Y_t\mid Y_{t-1}]=\phi Y_{t-1}.
\]
Equivalently,
\[
Y_t
=
\underbrace{\phi Y_{t-1}}_{\text{conditional mean}}
+
\underbrace{\sigma}_{\text{conditional scale}}e_t,
\qquad
e_t\iid(0,1).
\]
In contrast, ARCH/GARCH models allow the conditional scale to change over time.

\subsection{ARCH(1)}

\begin{definition}
An ARCH(1) model is
\[
Y_t=\sigma_te_t,
\qquad
e_t\iid(0,1),
\]
with
\[
\sigma_t^2=\alpha_0+\alpha_1Y_{t-1}^2,
\]
where
\[
\alpha_0>0,
\qquad
0\leq \alpha_1<1.
\]
\end{definition}

This model is conditionally heteroskedastic because
\[
\Var(Y_t\mid Y_{t-1})=\sigma_t^2=\alpha_0+\alpha_1Y_{t-1}^2
\]
depends on the past value $Y_{t-1}$.

\subsubsection{Stationarity of ARCH(1)}

Although the conditional variance changes over time, the unconditional variance can be constant. Since
\[
\E[Y_t]=\E[\sigma_te_t]=0,
\]
and
\[
\Var(Y_t)=\E[Y_t^2],
\]
we compute
\[
\E[Y_t^2]
=
\E[\sigma_t^2e_t^2]
=
\E[\sigma_t^2]\E[e_t^2]
=
\E[\sigma_t^2],
\]
because $e_t$ is independent of the past and $\E[e_t^2]=1$.

Now
\[
\E[\sigma_t^2]
=
\E[\alpha_0+\alpha_1Y_{t-1}^2]
=
\alpha_0+\alpha_1\E[Y_{t-1}^2].
\]
Under stationarity,
\[
\E[Y_t^2]=\E[Y_{t-1}^2].
\]
Let
\[
v=\Var(Y_t)=\E[Y_t^2].
\]
Then
\[
v=\alpha_0+\alpha_1v.
\]
Solving,
\[
v=\frac{\alpha_0}{1-\alpha_1}.
\]
Thus $\alpha_1<1$ is required for a finite unconditional variance.

\subsubsection{Financial Returns}

Let $X_t$ be the price of an asset. The return can be defined as
\[
Y_t
=
\log(X_t)-\log(X_{t-1})
\approx
\frac{X_t-X_{t-1}}{X_{t-1}}.
\]
Returns are often modeled as having approximately zero mean and time-varying volatility. In the ARCH(1) model,
\[
Y_t=\sigma_te_t,
\qquad
\sigma_t^2=\alpha_0+\alpha_1Y_{t-1}^2.
\]
If $|Y_{t-1}|$ is large, then $\sigma_t^2$ is large, which makes another large movement in $Y_t$ more likely. This captures volatility clustering.

\subsubsection{Properties of ARCH(1)}

For ARCH(1),
\[
\E[Y_t]=0,
\]
and
\[
\Var(Y_t)=\frac{\alpha_0}{1-\alpha_1}.
\]

For $h>0$, $Y_t$ and $Y_{t-h}$ are uncorrelated. For example,
\[
\Cov(Y_t,Y_{t-1})
=
\E[Y_tY_{t-1}]
=
\E[\sigma_te_tY_{t-1}].
\]
Since $e_t$ is independent of the past and has mean zero,
\[
\E[\sigma_te_tY_{t-1}]
=
\E[\sigma_tY_{t-1}\E(e_t\mid \mathcal{F}_{t-1})]
=
0.
\]
Thus the process is uncorrelated, but it is not independent.

The dependence appears in the squared process. Since
\[
Y_t^2=\sigma_t^2e_t^2,
\]
large $Y_{t-1}^2$ increases $\sigma_t^2$, which makes large $Y_t^2$ more likely.

\subsubsection{Kurtosis of ARCH(1)}

Assume
\[
e_t\sim \Normal(0,1).
\]
Let
\[
\mu_4=\E[Y_t^4].
\]
Since
\[
Y_t^4=\sigma_t^4e_t^4
\]
and $\E[e_t^4]=3$, we get
\[
\mu_4=3\E[\sigma_t^4].
\]
Now
\[
\sigma_t^2=\alpha_0+\alpha_1Y_{t-1}^2,
\]
so
\[
\sigma_t^4
=
(\alpha_0+\alpha_1Y_{t-1}^2)^2
=
\alpha_0^2+2\alpha_0\alpha_1Y_{t-1}^2+\alpha_1^2Y_{t-1}^4.
\]
Therefore,
\[
\mu_4
=
3\left(
\alpha_0^2
+
2\alpha_0\alpha_1\E[Y_{t-1}^2]
+
\alpha_1^2\E[Y_{t-1}^4]
\right).
\]
By stationarity,
\[
\E[Y_{t-1}^2]=\Var(Y_t)=\frac{\alpha_0}{1-\alpha_1},
\qquad
\E[Y_{t-1}^4]=\mu_4.
\]
Thus
\[
\mu_4
=
3\left(
\alpha_0^2
+
2\alpha_0\alpha_1\frac{\alpha_0}{1-\alpha_1}
+
\alpha_1^2\mu_4
\right).
\]
Solving,
\[
\mu_4(1-3\alpha_1^2)
=
3\alpha_0^2
\left(
1+\frac{2\alpha_1}{1-\alpha_1}
\right).
\]
Since
\[
1+\frac{2\alpha_1}{1-\alpha_1}
=
\frac{1+\alpha_1}{1-\alpha_1},
\]
we get
\[
\mu_4
=
\frac{3\alpha_0^2(1+\alpha_1)}
{(1-\alpha_1)(1-3\alpha_1^2)}.
\]
This requires
\[
3\alpha_1^2<1.
\]

The kurtosis is
\[
K
=
frac{\mu_4}{\Var(Y_t)^2}
=
\frac{
\frac{3\alpha_0^2(1+\alpha_1)}
{(1-\alpha_1)(1-3\alpha_1^2)}
}{
\left(\frac{\alpha_0}{1-\alpha_1}\right)^2
}.
\]
Simplifying,
\[
K
=
3\frac{1-\alpha_1^2}{1-3\alpha_1^2}.
\]
The excess kurtosis is
\[
K-3
=
\frac{6\alpha_1^2}{1-3\alpha_1^2}>0.
\]
Thus ARCH(1) has heavier tails than a normal distribution.

\subsubsection{Connection to AR(1)}

The squared ARCH(1) process behaves like an AR(1). Starting from
\[
Y_t^2=\sigma_t^2e_t^2,
\]
and
\[
\sigma_t^2=\alpha_0+\alpha_1Y_{t-1}^2,
\]
we can write
\[
Y_t^2
=
\alpha_0+\alpha_1Y_{t-1}^2+\eta_t,
\]
where
\[
\eta_t
=
\sigma_t^2(e_t^2-1).
\]
Since
\[
\E[e_t^2-1]=0,
\]
the innovation $\eta_t$ has mean zero. It is uncorrelated over time but not independent. Hence $Y_t^2$ follows an AR(1)-type recursion with nonnormal innovations.

\subsubsection{ARCH Effects}

If $\{Y_t\}$ has mean zero and the sample ACF of $Y_t$ is approximately zero, but the sample ACF of $Y_t^2$ is nonzero, this is evidence of ARCH effects. In practice, this means the original returns may look uncorrelated, while their volatility displays dependence.

\subsection{ARCH($p$)}

The ARCH(1) model generalizes to ARCH($p$):
\[
Y_t=\sigma_te_t,
\]
where
\[
\sigma_t^2
=
\alpha_0+\sum_{i=1}^p\alpha_iY_{t-i}^2.
\]
The usual parameter restrictions are
\[
\alpha_0>0,
\qquad
\alpha_i\geq 0.
\]
Larger squared observations in recent lags lead to larger conditional variance.

A sufficient stationarity condition is
\[
\sum_{i=1}^p\alpha_i<1.
\]

\subsection{GARCH($p,q$)}

GARCH models extend ARCH models by including lagged conditional variances:
\[
Y_t=\sigma_te_t,
\]
where
\[
\sigma_t^2
=
\alpha_0
+
\sum_{i=1}^p\alpha_iY_{t-i}^2
+
\sum_{j=1}^q\beta_j\sigma_{t-j}^2.
\]
The usual parameter restrictions are
\[
\alpha_0>0,
\qquad
\alpha_i\geq 0,
\qquad
\beta_j\geq 0.
\]
A sufficient stationarity condition is
\[
\sum_{i=1}^p\alpha_i+\sum_{j=1}^q\beta_j<1.
\]

The squared process has an ARMA-like representation. For GARCH($p,q$), one can write
\[
Y_t^2-\sigma_t^2=\eta_t,
\]
and after substitution obtain an ARMA-type model for $Y_t^2$ with nonnormal innovations.

\section{Fitting an ARIMA--GARCH Model}

Given a time series $\{X_t\}$, an ARIMA--GARCH modeling strategy is:

\begin{enumerate}
\item Transform or difference $X_t$ if needed to obtain a stationary series $W_t$.
\item Fit an ARMA$(p,q)$ model to the conditional mean:
\[
W_t\sim \mathrm{ARMA}(p,q).
\]
\item Check whether the residuals exhibit ARCH effects by examining the sample ACF of $\widehat{\varepsilon}_t^2$ or using an ARCH test.
\item If ARCH effects are present, fit a GARCH model to the residual volatility:
\[
\widehat{\varepsilon}_t
=
\sigma_te_t,
\qquad
\sigma_t^2
=
\alpha_0+\sum_{i=1}^{p_0}\alpha_i\widehat{\varepsilon}_{t-i}^2
+
\sum_{j=1}^{q_0}\beta_j\sigma_{t-j}^2.
\]
\end{enumerate}

The combined model can be summarized as
\[
W_t\sim \mathrm{ARMA}(p,q)+\mathrm{GARCH}(p_0,q_0).
\]

\begin{remark}
In practice, small orders such as GARCH(1,1) are often sufficient. The conditional mean and conditional variance should be checked separately: first diagnose the ARMA residuals for remaining serial correlation, then diagnose squared residuals for volatility dependence.
\end{remark}

\chapter{Special Topics: Exogenous Predictors}

\section{Regression with ARMA Errors}

Consider the model
\[
Y_t = \beta_0 + \sum_{j=1}^k \beta_j X_{j,t} + \varepsilon_t,
\qquad
\Phi(B)\varepsilon_t = \Theta(B) u_t,\quad u_t \iid \Normal(0,\sigma^2).
\]

\subsection{Motivation}

If $\varepsilon_t$ is autocorrelated, then OLS remains unbiased but:
\begin{itemize}
\item inefficient,
\item incorrect standard errors,
\item invalid inference.
\end{itemize}

\subsection{Whitening Transformation (Derivation)}

Apply $\Phi(B)$:
\[
\Phi(B)Y_t = \beta_0 \Phi(B)1 + \sum_{j=1}^k \beta_j \Phi(B)X_{j,t} + \Theta(B)u_t.
\]

Define
\[
Y_t^\ast = \Phi(B)Y_t, \quad X_{j,t}^\ast = \Phi(B)X_{j,t}.
\]

Then
\[
Y_t^\ast = \beta_0^\ast + \sum_{j=1}^k \beta_j X_{j,t}^\ast + \Theta(B)u_t.
\]

If $q=0$, then
\[
Y_t^\ast = \beta_0^\ast + \sum \beta_j X_{j,t}^\ast + u_t,
\]
which is standard regression.

\begin{example}[AR(1) errors]
\[
\varepsilon_t = \phi \varepsilon_{t-1} + u_t.
\]
Then
\[
Y_t - \phi Y_{t-1}
=
\beta_0(1-\phi)
+
\sum_{j=1}^k \beta_j (X_{j,t} - \phi X_{j,t-1})
+
u_t.
\]
\end{example}

\subsection{Estimation}

Two main approaches:

\begin{itemize}
\item Iterative GLS (Cochrane–Orcutt)
\item Maximum likelihood
\end{itemize}

\section{Transfer Function Models}

\begin{definition}
\[
Y_t = \sum_{j=0}^\infty \psi_j X_{t-j} + N_t,
\]
where $N_t$ is ARMA noise.
\end{definition}

\subsection{Rational Representation}

\[
\frac{\omega(B)}{\delta(B)} = \sum_{j=0}^\infty \psi_j B^j,
\]
so
\[
\delta(B)Y_t = \omega(B)X_t + \delta(B)N_t.
\]

\subsection{Impulse Response Interpretation}

$\psi_j$ measures effect of a unit shock in $X_t$ on $Y_{t+j}$.

\begin{example}[Geometric decay]
\[
\omega(B)=\omega_0,\quad \delta(B)=1-\delta_1 B.
\]
Then
\[
\psi_j = \omega_0 \delta_1^j.
\]

Thus
\[
Y_t = \omega_0 \sum_{j=0}^\infty \delta_1^j X_{t-j} + N_t.
\]
\end{example}

\subsection{Identification via CCF}

Use:
\begin{itemize}
\item prewhitening of $X_t$
\item cross-correlation between filtered $X_t$ and $Y_t$
\end{itemize}

\section{ARMAX Models}

\begin{definition}
\[
\Phi(B)Y_t = \beta_0 + \sum_{j=1}^k \beta_j X_{j,t} + \Theta(B)\varepsilon_t.
\]
\end{definition}

\subsection{Expanded Form}

\[
Y_t =
\sum_{i=1}^p \phi_i Y_{t-i}
+
\sum_{j=1}^k \beta_j X_{j,t}
+
\varepsilon_t
-
\sum_{i=1}^q \theta_i \varepsilon_{t-i}.
\]

\subsection{Forecast Derivation}

For ARMAX(1,1):
\[
Y_t = \phi Y_{t-1} + \beta X_t + \varepsilon_t - \theta \varepsilon_{t-1}.
\]

Then
\begin{align*}
\widehat{Y}_t(1)
&= \E[Y_{t+1} \mid \mathcal{F}_t] \\
&= \phi Y_t + \beta X_{t+1} - \theta \varepsilon_t.
\end{align*}

\section{Dynamic Regression}

\begin{definition}
\[
Y_t = \sum_{j=1}^k \sum_{\ell=0}^{L_j} \beta_{j,\ell} X_{j,t-\ell} + \varepsilon_t.
\]
\end{definition}

\subsection{Interpretation}

\begin{itemize}
\item Immediate effects: $\beta_{j,0}$
\item Lagged effects: $\beta_{j,\ell}$
\end{itemize}

\begin{example}
\[
Y_t = \beta_0 + \beta_1 X_t + \beta_2 X_{t-1} + \varepsilon_t.
\]
Total effect of a unit increase in $X_t$:
\[
\text{impact over time} = \beta_1 + \beta_2.
\]
\end{example}

\subsection{Collinearity Issue}

Lagged predictors:
\[
X_t, X_{t-1}, X_{t-2}
\]
are highly correlated → unstable estimates.

\section{Intervention Analysis}

\begin{definition}
Intervention variable:
\[
I_t =
\begin{cases}
0, & t < T_0,\\
1, & t \ge T_0.
\end{cases}
\]
\end{definition}

\subsection{Model}

\[
Y_t = \mu + \omega I_t + \varepsilon_t.
\]

\begin{example}[Step change]
\[
\E[Y_t] =
\begin{cases}
\mu, & t<T_0,\\
\mu+\omega, & t\ge T_0.
\end{cases}
\]
\end{example}

\begin{example}[Pulse]
\[
I_t = \ind_{\{t=T_0\}},
\quad
Y_t = \mu + \omega I_t + \varepsilon_t.
\]
\end{example}

\begin{example}[Dynamic intervention]
\[
Y_t = \frac{\omega(B)}{\delta(B)} I_t + N_t.
\]
Captures gradual adjustment.
\end{example}

\section{Forecasting with Exogenous Variables}

\subsection{General Formula}

\[
\widehat{Y}_t(h)
=
\E[Y_{t+h} \mid \mathcal{F}_t, X_{t+1},\ldots,X_{t+h}].
\]

\subsection{Key Issue}

Future $X_t$ must be:
\begin{itemize}
\item known, or
\item forecasted separately.
\end{itemize}

\begin{example}[ARX(1)]
\[
Y_t = \phi Y_{t-1} + \beta X_t + \varepsilon_t.
\]

Then
\[
\widehat{Y}_t(1) = \phi Y_t + \beta X_{t+1},
\]
\[
\widehat{Y}_t(h) = \phi \widehat{Y}_t(h-1) + \beta X_{t+h}.
\]
\end{example}

\subsection{Forecast Error Variance}

If $X_t$ known:
\[
\Var(e_t(h)) = \text{same as ARMA part}.
\]

If $X_t$ forecasted:
\[
\Var(e_t(h)) = \text{ARMA variance} + \text{error from } X_t.
\]

\subsection{Practical Workflow}

\begin{enumerate}
\item Make $X_t$, $Y_t$ stationary
\item Prewhiten $X_t$
\item Examine CCF
\item Specify lag structure
\item Fit ARMAX / transfer model
\item Diagnose residuals
\end{enumerate}

\section{Connection Between Models}

\begin{itemize}
\item Dynamic regression = finite lag transfer function
\item Transfer function = infinite lag structured model
\item ARMAX = dynamic regression + ARMA errors
\end{itemize}

\begin{remark}
These models form a hierarchy from simple regression to fully dynamic stochastic systems with external inputs.
\end{remark}

\chapter{Special Topics: Multivariate Time Series}

\section{Vector Autoregressions}

A vector autoregression models several time series jointly. Instead of treating one variable as the response and the others as fixed exogenous predictors, a VAR treats all variables as potentially endogenous.

Let
\[
Y_t =
\begin{bmatrix}
Y_{1,t}\\
Y_{2,t}\\
\vdots\\
Y_{k,t}
\end{bmatrix}
\in \R^k.
\]
A VAR model explains $Y_t$ using lagged values of the full vector $Y_{t-1},Y_{t-2},\ldots$.

\begin{definition}
A VAR($p$) model is
\[
Y_t
=
c
+
A_1Y_{t-1}
+
A_2Y_{t-2}
+
\cdots
+
A_pY_{t-p}
+
\varepsilon_t,
\]
where
\[
c\in\R^k,
\qquad
A_i\in\R^{k\times k},
\qquad
\varepsilon_t\iid (0,\Sigma).
\]
The covariance matrix $\Sigma$ is a $k\times k$ positive definite matrix.
\end{definition}

The innovation vector is
\[
\varepsilon_t =
\begin{bmatrix}
\varepsilon_{1,t}\\
\varepsilon_{2,t}\\
\vdots\\
\varepsilon_{k,t}
\end{bmatrix}.
\]
The components of $\varepsilon_t$ may be contemporaneously correlated. That is,
\[
\Cov(\varepsilon_{i,t},\varepsilon_{j,t})
\]
may be nonzero for $i\neq j$. However, the innovations are serially uncorrelated:
\[
\Cov(\varepsilon_t,\varepsilon_{t-h})=0
\qquad
\text{for } h\neq 0.
\]

\begin{remark}
A VAR model is a multivariate generalization of an AR model. The important difference is that each variable is allowed to depend on lags of every variable in the system.
\end{remark}

\subsection{Example: Bivariate VAR(1)}

Let
\[
Y_t=
\begin{bmatrix}
Y_{1,t}\\
Y_{2,t}
\end{bmatrix}.
\]
A bivariate VAR(1) is
\[
Y_t=c+A_1Y_{t-1}+\varepsilon_t,
\]
where
\[
c=
\begin{bmatrix}
c_1\\
c_2
\end{bmatrix},
\qquad
A_1=
\begin{bmatrix}
a_{11} & a_{12}\\
a_{21} & a_{22}
\end{bmatrix}.
\]
Thus
\[
\begin{bmatrix}
Y_{1,t}\\
Y_{2,t}
\end{bmatrix}
=
\begin{bmatrix}
c_1\\
c_2
\end{bmatrix}
+
\begin{bmatrix}
a_{11} & a_{12}\\
a_{21} & a_{22}
\end{bmatrix}
\begin{bmatrix}
Y_{1,t-1}\\
Y_{2,t-1}
\end{bmatrix}
+
\begin{bmatrix}
\varepsilon_{1,t}\\
\varepsilon_{2,t}
\end{bmatrix}.
\]
Equivalently,
\[
Y_{1,t}
=
c_1+a_{11}Y_{1,t-1}+a_{12}Y_{2,t-1}+\varepsilon_{1,t},
\]
and
\[
Y_{2,t}
=
c_2+a_{21}Y_{1,t-1}+a_{22}Y_{2,t-1}+\varepsilon_{2,t}.
\]

The coefficient $a_{12}$ measures whether lagged $Y_2$ helps predict $Y_1$, holding lagged $Y_1$ fixed. The coefficient $a_{21}$ measures whether lagged $Y_1$ helps predict $Y_2$, holding lagged $Y_2$ fixed.

\section{VAR($p$)}

The VAR($p$) model is
\[
Y_t
=
c
+
A_1Y_{t-1}
+
\cdots
+
A_pY_{t-p}
+
\varepsilon_t.
\]

For $k$ variables, each $A_i$ is a $k\times k$ matrix:
\[
A_i=
\begin{bmatrix}
a_{11}^{(i)} & a_{12}^{(i)} & \cdots & a_{1k}^{(i)}\\
a_{21}^{(i)} & a_{22}^{(i)} & \cdots & a_{2k}^{(i)}\\
\vdots & \vdots & \ddots & \vdots\\
a_{k1}^{(i)} & a_{k2}^{(i)} & \cdots & a_{kk}^{(i)}
\end{bmatrix}.
\]
The $j$th equation is
\[
Y_{j,t}
=
c_j
+
\sum_{i=1}^p
\sum_{\ell=1}^k
a_{j\ell}^{(i)}Y_{\ell,t-i}
+
\varepsilon_{j,t}.
\]

\subsection{Mean of a VAR($p$)}

Assume the VAR($p$) process is stationary. Let
\[
\mu=\E[Y_t].
\]
Taking expectations in
\[
Y_t
=
c+A_1Y_{t-1}+\cdots+A_pY_{t-p}+\varepsilon_t
\]
gives
\[
\mu
=
c+A_1\mu+\cdots+A_p\mu.
\]
Therefore,
\[
\left(I-A_1-\cdots-A_p\right)\mu=c.
\]
If the matrix $I-A_1-\cdots-A_p$ is invertible, then
\[
\mu
=
\left(I-A_1-\cdots-A_p\right)^{-1}c.
\]

\begin{example}
For a VAR(1),
\[
Y_t=c+AY_{t-1}+\varepsilon_t.
\]
If the process is stationary, then
\[
\mu=c+A\mu.
\]
Thus
\[
(I-A)\mu=c,
\]
so
\[
\mu=(I-A)^{-1}c.
\]
\end{example}

\subsection{Centering a VAR($p$)}

If the stationary mean is $\mu$, define
\[
X_t=Y_t-\mu.
\]
Then
\[
Y_t=\mu+X_t.
\]
Substitute into the VAR($p$) model:
\[
\mu+X_t
=
c
+
A_1(\mu+X_{t-1})
+\cdots+
A_p(\mu+X_{t-p})
+\varepsilon_t.
\]
Collect constants:
\[
X_t
=
c+(A_1+\cdots+A_p)\mu-\mu
+
A_1X_{t-1}
+\cdots+
A_pX_{t-p}
+\varepsilon_t.
\]
Since
\[
\mu=c+(A_1+\cdots+A_p)\mu,
\]
the constant terms cancel, leaving
\[
X_t=A_1X_{t-1}+\cdots+A_pX_{t-p}+\varepsilon_t.
\]
Therefore, after centering, the VAR($p$) has mean zero.

\section{Companion Form}

The companion form rewrites a VAR($p$) as a VAR(1) in a higher-dimensional state vector.

Define
\[
Z_t=
\begin{bmatrix}
Y_t\\
Y_{t-1}\\
\vdots\\
Y_{t-p+1}
\end{bmatrix}
\in \R^{kp}.
\]
Then
\[
Z_t
=
\mathcal{A}Z_{t-1}
+
\eta_t,
\]
where
\[
\mathcal{A}
=
\begin{bmatrix}
A_1 & A_2 & \cdots & A_{p-1} & A_p\\
I_k & 0 & \cdots & 0 & 0\\
0 & I_k & \cdots & 0 & 0\\
\vdots & \vdots & \ddots & \vdots & \vdots\\
0 & 0 & \cdots & I_k & 0
\end{bmatrix},
\]
and
\[
\eta_t=
\begin{bmatrix}
\varepsilon_t\\
0\\
\vdots\\
0
\end{bmatrix}.
\]

\begin{example}
For VAR(2),
\[
Y_t=A_1Y_{t-1}+A_2Y_{t-2}+\varepsilon_t.
\]
Define
\[
Z_t=
\begin{bmatrix}
Y_t\\
Y_{t-1}
\end{bmatrix}.
\]
Then
\[
Z_t
=
\begin{bmatrix}
A_1 & A_2\\
I_k & 0
\end{bmatrix}
\begin{bmatrix}
Y_{t-1}\\
Y_{t-2}
\end{bmatrix}
+
\begin{bmatrix}
\varepsilon_t\\
0
\end{bmatrix}.
\]
Thus
\[
Z_t=\mathcal{A}Z_{t-1}+\eta_t.
\]
\end{example}

\section{Stationarity of VAR Models}

The companion form allows us to state the stationarity condition cleanly.

\begin{theorem}
A VAR($p$) process is stationary if all eigenvalues of the companion matrix $\mathcal{A}$ lie inside the unit circle:
\[
|\lambda_i(\mathcal{A})|<1
\qquad
\text{for all } i.
\]
Equivalently, all roots of
\[
\det(I_k-A_1z-\cdots-A_pz^p)=0
\]
lie outside the unit circle.
\end{theorem}

\begin{proof}
Using companion form,
\[
Z_t=\mathcal{A}Z_{t-1}+\eta_t.
\]
Iterating backward,
\[
Z_t
=
\mathcal{A}Z_{t-1}+\eta_t,
\]
\[
Z_t
=
\mathcal{A}(\mathcal{A}Z_{t-2}+\eta_{t-1})+\eta_t
=
\mathcal{A}^2Z_{t-2}
+
\mathcal{A}\eta_{t-1}
+
\eta_t.
\]
Continuing for $m$ steps,
\[
Z_t
=
\mathcal{A}^{m+1}Z_{t-m-1}
+
\sum_{j=0}^{m}\mathcal{A}^j\eta_{t-j}.
\]
If all eigenvalues of $\mathcal{A}$ have modulus less than $1$, then
\[
\mathcal{A}^{m+1}\to 0
\]
as $m\to\infty$. Therefore,
\[
Z_t
=
\sum_{j=0}^{\infty}\mathcal{A}^j\eta_{t-j}.
\]
This is a vector general linear process. Its coefficient matrices are absolutely summable under the eigenvalue condition, so the process is weakly stationary.

Conversely, if at least one eigenvalue of $\mathcal{A}$ has modulus greater than or equal to $1$, then the powers $\mathcal{A}^j$ do not decay to zero in the required way. The infinite moving average representation fails to be stable, so the VAR is not stationary.
\end{proof}

\subsection{VAR(1) Stationarity}

For VAR(1),
\[
Y_t=AY_{t-1}+\varepsilon_t.
\]
The stationarity condition is
\[
|\lambda_i(A)|<1
\qquad
\text{for all eigenvalues } \lambda_i(A).
\]

\begin{example}
Let
\[
A=
\begin{bmatrix}
0.5 & 0\\
0 & 0.2
\end{bmatrix}.
\]
The eigenvalues are $0.5$ and $0.2$, both inside the unit circle. Therefore, the VAR(1) is stationary.

If instead
\[
A=
\begin{bmatrix}
1.1 & 0\\
0 & 0.2
\end{bmatrix},
\]
then one eigenvalue is $1.1$, which is outside the unit circle. The process is not stationary.
\end{example}

\section{VAR Estimation}

Suppose we observe
\[
Y_1,\ldots,Y_T,
\]
where each $Y_t\in\R^k$. Consider the VAR($p$)
\[
Y_t=c+A_1Y_{t-1}+\cdots+A_pY_{t-p}+\varepsilon_t.
\]
Let
\[
x_t=
\begin{bmatrix}
1\\
Y_{t-1}\\
Y_{t-2}\\
\vdots\\
Y_{t-p}
\end{bmatrix}
\in \R^{1+kp}.
\]
Let
\[
B=
\begin{bmatrix}
c^\top\\
A_1^\top\\
A_2^\top\\
\vdots\\
A_p^\top
\end{bmatrix}
\in \R^{(1+kp)\times k}.
\]
Then the VAR can be written as
\[
Y_t^\top=x_t^\top B+\varepsilon_t^\top.
\]

Stack observations for $t=p+1,\ldots,T$. Define
\[
\mathcal{Y}
=
\begin{bmatrix}
Y_{p+1}^\top\\
Y_{p+2}^\top\\
\vdots\\
Y_T^\top
\end{bmatrix},
\qquad
\mathcal{X}
=
\begin{bmatrix}
x_{p+1}^\top\\
x_{p+2}^\top\\
\vdots\\
x_T^\top
\end{bmatrix},
\qquad
\mathcal{E}
=
\begin{bmatrix}
\varepsilon_{p+1}^\top\\
\varepsilon_{p+2}^\top\\
\vdots\\
\varepsilon_T^\top
\end{bmatrix}.
\]
Then
\[
\mathcal{Y}=\mathcal{X}B+\mathcal{E}.
\]

\subsection{OLS Estimator}

The least squares objective is
\[
S(B)
=
\tr\left[
(\mathcal{Y}-\mathcal{X}B)^\top
(\mathcal{Y}-\mathcal{X}B)
\right].
\]
Equivalently,
\[
S(B)
=
\sum_{t=p+1}^T
\norm{Y_t^\top-x_t^\top B}^2.
\]

Differentiating with respect to $B$ gives
\[
\frac{\partial S(B)}{\partial B}
=
-2\mathcal{X}^\top\mathcal{Y}
+
2\mathcal{X}^\top\mathcal{X}B.
\]
Setting the derivative equal to zero,
\[
\mathcal{X}^\top\mathcal{X}B
=
\mathcal{X}^\top\mathcal{Y}.
\]
Assuming $\mathcal{X}^\top\mathcal{X}$ is invertible,
\[
\widehat{B}
=
(\mathcal{X}^\top\mathcal{X})^{-1}\mathcal{X}^\top\mathcal{Y}.
\]

\begin{remark}
A VAR can be estimated equation-by-equation using OLS because every equation has the same regressors. This is why VAR estimation is computationally simple.
\end{remark}

\subsection{Residual Covariance Estimator}

The fitted residuals are
\[
\widehat{\mathcal{E}}
=
\mathcal{Y}-\mathcal{X}\widehat{B}.
\]
Let
\[
N=T-p
\]
be the number of usable observations. A common estimator of $\Sigma$ is
\[
\widehat{\Sigma}
=
\frac{1}{N}
\widehat{\mathcal{E}}^\top\widehat{\mathcal{E}}.
\]
A degrees-of-freedom adjusted estimator is
\[
\widehat{\Sigma}_{\mathrm{adj}}
=
\frac{1}{N-(1+kp)}
\widehat{\mathcal{E}}^\top\widehat{\mathcal{E}}.
\]

\section{OLS as Maximum Likelihood}

Assume
\[
\varepsilon_t\iid \Normal(0,\Sigma).
\]
Conditional on initial values $Y_1,\ldots,Y_p$, the likelihood is
\[
L(B,\Sigma)
=
\prod_{t=p+1}^{T}
(2\pi)^{-k/2}
|\Sigma|^{-1/2}
\exp\left[
-\frac{1}{2}
\varepsilon_t^\top\Sigma^{-1}\varepsilon_t
\right],
\]
where
\[
\varepsilon_t^\top=Y_t^\top-x_t^\top B.
\]
The log-likelihood is
\[
\ell(B,\Sigma)
=
-\frac{Nk}{2}\log(2\pi)
-\frac{N}{2}\log|\Sigma|
-\frac{1}{2}
\sum_{t=p+1}^{T}
\varepsilon_t^\top\Sigma^{-1}\varepsilon_t.
\]
Using matrix notation,
\[
\ell(B,\Sigma)
=
-\frac{Nk}{2}\log(2\pi)
-\frac{N}{2}\log|\Sigma|
-\frac{1}{2}
\tr\left[
\Sigma^{-1}
(\mathcal{Y}-\mathcal{X}B)^\top
(\mathcal{Y}-\mathcal{X}B)
\right].
\]

For fixed $\Sigma$, maximizing the likelihood with respect to $B$ is equivalent to minimizing
\[
\tr\left[
\Sigma^{-1}
(\mathcal{Y}-\mathcal{X}B)^\top
(\mathcal{Y}-\mathcal{X}B)
\right].
\]
Because every equation has the same regressors, the minimizer is the OLS estimator:
\[
\widehat{B}_{\mathrm{MLE}}
=
(\mathcal{X}^\top\mathcal{X})^{-1}\mathcal{X}^\top\mathcal{Y}.
\]

For fixed $B$, maximizing with respect to $\Sigma$ gives
\[
\widehat{\Sigma}_{\mathrm{MLE}}
=
\frac{1}{N}
(\mathcal{Y}-\mathcal{X}\widehat{B})^\top
(\mathcal{Y}-\mathcal{X}\widehat{B}).
\]

\begin{remark}
Under Gaussian innovations, conditional MLE for the VAR coefficient matrices equals OLS.
\end{remark}

\section{Forecasting with VAR Models}

Consider the VAR($p$)
\[
Y_t=c+A_1Y_{t-1}+\cdots+A_pY_{t-p}+\varepsilon_t.
\]
Given information at time $t$, the one-step forecast is
\[
\widehat{Y}_t(1)
=
c+A_1Y_t+A_2Y_{t-1}+\cdots+A_pY_{t-p+1}.
\]
For $h\geq 2$, the forecast is computed recursively:
\[
\widehat{Y}_t(h)
=
c
+
A_1\widehat{Y}_t(h-1)
+
A_2\widehat{Y}_t(h-2)
+\cdots+
A_p\widehat{Y}_t(h-p),
\]
where for negative or zero forecast horizons, observed values are used:
\[
\widehat{Y}_t(0)=Y_t,
\qquad
\widehat{Y}_t(-j)=Y_{t-j}.
\]

\begin{example}
For VAR(1),
\[
Y_t=c+AY_{t-1}+\varepsilon_t.
\]
The one-step forecast is
\[
\widehat{Y}_t(1)=c+AY_t.
\]
The two-step forecast is
\[
\widehat{Y}_t(2)=c+A\widehat{Y}_t(1)
=c+A(c+AY_t)
=c+Ac+A^2Y_t.
\]
In general,
\[
\widehat{Y}_t(h)
=
\sum_{j=0}^{h-1}A^jc+A^hY_t.
\]
If the VAR is stationary, then $A^h\to 0$, and
\[
\widehat{Y}_t(h)\to (I-A)^{-1}c=\mu.
\]
\end{example}

\section{Granger Causality}

Granger causality is about predictability. One variable Granger-causes another if its past values improve prediction after controlling for the past values already in the model.

\begin{definition}
In a VAR, variable $X_t$ Granger-causes variable $Y_t$ if lagged values of $X_t$ improve the prediction of $Y_t$ after controlling for lagged values of $Y_t$ and the other variables.
\end{definition}

\subsection{Bivariate VAR($p$)}

Consider
\[
\begin{bmatrix}
Y_t\\
X_t
\end{bmatrix}
=
c
+
A_1
\begin{bmatrix}
Y_{t-1}\\
X_{t-1}
\end{bmatrix}
+
\cdots
+
A_p
\begin{bmatrix}
Y_{t-p}\\
X_{t-p}
\end{bmatrix}
+
\varepsilon_t.
\]
The first equation is
\[
Y_t
=
c_1
+
\sum_{i=1}^p a_{11}^{(i)}Y_{t-i}
+
\sum_{i=1}^p a_{12}^{(i)}X_{t-i}
+
\varepsilon_{1,t}.
\]
To test whether $X_t$ Granger-causes $Y_t$, test
\[
H_0:
a_{12}^{(1)}=a_{12}^{(2)}=\cdots=a_{12}^{(p)}=0.
\]
If this null is rejected, then lagged $X_t$ helps predict $Y_t$.

\subsection{F-Test Form}

Let $\mathrm{SSR}_R$ be the residual sum of squares from the restricted model, where the lagged $X_t$ terms are excluded. Let $\mathrm{SSR}_U$ be the residual sum of squares from the unrestricted model. If $p$ restrictions are tested, then
\[
F
=
\frac{(\mathrm{SSR}_R-\mathrm{SSR}_U)/p}
{\mathrm{SSR}_U/(N-m)},
\]
where $m$ is the number of parameters in the unrestricted equation. Under the null,
\[
F\approx F_{p,N-m}.
\]

\begin{remark}
Granger causality is not the same as structural causality. It only means that one variable contains useful predictive information for another variable.
\end{remark}

\section{Impulse Response Functions}

Impulse response functions describe the dynamic effect of shocks.

Consider a mean-zero VAR(1):
\[
Y_t=AY_{t-1}+\varepsilon_t.
\]
Iterating forward:
\[
Y_{t+1}=AY_t+\varepsilon_{t+1},
\]
\[
Y_{t+2}=A(AY_t+\varepsilon_{t+1})+\varepsilon_{t+2}
=
A^2Y_t+A\varepsilon_{t+1}+\varepsilon_{t+2}.
\]
In general,
\[
Y_{t+h}
=
A^hY_t+\sum_{j=0}^{h-1}A^j\varepsilon_{t+h-j}.
\]
Thus the effect of a shock $\varepsilon_t$ on $Y_{t+h}$ is
\[
A^h.
\]

\begin{definition}
For a VAR(1), the impulse response matrix at horizon $h$ is
\[
\Psi_h=A^h.
\]
The $(i,j)$ entry of $\Psi_h$ measures the response of variable $i$ at horizon $h$ to a one-unit shock in innovation $j$ at time $t$.
\end{definition}

\subsection{VAR($p$) Impulse Responses}

Using companion form,
\[
Z_t=\mathcal{A}Z_{t-1}+\eta_t.
\]
Then
\[
Z_{t+h}
=
\mathcal{A}^hZ_t
+
\sum_{j=0}^{h-1}\mathcal{A}^j\eta_{t+h-j}.
\]
The impulse responses for $Y_t$ are obtained from the first $k$ rows of $\mathcal{A}^h$.

Equivalently, a causal VAR($p$) has a vector moving average representation:
\[
Y_t
=
\sum_{j=0}^{\infty}\Psi_j\varepsilon_{t-j},
\]
where
\[
\Psi_0=I_k.
\]
The coefficient matrices satisfy
\[
\Psi_h
=
A_1\Psi_{h-1}
+
A_2\Psi_{h-2}
+\cdots+
A_p\Psi_{h-p},
\]
with
\[
\Psi_0=I_k,
\qquad
\Psi_h=0
\quad
\text{for } h<0.
\]

\begin{exercise}
Derive the recursion for $\Psi_h$.
\end{exercise}

\begin{proof}
Assume
\[
Y_t=\sum_{j=0}^{\infty}\Psi_j\varepsilon_{t-j}.
\]
Substitute this into
\[
Y_t=A_1Y_{t-1}+\cdots+A_pY_{t-p}+\varepsilon_t.
\]
Then
\[
\sum_{j=0}^{\infty}\Psi_j\varepsilon_{t-j}
=
A_1\sum_{j=0}^{\infty}\Psi_j\varepsilon_{t-1-j}
+\cdots+
A_p\sum_{j=0}^{\infty}\Psi_j\varepsilon_{t-p-j}
+\varepsilon_t.
\]
The coefficient of $\varepsilon_t$ on the left is $\Psi_0$, while on the right it is $I_k$. Hence
\[
\Psi_0=I_k.
\]
For $h\geq 1$, match coefficients of $\varepsilon_{t-h}$. The right-hand side contributes
\[
A_1\Psi_{h-1}+A_2\Psi_{h-2}+\cdots+A_p\Psi_{h-p}.
\]
Therefore,
\[
\Psi_h
=
A_1\Psi_{h-1}
+
A_2\Psi_{h-2}
+\cdots+
A_p\Psi_{h-p}.
\]
\end{proof}

\section{Cholesky Orthogonalized Impulse Responses}

In a VAR, innovations can be contemporaneously correlated:
\[
\Sigma=
\Cov(\varepsilon_t)
\]
is usually not diagonal. Therefore, a one-unit shock to one component of $\varepsilon_t$ may not be independent of shocks to the other components.

To create orthogonal shocks, use the Cholesky decomposition:
\[
\Sigma=PP^\top,
\]
where $P$ is lower triangular.

Define
\[
\varepsilon_t=Pu_t,
\]
where
\[
u_t\iid (0,I_k).
\]
Then
\[
\Cov(\varepsilon_t)
=
\Cov(Pu_t)
=
P\Cov(u_t)P^\top
=
PP^\top
=
\Sigma.
\]

The moving average representation is
\[
Y_t
=
\sum_{j=0}^{\infty}\Psi_j\varepsilon_{t-j}.
\]
Substitute $\varepsilon_t=Pu_t$:
\[
Y_t
=
\sum_{j=0}^{\infty}\Psi_jP u_{t-j}.
\]
Thus the orthogonalized impulse response at horizon $h$ is
\[
\Theta_h=\Psi_hP.
\]

\begin{definition}
The Cholesky orthogonalized impulse response matrix at horizon $h$ is
\[
\Theta_h=\Psi_hP,
\]
where $P$ is the lower triangular Cholesky factor satisfying $\Sigma=PP^\top$.
\end{definition}

The $(i,j)$ entry of $\Theta_h$ measures the response of variable $i$ at horizon $h$ to a one-standardized-unit orthogonal shock in variable $j$.

\begin{remark}
Cholesky impulse responses depend on variable ordering. The variable ordered first can affect variables ordered later contemporaneously, but variables ordered later do not affect variables ordered earlier contemporaneously.
\end{remark}

\subsection{Bivariate Cholesky Example}

Let
\[
\Sigma=
\begin{bmatrix}
\sigma_1^2 & \sigma_{12}\\
\sigma_{12} & \sigma_2^2
\end{bmatrix}.
\]
The Cholesky factor is
\[
P=
\begin{bmatrix}
p_{11} & 0\\
p_{21} & p_{22}
\end{bmatrix}.
\]
Since
\[
PP^\top
=
\begin{bmatrix}
p_{11}^2 & p_{11}p_{21}\\
p_{11}p_{21} & p_{21}^2+p_{22}^2
\end{bmatrix},
\]
matching entries gives
\[
p_{11}=\sigma_1,
\]
\[
p_{21}=\frac{\sigma_{12}}{\sigma_1},
\]
and
\[
p_{22}
=
\sqrt{\sigma_2^2-\frac{\sigma_{12}^2}{\sigma_1^2}}.
\]
Thus
\[
P=
\begin{bmatrix}
\sigma_1 & 0\\
\dfrac{\sigma_{12}}{\sigma_1}
&
\sqrt{\sigma_2^2-\dfrac{\sigma_{12}^2}{\sigma_1^2}}
\end{bmatrix}.
\]

If
\[
Y_t=AY_{t-1}+\varepsilon_t,
\]
then
\[
\Psi_h=A^h,
\]
so the orthogonalized impulse response is
\[
\Theta_h=A^hP.
\]

\section{Forecast Error Variance Decomposition}

For a causal VAR, write the moving average representation:
\[
Y_t=\sum_{j=0}^{\infty}\Psi_j\varepsilon_{t-j}.
\]
The $h$-step forecast error is
\[
e_t(h)=Y_{t+h}-\widehat{Y}_t(h).
\]
Using the moving average representation,
\[
Y_{t+h}
=
\sum_{j=0}^{h-1}\Psi_j\varepsilon_{t+h-j}
+
\sum_{j=h}^{\infty}\Psi_j\varepsilon_{t+h-j}.
\]
The second sum is known at time $t$, while the first sum contains future innovations. Therefore,
\[
e_t(h)
=
\sum_{j=0}^{h-1}\Psi_j\varepsilon_{t+h-j}.
\]
The forecast error covariance matrix is
\[
\Var(e_t(h))
=
\sum_{j=0}^{h-1}
\Psi_j\Sigma\Psi_j^\top.
\]

Using the Cholesky factorization $\Sigma=PP^\top$, define
\[
\Theta_j=\Psi_jP.
\]
Then
\[
\Var(e_t(h))
=
\sum_{j=0}^{h-1}
\Theta_j\Theta_j^\top.
\]

The contribution of orthogonal shock $\ell$ to the forecast error variance of variable $i$ is
\[
\sum_{j=0}^{h-1}
\left(\Theta_j\right)_{i\ell}^2.
\]
Therefore, the forecast error variance decomposition is
\[
\mathrm{FEVD}_{i\ell}(h)
=
\frac{
\sum_{j=0}^{h-1}
\left(\Theta_j\right)_{i\ell}^2
}{
\sum_{\ell'=1}^{k}
\sum_{j=0}^{h-1}
\left(\Theta_j\right)_{i\ell'}^2
}.
\]
The denominator is the $i$th diagonal entry of $\Var(e_t(h))$.

\begin{remark}
For each variable $i$ and horizon $h$,
\[
\sum_{\ell=1}^k \mathrm{FEVD}_{i\ell}(h)=1.
\]
Thus FEVD gives the fraction of forecast uncertainty in variable $i$ attributable to each orthogonal shock.
\end{remark}

\section{Cointegration and VECM}

Many multivariate time series are nonstationary. If each component of $Y_t$ is integrated of order $1$, then
\[
Y_t\sim I(1).
\]
This means that $Y_t$ is nonstationary, but
\[
\Delta Y_t=Y_t-Y_{t-1}
\]
is stationary.

\begin{definition}
The components of $Y_t$ are cointegrated if there exists a nonzero vector $\beta$ such that
\[
\beta^\top Y_t
\]
is stationary, even though the individual components of $Y_t$ are nonstationary.
\end{definition}

The vector $\beta$ is called a cointegrating vector.

\begin{example}
Suppose $X_t$ and $Z_t$ are both $I(1)$, but
\[
X_t-Z_t
\]
is stationary. Define
\[
Y_t=
\begin{bmatrix}
X_t\\
Z_t
\end{bmatrix}.
\]
Then
\[
\beta=
\begin{bmatrix}
1\\
-1
\end{bmatrix}
\]
is a cointegrating vector because
\[
\beta^\top Y_t=X_t-Z_t
\]
is stationary.
\end{example}

\subsection{Deriving the VECM from VAR($p$)}

Start with the VAR($p$):
\[
Y_t
=
A_1Y_{t-1}
+
A_2Y_{t-2}
+
\cdots
+
A_pY_{t-p}
+
\varepsilon_t.
\]
Subtract $Y_{t-1}$ from both sides:
\[
\Delta Y_t
=
(A_1-I)Y_{t-1}
+
A_2Y_{t-2}
+
\cdots
+
A_pY_{t-p}
+
\varepsilon_t.
\]
Now express lagged levels in terms of $Y_{t-1}$ and differences:
\[
Y_{t-2}=Y_{t-1}-\Delta Y_{t-1},
\]
\[
Y_{t-3}=Y_{t-1}-\Delta Y_{t-1}-\Delta Y_{t-2},
\]
and so on. After collecting terms, the model can be written as
\[
\Delta Y_t
=
\Pi Y_{t-1}
+
\Gamma_1\Delta Y_{t-1}
+
\cdots
+
\Gamma_{p-1}\Delta Y_{t-p+1}
+
\varepsilon_t,
\]
where
\[
\Pi=A_1+\cdots+A_p-I,
\]
and
\[
\Gamma_i
=
-\sum_{j=i+1}^{p}A_j,
\qquad
i=1,\ldots,p-1.
\]

\begin{definition}
The vector error correction model is
\[
\Delta Y_t
=
\Pi Y_{t-1}
+
\sum_{i=1}^{p-1}\Gamma_i\Delta Y_{t-i}
+
\varepsilon_t.
\]
If $\Pi$ has rank $r$, where $0<r<k$, then
\[
\Pi=\alpha\beta^\top,
\]
where $\beta$ contains the cointegrating vectors and $\alpha$ contains adjustment coefficients.
\end{definition}

The term
\[
\beta^\top Y_{t-1}
\]
is the long-run equilibrium error. The matrix $\alpha$ determines how each variable adjusts when the system deviates from equilibrium.

\subsection{Simple VECM Example}

Suppose
\[
Y_t=
\begin{bmatrix}
X_t\\
Z_t
\end{bmatrix},
\]
and $X_t-Z_t$ is stationary. Then
\[
\beta=
\begin{bmatrix}
1\\
-1
\end{bmatrix}.
\]
A simple VECM is
\[
\Delta Y_t
=
\alpha(X_{t-1}-Z_{t-1})
+
\varepsilon_t,
\]
where
\[
\alpha=
\begin{bmatrix}
\alpha_1\\
\alpha_2
\end{bmatrix}.
\]
Equivalently,
\[
\Delta X_t=\alpha_1(X_{t-1}-Z_{t-1})+\varepsilon_{1,t},
\]
\[
\Delta Z_t=\alpha_2(X_{t-1}-Z_{t-1})+\varepsilon_{2,t}.
\]
If $X_{t-1}>Z_{t-1}$, then the equilibrium error is positive. The signs of $\alpha_1$ and $\alpha_2$ determine whether $X_t$ falls, $Z_t$ rises, or both adjust to restore equilibrium.

\section{R Code Examples}

\subsection{Simulating a VAR(1)}

\begin{lstlisting}[language=R, caption={Simulating a bivariate VAR(1) process}]
set.seed(123)

T <- 300
k <- 2

A <- matrix(c(0.6, 0.2,
              0.1, 0.4), nrow = 2, byrow = TRUE)

Sigma <- matrix(c(1.0, 0.5,
                  0.5, 1.0), nrow = 2, byrow = TRUE)

eps <- MASS::mvrnorm(n = T, mu = rep(0, k), Sigma = Sigma)

Y <- matrix(0, nrow = T, ncol = k)

for (t in 2:T) {
  Y[t, ] <- A %*% Y[t - 1, ] + eps[t, ]
}

matplot(Y, type = "l", lty = 1,
        xlab = "t", ylab = "Value",
        main = "Simulated VAR(1)")
legend("topright", legend = c("Y1", "Y2"), lty = 1, col = 1:2)
\end{lstlisting}

\subsection{Fitting a VAR Model in R}

\begin{lstlisting}[language=R, caption={Fitting a VAR model in R}]
library(vars)

data <- as.data.frame(Y)
names(data) <- c("Y1", "Y2")

fit <- VAR(data, p = 1, type = "const")

summary(fit)
roots(fit)
serial.test(fit, lags.pt = 12, type = "PT.asymptotic")
\end{lstlisting}

\subsection{Impulse Responses and FEVD in R}

\begin{lstlisting}[language=R, caption={Impulse responses and FEVD in R}]
irf_fit <- irf(fit, impulse = "Y1", response = "Y2",
               n.ahead = 20, ortho = TRUE, boot = TRUE)

plot(irf_fit)

fevd_fit <- fevd(fit, n.ahead = 20)
plot(fevd_fit)
\end{lstlisting}

\end{document}